% 本文件由 LaTeX 排版 Agent 根据有效 translation.json 自动生成。
% author=John Cassian (c. 360-c. 435); work_id=3509; title=论降生成人

\chapter{论\OriginalTerm{降生成人}{Incarnation}}

% structure: p0001 source=h1 target=omitted reason=page-title-already-rendered-by-parent

第一卷\newline{}第二卷\newline{}第三卷\newline{}第四卷\newline{}第五卷\newline{}第六卷\newline{}第七卷\par

\SourceRecord{Translated by C.S. Gibson.；Nicene and Post-Nicene Fathers, Second Series；Vol. 11.；Edited by Philip Schaff and Henry Wace.；Buffalo, NY: Christian Literature Publishing Co.,；1894.；<http://www.newadvent.org/fathers/3509.htm>.}

% structure: p0001 source=h1 target=section reason=page-title

\section{论降生成人（第一卷）}

% structure: p0002 source=h2 target=subsection reason=relative-heading-rank; base=h2

\subsection{序言}

及我既完成\WorkTitle{属灵论谈}各书——其价值在于所表达的思想而非所用的语言（因我粗拙的言辞不足以匹配圣人们深邃的思想）——我本已考虑并几乎决定退隐于沉默之中（因我羞于暴露自己的无知），以便尽可能以未来谨慎的缄默弥补我说话的大胆。然而，我亲爱的良，我尊贵而备受敬重的朋友，罗马教会及神圣职务的荣耀，你以可赞的热忱和最恳切的爱情克服了我的决心和意图，因为你将我从所决定的沉默幽暗中拖出，进入一个我所惧怕的公开法庭，并迫使我在仍为过去之事羞愧之际承担新的劳苦。我虽不足以承担较小的任务，你却强逼我与更大的任务相称。因为在那些微小的作品中，我们以微薄之能向主献上少许祭品时，若非主教的命令，我决不敢尝试或致力于任何事。因此，藉着你，我们的主题和语言的重要性都增加了。因为以前我们受命谈论主的事务，如今你却要求我们谈论主亲自的降生成人和荣耀。于是，从前我们仿佛由司铎之手被带入圣殿的圣地，如今在你指引和保护下，我们好似进入至圣所。这份荣誉伟大，但任务极其危险，因为圣所之奖赏和天主的酬报只能通过战胜仇敌而获得。因此，你要求和命令我们举起软弱的手反对一个新异端和新信德的敌人，并使我们站定，好比对抗那致命毒蛇可怕的张口，使我召唤时，预言之力和福音之言的神圣能力就能摧毁那蜿蜒而行反对天主教会的那龙。我听从你的恳求，顺从你的命令；因为我宁愿在有关我自己的事上信赖你而非信赖我自己，尤其我主耶稣基督的爱情命令我如此，如同命令你一样，因为祂亲自在你身上给我这一命令。因为在这事上，你比我更关切，因你的判断比我的职责更危险。因为在我身上，无论我是否胜任你所命令的，我服从和谦卑的事实本身多少会成为我的辩辞；若我尚可主张：我的服从更有价值，若我能做之事更少。因为出于我们的富足，我们易于顺从任何人的命令；但那种渴望超越其能力的人，其工作才是伟大而奇妙的。因此，这是你的工作和事务，也是你以此为羞之事。请祈祷和恳求，免得我的笨拙使你的选择蒙羞；并且，若我们未能达到你对我们的期望，你不至于因轻率决定而命令错误，而我因服从之故，顺从是正确的。\par

% structure: p0004 source=h2 target=subsection reason=relative-heading-rank; base=h2

\subsection{第一章}

\Emph{异端与诗人所咏之九头蛇比较。}\par

诗人的传说告诉我们，古时九头蛇被斩首后，反因受伤而滋生更茂；因此，藉着一种奇特且前所未闻的奇迹，它的损失竟成为那怪物的增益，它因死亡而增多，而那种非凡的繁殖力使刽子手的刀所斩下的一切加倍，直到那急切寻求毁灭它的人，劳苦流汗，发现自己的努力屡屡因徒劳无功而受挫，便在战斗的勇气上添加了狡诈的技巧，并且如他们所说，用火施以火剑，斩断了那怪物身体的众多子孙；因此，当内脏被焚烧，那可怕繁殖力的叛逆悸动被灼伤后，那些怪异的产物终于终结。同样，教会的异端也与诗人幻想所造的九头蛇有几分相似；因为\Emph{它们}也用致命的舌头向我们嘶嘶作响；\Emph{它们}也喷出致命的毒液，并且在斩首后又复生。但因疾病复发时药物不可或缺，且病势越凶险，治疗就应越迅速，我主天主能成就此事：在外邦人的虚构中关于九头蛇之死所想象的，在教会的争战中成为真实，即圣神的火剑能灼伤那最危险的新异端的内在部分，使其被摧毁，以致其怪异的繁殖力最终不再响应其垂死的悸动。\par

% structure: p0007 source=h2 target=subsection reason=relative-heading-rank; base=h2

\subsection{第二章}

\Emph{描述彼此孳生的不同异端怪物。}\par

因为这些不自然的种子所生的枝条在教会中并非新鲜事。主田地的收成一直不得不忍受牛蒡和荆棘，并且其中不断长出令人窒息的莠子的枝条。因为由此产生了\OriginalTerm{厄彼雍派}{Ebionites}、\OriginalTerm{撒伯流派}{Sabellians}、\OriginalTerm{亚略派}{Arians}，以及\OriginalTerm{欧诺米派}{Eunomians}和\OriginalTerm{马其顿尼派}{Macedonians}，以及\OriginalTerm{福提诺派}{Photinians}和\OriginalTerm{阿波林派}{Apollinarians}，以及教会中所有其他的莠子和毁坏信德果实的蓟草。这些当中最早的是\Person{厄彼雍}，他过于急于维护我们主的人性，却剥夺了它与天主性的结合。但在他之后，\Person{撒伯流}的分裂派作为对上述异端的反作用爆发出来，他宣称圣父、圣子、圣神之间没有区别，从而在可能的范围内不敬地混淆了各位，未能区分神圣而不可言说的天主圣三。在我们提到的他之后，接下来是亚略乖张的亵渎，它为了避免混淆神圣各位的表象，宣称天主圣三中有不同且不相似的本体。但在时间上于他之后，虽在邪恶上相似，来了\Person{欧诺米}，他虽允许天主圣三的各位是神圣且彼此相似的，却坚持它们彼此分离；因此，在承认相似的同时否定了平等。\Person{马其顿尼}也以不可宽恕的邪恶亵渎圣神，虽然允许圣父和圣子同体，却称圣神为受造物，从而犯罪得罪了整个天主性，因为对天主圣三中任何一者施加伤害不可能不影响整个天主圣三。但\Person{福提诺}虽然允许生于童贞女的耶稣是天主，却在观念上错谬，认为祂的天主性始于祂人性开始之时；而\Person{阿波林}由于不精确地理解天主与人的结合，错误地相信祂没有人灵魂。因为给我们主耶稣基督添加不属于祂的东西与剥夺属于祂的东西是同样严重的错误。因为当祂被以不同于祂实际情况的方式谈论时——即使这似乎增添了祂的光荣——却是冒犯。就这样，一个接一个，作为对异端的反作用，他们产生了异端，并且都教导彼此不同、但同样敌对信仰的东西。而就在最近，即在我们这个时代，我们看到一种最具毒害的异端从\Place{比利时}最大的城市兴起，虽然它的错误毫无疑问，但它的名称却有疑问，因为它从厄彼雍派的旧根上长出了新头，因此它究竟该被称为古老还是新近，仍是一个问题。因为就其支持者而言，它是新的；但就其错误的性质而言，它是古老的。的确，它亵渎地教导说，我们的主耶稣基督是作为一个纯粹的人而生的，并坚持认为祂后来获得天主性的光荣和能力源于祂人性的价值，而非源于祂的天主性；并以此教导说，祂并非凭借属于祂自己的天主性的权利而始终拥有天主性，而是后来才作为祂辛劳和受苦的赏报获得天主性。然而，当它亵渎地教导说，我们的主和救主在祂出生时并非天主，后来才被接纳进入天主性时，它确实接近于现在兴起的这个异端，可以说是它的第一表亲，与它相似，并且与厄彼雍主义和这些新异端都协调一致，在时间上居于它们之间，在邪恶上与之相连。虽然还有其他一些像我们提到过的，但要全部描述它们就太长了。我们现在也没有着手列举那些已死去的，而是要驳斥那些新生的。\par

% structure: p0010 source=h2 target=subsection reason=relative-heading-rank; base=h2

\subsection{第三章。}

\Emph{他描述贝拉奇派的致命错误。}\par

无论如何，我们认为不应遗漏这一事实，它对于上述源于白拉奇谬误的异端是独特而特有的；即，他们声称耶稣基督作为纯粹的人毫无罪污地生活，竟进而宣称人若愿意也能无罪。因为他们幻想，如果耶稣基督作为纯粹的人而无罪，那么所有人在没有天主帮助的情况下，也能成为这位作为纯粹的人、未分享\OriginalTerm{天主性}{Godhead}者所能成为的一切。于是他们断定，任何人与我们的主耶稣基督之间没有区别，因为任何人都能通过努力和奋斗，获得基督通过其热忱和努力所获得的同样成果。由此，他们陷入了更严重、更反常的疯狂，竟说我们的主耶稣基督来到这世界，并非为给人类带来救赎，而是为善行树立榜样，即人通过跟随祂的教导，走在同样的美德道路上，便可获得同样的美德奖赏：从而尽其所能地毁灭了祂神圣降临的一切益处和神圣救赎的一切恩宠，因为他们宣称人可凭己力获得天主为人类的救恩而死所成就的一切。他们还补充说，我们的主和救主在圣洗后成为基督，在复活后成为天主，前者归因于祂傅油的奥迹，后者归因于祂苦难的功绩。由此，这位宣称我们的主和救主生为纯粹之人的新异端（并非新的）作者，指出他所说的恰好与\OriginalTerm{白拉奇主义者}{Pelagians}在他之前所说的相同，并承认从他的谬误可推知：既然他断言我们的主耶稣基督作为纯粹之人完全无罪，那么他在亵渎中必须坚持所有人靠自己都能无罪，他也不愿承认主的救赎是对其榜样而言必要的事，因为人（如他们所说）能靠自己的努力到达天国。这一点毫无疑问，事实本身向我们表明。因此，他通过干涉鼓励白拉奇主义者的抱怨，并将他们的案例引入他的著作，因为他狡猾地（或更真实地说，诡诈地）庇护他们，并因对之邪恶的喜好而推荐与他们相近的毒害教导，因为他深知自己与他们是同一见解、同一精神，故而对一个与己相近的异端被逐出教会感到痛苦，因为他知道这异端在邪恶上与自己的完全相连。\par

% structure: p0013 source=h2 target=subsection reason=relative-heading-rank; base=h2

\subsection{第四章。}

\Emph{\Person{莱波留斯}与一些其他人一同撤回其\OriginalTerm{白拉奇主义}{Pelagianism}。}\par

然而，既然这一毒刺丛的产物已藉天主的助佑和仁慈得医治，我们现在也当向我们的主天主祈祷：正如那较早的异端与这新异端在某些方面相近那样，愿祂赐予那些有相似恶劣开端者同样幸福的结局。因为那时尚为隐修士、现今成为\OriginalTerm{司铎}{presbyter}的\Person{莱波留斯}，如上所述，他曾追随\Person{白拉奇}的教导（或更确切地说，恶行），并在高卢是上述异端最早、最大的捍卫者之一，经我们劝诫、受天主纠正后，如此高尚地谴责了他先前的错误信念，以致他的改过几乎与许多人未受损的信德一样值得庆幸。因为从不犯错是最好的；其次便是良好地弃绝错误。他于是醒悟过来，怀着忧愁却无羞耻地承认了自己的错误，不仅在当时和现今所在的非洲，而且向高卢各城寄出了包含其认罪和忧愁的忏悔信，以便他回归信德的消息能传达到他最先背离信德的地方，并使那些曾见证他错误的人此后也见证他的改过。\par

% structure: p0016 source=h2 target=subsection reason=relative-heading-rank; base=h2

\subsection{第五章。}

\Emph{通过\Person{莱波留斯}的例子，他确立了一个事实，即公开的罪应由公开的认罪来补偿；并通过\Person{莱波留斯}的话教导了关于降生成人应持的正确观点。}\par

我们因此认为从他的忏悔，更确切地说是哀叹中，引用一部分是适当的，有两个原因：使他们的反悔成为对我们的见证，对软弱者的榜样，并使那些在错误中不耻于跟随他们的人，在改正中也不耻于跟随他们；并使他们如遭受同样的疾病一样，因同样的疗法而得治愈。他于是承认自己观点的乖谬，见到信德之光，写信给\Place{高卢}主教们，开始说：\Quote{我几乎不知道，我最可敬的主教们和蒙福的司铎们，我应当先控告自己什么，先为自己辩解什么。笨拙、骄傲和愚昧无知，连同错误的观念、不谨慎的热忱，以及（说实话）逐渐衰落的软弱信德，这一切都被我接纳，并且如此盛行，以致我为屈服于如此众多的罪而感到羞愧，同时我又为能够将它们从我的灵魂中驱逐出去而深深感恩。}稍后他又说：\Quote{如果，我们不理解天主这种德能，在自己的臆想和意见中自作聪明，因为害怕天主似乎会扮演不符合祂身份的角色，而认为一个人是与天主一同出生的，以致我们将属于天主的单独归于天主，将属于人的单独归于人，那么我们显然为天主圣三增添了第四个位格，并从独一天主子中开始制造两个基督，而不是一个；愿我们的主、天主耶稣基督亲自保护我们免受此害。因此我们宣认，我们的主、天主子耶稣基督，天主的独生子，祂为自身在万世之前由圣父所生，当时期圆满时，祂为了我们由圣神和卒世童贞玛利亚取得人性而成孕出生，祂出生时就是天主；在我们宣认肉躯与圣言两种本性的同时，我们始终以虔诚的信念和信德承认，同一位格不可分割地既是天主又是人；我们说，从祂取肉躯之时起，凡属于天主的都赐给了人，凡属于人的都与天主结合。在这意义上，\BibleQuote{圣言成了血肉}\BibleRef{若 1:14}——不是祂通过任何转换或变化开始成为祂所不是的，而是通过神圣的\OriginalTerm{经纶}{economy}，圣父的圣言从未离开圣父，却屈尊成为真人，独生子通过那只有祂才理解的奥秘而降生成人（因为我们负责\Emph{相信}，祂负责\Emph{理解}）。因此，天主\InnerQuote{圣言}亲自接受了属于人的一切，成为人，而所取的人性接受了属于天主的一切，不能不是天主；但虽然祂被称为降生成人且不相混合，我们不可认为祂的本体有任何减损：因为天主知道如何不遭受任何腐化而通传自己，并真正地通传自己。祂知道如何接收到自己内而不因此增加，正如祂知道如何分施自己而不受损失。我们不应以软弱的理智，根据可见的证据和实验，从平等且相互进入的受造物的情况中做出猜测，也不应认为天主和人混合在一起，并将肉躯与圣言（即天主性与人性）的如此融合产生某种身体。天主禁止我们想象两个本性这样被模铸在一起而成为同一本体。因为这种混合对两者都是破坏性的。因为天主包含万物而不被包含，进入万物而不被进入，充满万物而不被充满，同时完整地处处存在并弥漫各处，祂通过灌注自己的德能而恩赐地通传自己于人性。}稍后又说：\Quote{因此，天主而人，耶稣基督，天主子，真正为我们由圣神和卒世童贞玛利亚而生。所以在两个本性中，圣言与肉躯成为一体，以致每个本性在其自身内继续完美地存在，天主性毫不受损地通传给人性，而人性参与天主性；并非一位是天主，另一位是人，而是同一位既是天主也是人；再者，那同时也是人的天主被称为并确实是耶稣基督，天主的独生子；因此我们必须始终小心并相信，不否认我们的主耶稣基督，天主子，真天主（我们宣认祂在万世之前与圣父同在，与圣父平等），从祂取肉躯之时起，就成为天主而人。我们也不可想象祂逐渐随着时间成为天主，认为祂在复活前是一种状态，在复活后是另一种状态，而是祂始终具有同样的圆满和德能。}又稍后说：\Quote{但因为天主的圣言屈尊通过取人性而降于人性，而人性通过被天主所取而被提升至圣言，天主圣言完全成为完整的人。因为不是天主圣父成为人，也不是圣神，而是圣父的独生子；因此我们必须坚持，肉躯与圣言是同一个位格：以便忠信无疑地相信，同一天主子永不可分割，在两种本性中存在（祂在祂肉躯的日子也被称为\InnerQuote{巨人}），真正取了属于人的一切，并永远真正拥有属于天主的：因为虽然祂在软弱中被钉死，但祂却因天主的德能而生活。}\par

% structure: p0019 source=h2 target=subsection reason=relative-heading-rank; base=h2

\subsection{第六章}

\Emph{天主教的合一教义应被视为正统信仰。}\par

因此，他的这一宣认——即所有天主教信友的信德——得到了\Place{非洲}的所有主教和\Place{高卢}的所有主教的赞同。从未有人反对这信德而不犯不信之罪：因为否认正确且被证实的事就是承认错误。因此，所有人的一致意见本身应足以驳倒异端：因为所有人的权威显示无可置疑的真理，无人争议之处就产生完美的理由。因此，如果有人试图持有与此相反的意见，我们首先应谴责他的执拗，而非听他的主张；因为攻击众人判断的人预先宣布了自己的定罪；扰乱已被众人决定之事的人，甚至不被听取。当真理已由所有人一劳永逸地确立时，凡与它相反之物，凭此事实立即被辨认为虚假，因为它与真理不同。因此，公认单凭这一点就足以定罪，即他与真理的判断不同。然而，因为对某一体系的解释无损于该体系，真理充分探讨时总是更加闪亮，且最好让犯错者通过讨论被纠正，而非被严厉指责定罪，所以我们应尽己所能，藉天主的助佑，治愈这在新异端者身上出现的古老异端，好使他们因天主的仁慈恢复健康时，他们的治愈可为我们的圣善信德作证，而不是定罪成为公义严厉的实例。惟愿真理确实临于我们关于它的讨论和话语中，并以天主屈尊就卑来到人间的良善扶助我们人性的软弱；因为祂正是为此首要目的而愿意生于世上、生于人间，即不让虚假再有容身之地。\par

\SourceRecord{Translated by C.S. Gibson.；Nicene and Post-Nicene Fathers, Second Series；Vol. 11.；Edited by Philip Schaff and Henry Wace.；Buffalo, NY: Christian Literature Publishing Co.,；1894.；<http://www.newadvent.org/fathers/35091.htm>.}

% structure: p0001 source=h1 target=section reason=page-title

\section{论降生成人（卷二）}

% structure: p0002 source=h2 target=subsection reason=relative-heading-rank; base=h2

\subsection{第一章}

\Emph{后来的异端者的错误如何在他们的作者和发起人身上受到谴责和驳斥。}\par

正如我们在第一卷开始时列出一些东西，以此表明我们的新异端者不过是古代异端枝条的一个分支，对先前异端者的应有谴责足以确保对他做出应有的判决。因为他具有相同的根源，从相同的休耕地中生长出来，所以他已经在那些前任身上受到了充分的谴责，尤其是那些在这些之前立即犯错的人非常恰当地谴责了这些现在正在断言的事情，以至于他们自己党派的例子在两个方面都应该足够给他们：即那些被恢复的人和那些被谴责的人。因为如果他们能够改正，他们在自己党派的改正中就有了补救措施。如果他们无法改正，他们在自己人的谴责中接受了判决。但为了不让我们被认为对他们有偏见而不是公正地判断，我们将提出他们那些实际有害的断言，或者更确切地说，是他们的亵渎愚昧：拿起\BibleQuote{首先拿起信德的盾牌，以及圣神的剑，即是天主的圣言} \BibleRef{弗 6:16}，以便当古蛇的头再次抬起时，曾经在那些古龙身上砍掉它的同一把圣言的剑现在也可能在这些新蛇身上砍掉它。因为既然这些人的错误与那些先前的错误相同，那些人的斩首应该算作这些人的斩首；并且由于毒蛇复活并发出毒气对抗主的教会，并通过它们的嘶嘶声使一些人跌倒，我们必须针对这些新疾病在那些旧疗法之外添加一种新疗法，以便即使已经做过的事情不足以治愈疾病，我们现在正在做的事情足以恢复那些患病的人。\par

% structure: p0005 source=h2 target=subsection reason=relative-heading-rank; base=h2

\subsection{第二章}

\Emph{证明天主之母童贞玛利亚不仅是\OriginalTerm{基督之母}{Christotocos}，也是\OriginalTerm{天主之母}{Theotocos}，以及基督是真天主。}\par

异端者，无论你是谁，你否认天主由童贞女所生，并声称吾主耶稣基督之母玛利亚不应被称为\OriginalTerm{天主之母}{Theotocos}，而应称为\OriginalTerm{基督之母}{Christotocos}，即仅为基督之母，而非天主之母。你说：没有人能生出时间上在先者。关于这个极其愚蠢的论点——你竟以为血肉之心能理解天主的诞生，并幻想人性推理能解释祂威严的奥迹——若天主允许，我们稍后将加以讨论。此时，我们将以神圣的见证来证明基督是天主，玛利亚是天主之母。请听，天主的使者如何对牧羊人谈论天主的诞生：\BibleQuote{今天在达味城中，为你们诞生了一位救世者，他是主基督。}\BibleRef{路 2:11}为使你不把基督当作一个普通人，他特意加上\Quote{主}和\Quote{救世者}的名号，好让你毫不怀疑你所承认的救世者就是天主，并且（因救恩的职分只属于天主的能力）你不可质疑祂具有天主的能力，因为你已知道救恩的能力在祂内。但或许这不足以说服你的不信，因为上主的天使称祂为主和救世者，而非天主或天主子，而你确实极其邪恶地否认祂是天主，尽管你承认祂是救世者。那么，请听总天使加俾额尔对童贞玛利亚的宣告：\BibleQuote{圣神要临于你，至高者的能力要庇荫你，因此，那要诞生的圣者，将称为天主子。}\BibleRef{路 1:35}你看见吗？当他即将指出天主的诞生时，他首先谈论天主的工程。因为他说：\BibleQuote{圣神要临于你，至高者的能力要庇荫你。}天使说得极妙，他以言语的神性来解释神性工程的尊威。因为圣神圣化了童女的子宫，以祂天主性的能力吹入其中，如此将祂自己赋予并通传给人性；使那原本与祂无关的成为祂自己的，并以祂自己的能力和尊威将其纳为己有。唯恐人性的软弱无法承受天主性的进入，至高者的能力加强了这位永受尊敬的童女，以庇荫的保护拥抱她，支持她身体的软弱；而人性的软弱因神圣庇荫的支持，足以完成圣怀孕的不可言喻奥迹。他说：\BibleQuote{因此，圣神要临于你，至高者的能力要庇荫你。}若仅仅是一个普通人由纯贞女所生，为何要如此仔细地提及天主的来临？为何要有天主性本身的介入？当然，如果仅仅是一个人从人而生，肉从肉而生，只需一道命令或天主的旨意即可。因为如果天主的旨意和祂的命令就足以形成诸天、大地、海洋、宝座、座位、天使、总天使、宰制者和掌权者，总之，足以创造天上所有军旅和无数千万的天主神军（\BibleQuote{祂一发言，它们就造成；祂一出命，它们就形成}），那么，为何那足以产生一切神圣事物、创造所有地上和天上事物的能力，却不足以（按你们的说法）创造一个人？而天主的能力和尊威为何不信任那诞生一个婴孩的工作，那曾创造一切事物的工作？但所有那些工作都是由天主的命令完成的，而诞生只能由祂的来临完成，理由是：除非天主允许，人不能怀祂；除非祂自己进入，人不能生祂。因此，总天使指出神圣的尊威将临于童女，我是说，因如此伟大的事件不能由人的安排实现，他宣布那要诞生的荣耀将临于怀孕。于是，圣言、圣子降临了；圣神的尊威临在；圣父的能力庇荫；以便在圣怀孕的奥迹中，整个天主圣三共同合作。他说：\BibleQuote{因此，那要诞生的圣者，将称为天主子。}他极妙地加上\Quote{因此}，以表明这事将\Emph{因此}发生，\Emph{因为}先前的事已发生；并且\Emph{因为}天主在她怀孕时已临于她，\Emph{因此}天主在诞生时也必临在。当童女不明白时，他为这大事给出了理由，说：\BibleQuote{因为圣神要临于你，因为至高者的能力要庇荫你，因此，那要诞生的圣者将称为天主子。}也就是说：为使你不致对如此伟大的工程和这伟大奥秘的安排一无所知，天主的尊威因此完全临于你；因为天主子将由你而生。对此还能有什么疑问？还有什么可说的？他说天主要临于她；天主子要诞生。现在，如果你愿意，请问：天主子怎能不是天主？那生出天主的，又怎能不是\OriginalTerm{天主之母}{Theotocos}？单是这点就该足够了；是的，这该完全满足你了。\par

% structure: p0008 source=h2 target=subsection reason=relative-heading-rank; base=h2

\subsection{第三章}

\Emph{以旧约经文继续同一论证}\par

但既然有大量为神圣诞生的证人，即所有为此而写的、为之作证的内容，让我们也稍微考察一下旧约中关于天主的预告，好使你知道：天主由童贞女诞生这件事，不仅在当时发生时被宣布，而且从世界之初就已预言，为了使这即将成就的、不可言说的事件，因着在它尚未来到之前就已被预告，从而消除它在实现时的怀疑。因此，先知\Person{依撒意亚}说：\Quote{看，一位童贞女将怀孕生子，人要称祂的名为厄玛奴耳，解释出来就是：天主与我们同在。}\BibleRef{依 7:14}这里还有什么怀疑的余地呢，你这怀疑的人？先知说童贞女将怀孕：童贞女\Emph{已经}怀孕了；说将生一个儿子：儿子\Emph{已经}生了；说祂要被称为天主：祂被称为天主。因为祂因着那本性而被称为那名。因此，当天主之神说祂应被称为天主时，祂就证实那使自己与神圣名号的一切共融隔绝的人是没有天主之神的。\Quote{看，}他说，\Quote{一位童贞女将怀孕生子，人要称祂的名为厄玛奴耳，解释出来就是：天主与我们同在。}但有一点，你那狡猾的怀疑可能抓住不放，即说先知所宣布的祂应被称为的，不是指祂的天主性的荣耀，而是指人们称呼祂的名字。但我们怎么办呢？因为基督在福音书中从未被称为这个名字，尽管不能称天主之神通过先知说了谎。那么这是怎么回事呢？显然，我们应当明白，那预言当时预告的是祂的神圣本性的名，而不是祂人性的名。因为在祂与天主性结合的人性中，祂在福音书中接受了另一个名字，显然这个名属于祂的人性，那个名属于祂的天主性。但让我们继续前进，召唤其他真实的证人来确立真理：因为当我们谈论天主性时，没有比祂自己的证人更能确立天主性的了。因此，同一位先知在别处说：\Quote{因为有一个儿子为我们诞生了，一个婴孩赐给了我们；权柄将担在祂的肩上；祂的名字要称为谋略天使、强有力的天主、永远之父、和平之王。}正如上文先知明确说祂要被称为厄玛奴耳，这里他说祂要被称为\Quote{谋略天使、强有力的天主、永远之父、和平之王}（尽管我们肯定从未读到祂在福音书中被称为这些名字）：当然，这是要我们明白，这些不是属于祂人性的术语，而是属于祂神圣本性的；而在福音书中使用的名字属于祂所取的人性，而这个名字属于祂固有的能力。因为天主将以人的形式诞生，这些名字在神圣的安排中如此分配，以致人性被赋予了一个人的名字，天主性被赋予了一个神圣的名字。因此他说：\Quote{祂要被称为谋略天使、强有力的天主、永远之父、和平之王。}不是的，异端者，无论你是谁，不是先知觉满圣神，效法你的榜样，把所生的比作无感觉的铸像和雕像。因为\Quote{一个儿子，}他说，\Quote{为我们诞生了，一个婴孩赐给了我们；权柄将担在祂的肩上；祂的名字要称为谋略天使、强有力的天主。}为使你不至于以为他所宣布为天主的另一位，与祂肉身所生的那位不同，他加了一个关于祂诞生的词，说：\Quote{一个婴孩为我们诞生了，一个儿子赐给了我们。}你看见先知用了多少称号来表明祂肉身诞生的真实吗？因为他特意称祂为儿子和婴孩，好使所生婴孩的方式通过一个指祂婴儿时期的名称更清楚地显示出来；而圣神无疑预见到亵渎神明的异端者的这种邪僻，用所用的术语和话语向全世界展示，所生的是天主；即使一个异端者决心发出亵渎，他也找不到任何亵渎的漏洞。因此他说：\Quote{一个儿子为我们诞生了，一个婴孩赐给了我们；权柄将担在祂的肩上；祂的名字要称为谋略天使、强有力的天主、永远之父、和平之王。}他教导说，所生的这个婴孩既是和平之王，又是永远之父和强有力的天主。那么还有什么推脱的余地呢？所生的这个婴孩不能与在祂内出生的天主分开，因为他称祂所论及的那位所生者为永远之父；他称祂为婴孩，却预告祂是强有力的天主。这是什么，异端者？你将逃往何处？每个地方都被围住堵住了：没有逃脱的可能。除非你最终被迫承认你\Emph{不愿}理解的错误。但我们不满足于这些确实足够的章节，让我们查考圣神通过另一位先知说了什么。\Quote{人岂能刺穿他的天主，因为你们在刺穿我？}为使预言的主题更加清楚，先知预言了他所宣布的主的苦难，仿佛从他所谈论的那位的口中说出。\Quote{人岂能刺穿他的天主，因为你们在刺穿我？}我问，当我们主天主被带去受十字架时，祂不就像说了这话吗？为什么你们不承认我是你们的救赎者？为什么你们不认识为你们穿上血肉的天主？你们在为你们的救主预备死亡吗？你们在把生命之作者带向死亡吗？我是你们的天主，你们举起的；你们的天主，你们钉死的。我问，这是什么错误，或什么疯狂？\Quote{人岂能刺穿他的天主，因为你们在刺穿我？}你看见这些话多么准确地描述了实际发生的事吗？你能要求更明确或更清楚的吗？你看见神圣的见证从摇篮一直跟随我们降生成人的主耶稣基督直到祂所背负的十字架吗？正如你在这里所见，祂在肉身中诞生时，你在别处读到祂为天主，祂在十字架上被刺时也为天主。因此，在那里，当论及祂的诞生时，祂被先知称为天主；而在这里，当涉及祂的十字架时，祂最清楚地被命名为天主；好使接受人性丝毫不损害祂天主性的尊荣，肉身的卑微和苦难的耻辱也不影响祂威严的荣耀；因为祂屈尊俯就如此卑微的诞生，以及祂甘愿忍受苦难的伟大善行，应当增加我们对祂的爱和虔敬；因为如果祂对我们倾注的爱越多，我们对祂尊敬越少，这无疑是一个大而可憎的罪。\par

% structure: p0011 source=h2 target=subsection reason=relative-heading-rank; base=h2

\subsection{第四章}

\Emph{他引用宗徒保禄的见证来证明同一道理。}\par

但是，略过这些无法尽述的事，因为祂所赐的恩惠没有限量，现在是我们请教宗徒保禄的时候了——他是最坚强、最清晰的见证者，能以最可信的方式向我们讲述一切关于天主的事，因为天主始终从他的胸中说话。他，这位万民的特选导师，被派遣来摧毁外邦迷信的谬误，以下方式见证我们的主天主的恩宠和降临：他说：\BibleQuote{天主的恩宠，以及我们的救主，向众人显现出来了，教导我们弃绝不虔敬的行为和世俗的私欲，在今世度着克己、公正、虔敬的生活，期待那幸福的希望，以及伟大的天主和我们的救主耶稣基督的荣耀显现。}\BibleRef{铎 2:11}他说：\Quote{\Emph{出现}了天主的恩宠，我们的救主。}他绝妙地使用了一个适合表明新恩宠和新诞生的降临的词；因为通过说\Quote{\Emph{出现}}，他指出了一种新恩宠和新诞生的来临，因为从那时起，当天主显现为生于世界时，新恩宠的恩赐开始显现。因此，通过使用一个恰当且极为合适的词，他几乎像用手指指着一样，指出了这新恩宠的光芒。因为最恰当地被称为\Emph{出现}的，是突然的光照使其显现的。正如我们在福音中读到，那颗星\Emph{出现}在东方的贤士面前；以及在出谷纪中：他说：\BibleQuote{\Emph{出现}了，在荆棘丛中，有火焰中的天使向梅瑟显现。}\BibleRef{出 3:2}因为在所有这一切以及圣经中其他异象的情况下，圣经特别规定使用这个词，用以指称那以不寻常的光辉照耀的事物为\Quote{出现}。因此，宗徒也深知那天上的恩宠，即在神圣诞生临近时出现的恩宠，便用了一个适用于光辉显现的词来指明它；特意要说它\Emph{出现}了，如同以新光的光辉照耀一样。因此，\Quote{出现了天主的恩宠，我们的救主。}难道在这个地方，你能就名号的模糊性提出任何诡辩，说\Quote{基督}是一回事，\Quote{天主}是另一回事，或者将\Quote{救主}从祂名号的荣耀中分离，将\Quote{主}与天主性分开吗？看，天主的器皿从天主说话，并以最清晰的陈述证明天主的恩宠从玛利亚显现。为了让你不至于否认天主从玛利亚显现，他立即加上了救主的名号，特意让你相信，那由玛利亚所生的是天主，而你不能否认祂是作为救主而生的，符合这段经文：\BibleQuote{因为今天在达味城中，为你们诞生了一位救主。}\BibleRef{路 2:11}啊，真正由天主赐给外邦人的卓越导师！因为他知道这野蛮的异端愚行将会兴起，它会将天主的名号用于争议，并且不会停止用祂自己的头衔来诽谤天主；因此，正是为了不让异端者将救主的名号与天主性分离，他首先放置了天主的名号，以便首要的天主之名可以声称所有后续的名号都属于祂，并且没有人会想象在后续的话中，基督被当作一个纯粹的人来谈论，因为通过使用的第一个词，他已经教导说祂是天主。同一宗徒说：\BibleQuote{期待那幸福的希望，以及伟大的天主和我们的救主耶稣基督的荣耀显现。}当然，那位神圣智慧的导师看到，单纯朴素的教导本身不足以应付魔鬼诡计的狡猾伎俩，除非他以极其谨慎的保护来加固信仰的神圣宣讲。因此，尽管他在上文已经使用了天主救主的名号，但在此他又加上\Quote{耶稣基督}，唯恐你以为仅仅救主的名号不足以向你指明我们的主耶稣基督，而未能明白你所承认的天主救主，就是同一耶稣基督。那么他说什么呢？他说：\BibleQuote{期待那幸福的希望，以及伟大的天主和我们的救主耶稣基督的荣耀显现。}这里关于我们主的名号，毫无欠缺；你在这里看到天主、救主、耶稣和基督。但当你看到这一切时，你看到它们都属于天主。因为你听到祂是天主，同时也是救主；你听到祂是天主，同时也是耶稣；你听到祂是天主，同时也是基督。那由天主性所结合并联合的，不能因名号的多样性而分离；因为无论你寻求其中哪一个，你都会找到一切。救主是天主，耶稣是天主，基督是天主。在你所听到的一切中，虽然使用的名号众多，但它们都属于同一权能的位格。因为既然救主是天主，耶稣是天主，基督是天主，很容易看出这些虽然称谓不同，但在威严上是合一的。而且当你清楚地听到同一个位格在每种情况下都被称为天主时，你当然可以清楚地看到，在所有这些情况下，所说的是唯一的天主。因此，你不能再试图从给予主的不同名字中找出权力的区别，也不能由于名号的多样性而造成位格的差异。你不能说：基督生于玛利亚，但天主不是；因为宗徒宣告天主是。你不能说：耶稣生于玛利亚，但天主不是；因为宗徒证明天主是。你不能说：救主生了，但天主不是；因为宗徒支持天主是在的事实。你没有逃脱之路。无论你取用主的哪个名号，祂就是你所讲的天主。你无话可说：无可主张：在你的邪恶谎言中无可杜撰。你可以以不虔信的不信拒绝相信：但在你的亵渎之事上，你没有任何可否认的。\par

% structure: p0014 source=h2 target=subsection reason=relative-heading-rank; base=h2

\subsection{第五章}

\Emph{从我们通过基督领受的神圣恩宠的恩赐中，他推论出祂确实是天主。}\par

虽然我们早些时候已开始谈论我们的主和救主的神圣恩宠，但我仍愿从圣经中就同一主题多讲一些。我们在\BibleBook{宗徒}{大事录}中读到，宗徒\Person{雅各伯}这样反驳那些认为接受福音后仍应负旧约法律之轭的人：\Quote{为什么，}他说，\Quote{你们试探天主，要把一个我们祖先和我们自己都不能负的轭放在门徒的颈上？但是，我们相信，我们因主耶稣基督的恩宠得救，正像他们一样。}\BibleRef{宗 15:10} 宗徒无疑是在谈论由耶稣基督所赐的这份恩宠之恩赐。现在请你回答我：你认为这为众人得救而赐予的恩宠，是由人还是由天主赐予的？若你说是由人，天主的器皿\Person{保禄}会向你疾呼：\BibleQuote{天主我们的救主的恩宠已经出现。}\BibleRef{铎 2:11} 他教导说，这恩宠来自神圣恩赐的结果，而非人性的软弱。即使神圣的见证尚不够充分，事情本身的真理也会作证，因为脆弱而属尘世的事物不可能提供持久不朽的价值；没有人能将自己所缺乏的给予别人，也不能提供充足的东西，而他承认自己正因缺乏那东西而受苦。因此，你不得不承认恩宠来自天主。那么，是天主赐予了它。但它是由我们的主耶稣基督赐予的。所以，主耶稣基督就是天主。如果祂确是天主，那么生产祂的那一位便是\OriginalTerm{天主之母}{Theotocos}，即天主的母亲。除非你打算躲进一个极其荒谬且亵渎的矛盾中，即否认那生天主的是天主之母，同时又无法否认那出生者是天主。然而，让我们看看天主的福音如何论及同一主题——我们主的恩宠：\Quote{恩宠和真理，}它说，\Quote{是由耶稣基督而来的。}\BibleRef{若 1:17} 如果基督只是一个凡人，这些怎能由基督而来？祂里面怎能有神圣的能力，如果如你所说，祂只有人的本性？天上的慷慨从何而来，如果祂拥有的是尘世的贫乏？因为没有人能给出他自己尚未拥有的东西。既然基督赐予了神圣恩宠，祂就已经拥有祂所赐予的。没有人能够同时承受如此截然不同的事物，既忍受穷人的匮乏，又显示慷慨者的宽宏。因此，知道天上一切财富的宝藏都在基督内的宗徒\Person{保禄}，正确地致书给众教会说：\BibleQuote{愿我们的主耶稣基督的恩宠与你们同在。}\BibleRef{格前 16:23} 尽管他已经多次教导天主与基督同一，所有天主性的荣耀都在祂内，所有神性的圆满都实际地住在祂内，但在这里他无疑正确地只祈求基督的恩宠，而未加上\Quote{天主}一词：因为他经常教导天主的恩宠与基督的恩宠相同，现在他完美地只祈求基督的恩宠，因为他知道在基督的恩宠内包含了天主全部的恩宠。因此他说：\BibleQuote{愿我们的主耶稣基督的恩宠与你们同在。} 如果耶稣基督只是一个凡人，那么在他希望基督的恩宠赐予众教会时，他是在希望一个人的恩宠被赐予；而说\Quote{愿基督的恩宠与你们同在}，意思就是：愿人的恩宠与你们同在，愿肉身的恩宠与你们同在，愿肉身软弱的恩宠，愿人性脆弱的恩宠与你们同在！或者，如果他所希望的是一个人的恩宠，为什么他竟提到\Quote{恩宠}这个词？因为如果所希望的东西不存在，就没有理由希望；也不应该祈祷将一个人的恩宠赐予他们，因为按照你们的说法，那人并不拥有所希望恩宠的实在。因此，你看这是多么荒谬可笑——不，不是可笑之事，而是可哭之事，因为对某些轻浮之人来说是笑料，对虔诚忠信的灵魂却是哭泣之事；他们因你无信的愚昧而流出爱德之泪，也为他人不虔诚的愚昧而流下虔诚的泪水。让我们暂时恢复理智，喘一口气，因为这种观念不仅缺乏智慧，也缺乏圣神，它显然缺少属神的智慧，与救恩的圣神毫无关系。\par

% structure: p0017 source=h2 target=subsection reason=relative-heading-rank; base=h2

\subsection{第6章}

\Emph{赐予神圣恩宠的能力并非在时间进程中才来到基督身上，而是从祂的诞生之初就内在于祂。}\par

然而你或许会说，主耶稣基督的这恩宠，即宗徒所写的那恩宠，并非与祂一同诞生的，而是后来藉着神性降临于祂而注入祂内的，因为你声称主耶稣基督这个人（你称之为一个纯粹的人）并非与天主一同诞生，而是后来被天主所接纳的；并且，藉着这恩宠，当神性赐予那人时，恩宠也同时赐予了祂。我们所说的也无非是：天主的恩宠是随同神性一同降临的，因为天主的恩宠在某种程度上是实际神性的赐予和丰盛恩宠的馈赠。那么，或许会有人认为我们之间的区别在于时间而非本质，因为我们说神性是随同耶稣基督一同诞生的，而你说神性是后来注入祂内的。但事实是，如果你否认神性是随同主一同诞生的，那么你后来便无法按照信德来宣认了；因为同一件事不可能部分不虔敬又部分虔敬，也不可能部分属于信德部分属于谬误。那么首先我问你：你是否认为由童贞玛利亚出生的主耶稣基督只是人子，还是祂同时也是天主子？因为我们，即所有持守大公信德的人，我们所有人都相信、确信、知道并宣认祂两者都是，即祂是人子，因为由女人所生，同时是天主子，因为因神性而受孕。那么，你承认祂两者都是，即天主子与人子，还是你说祂只是人子？如果只是人子，那么宗徒与先知，甚至那促成这受孕的圣神本身，都向你呼喊。你这最无耻的口被所有神圣法令的见证堵住了；被神圣著作和圣洁的见证堵住了；也被天主的福音本身如同被天主的手堵住了。那位大能的天使加俾额尔，在匝加利亚的事件中曾以言语的力量制止了不信的声音，他以自己的口唇更强烈地谴责了你的亵渎与罪恶之声，对天主之母童贞玛利亚说：\BibleQuote{圣神要临到你，至高者的能力要庇荫你，因此那要诞生的圣者，将称为天主子。} \BibleRef{路 1:35} 你看，耶稣基督首先被宣称为天主子，而后按照肉身才成为人子。因为当童贞玛利亚将要生出主时，她因圣神降临于她并至高者能力的合作而受孕。由此你可见，我们的主与救主的根源必来自祂受孕的源头；由于祂是因圆满完全的神性整个降临于童贞女而出生的，祂不可能成为人子，除非祂首先是天主子；因此，当天使被派遣来宣告祂的诞生与神圣出生时，在讲述了祂受孕的奥秘之后，又加了一句表达祂诞生的话，说：\BibleQuote{因此那要诞生的圣者，将称为天主子}（即祂将被称为祂所由生之父的子）。因此，耶稣基督是天主子，因为祂由天主所生、由天主所孕。但如果祂是天主子，那么祂无疑就是天主；但如果祂是天主，那么祂就不缺少天主的恩宠。事实上，祂从不缺少那祂自身所创造的东西。因为恩宠和真理是由耶稣基督而来的。\par

% structure: p0020 source=h2 target=subsection reason=relative-heading-rank; base=h2

\subsection{第七章}

\Emph{在基督内，天主性、威严、大能与德能如何自永远且恒常圆满地存在。}\par

\Emph{因此}，一切的恩宠、德能、大能、天主性，是的，还有实际神性的圆满与光荣，无论在天上或地上、在母腹中或诞生时，都与祂同在、在祂内。凡属于天主的，从未对天主有所欠缺。因为天主性与天主永远同在，无论何处何时都不与祂分离。因为天主处处圆满而完美地临在。祂不受分割、变化或减损；因为天主既不能有所增添，也不能有所减损，祂的神性既不会减少，也不会增加。祂在地上时同一位也在天上；祂在卑微时同一位也在至高之处；祂在人性渺小時同一位也在神性荣耀中。因此，宗徒论及基督的恩宠时，意指天主的恩宠，这是正确的。因为基督就是天主所是的一切。在祂作为人而受孕的那一刻，天主的一切德能、神性的一切圆满都降临了；因为神性的一切完美都来自祂的根源。祂的人性从未脱离神性，因为祂从神性获得存在的本身。因此，首先，无论你是否愿意，你不能否认这一点，即主耶稣基督是天主子，尤其是在福音中总领天使宣告说：\BibleQuote{那要诞生的圣者，将称为天主子。} 这一点确定之后，请记住：凡你所读到关于基督的，都是读到关于天主子；凡你所读到关于主或耶稣的，都属于天主子。因此，当你在这所有你听到的词语中认出一个神性的称号时，既然你看到在每种情况下你都应当理解所指的乃是天主子，那么，如果你愿意，请向我证明，你如何能将神性与天主子分开。\par

\SourceRecord{Translated by C.S. Gibson.；Nicene and Post-Nicene Fathers, Second Series；Vol. 11.；Edited by Philip Schaff and Henry Wace.；Buffalo, NY: Christian Literature Publishing Co.,；1894.；<http://www.newadvent.org/fathers/35092.htm>.}

% structure: p0001 source=h1 target=section reason=page-title

\section{论降生成人（第三卷）}

% structure: p0002 source=h2 target=subsection reason=relative-heading-rank; base=h2

\subsection{第一章}

\Emph{基督在位格的统一中既是天主又是人，按肉身出自以色列和童贞玛利亚。}\par

那位教会的神圣导师在致罗马人书中，当他责备，或更确切地说，哀叹犹太人的不信，即他自己的弟兄时，使用了这些话：\Quote{我竟情愿自己，}他说，\Quote{从基督被诅咒，为了我的弟兄，他们按肉身是我的亲属，他们是以色列人，义子的名分、荣耀、盟约、法律的颁布、礼仪和恩许都属于他们；圣祖是他们的，基督按肉身也是从他们而来的，祂是在万有之上，永远受赞颂的天主。} \BibleRef{罗 9:3} 哦，这位最忠信的宗徒和最仁慈的亲属的爱德！他在无边的爱德中情愿死去——作为亲属为亲人，作为导师为门徒。那么他为什么情愿死去呢？只有一个原因：即他们可以活着。但他们的生命在于什么呢？简单地说，正如他所说的，在于他们能认识到一位神圣的基督按肉身是从他们自己的肉身所生的。因此宗徒更加悲伤，因为那些应当因祂出自自己的血统而更爱祂的人，却不明白祂是由以色列而生的。\Quote{从他们而来的，}他说，\Quote{是按肉身的基督，祂是在万有之上，永远受赞颂的天主。} 他明确地指出，按肉身，那在万有之上、永远受赞颂的天主基督是从他们而生的。你当然不能否认基督按肉身是从他们而生的。但同一位从他们而生的人就是天主。你怎么能回避这一点？你怎么能逃避？宗徒说按肉身从以色列所生的基督就是天主。如果你能，教导我们祂在什么时间不存在。他说：\Quote{从他们而来的，是按肉身的基督，祂是在万有之上，天主。} 你看，因为宗徒把这两者联合并连接起来，\Quote{天主}绝不可能与\Quote{基督}分开。正如宗徒宣告基督是从他们而生的，同样他也断言天主在基督内。基督按肉身被说成是从他们而生的；但同一位被宗徒宣告为\Quote{天主在基督内}。因此他在别处也说：\Quote{因为天主在基督内使世界与自己和好。} \BibleRef{格后 5:19} 把两者分开是绝对不可能的。要么否认基督是从他们而生的，要么承认童贞女所生的是天主在基督内，\Quote{祂是，}正如他所说，\Quote{在万有之上，永远受赞颂的天主。}\par

% structure: p0005 source=h2 target=subsection reason=relative-heading-rank; base=h2

\subsection{第二章}

\Emph{天主的称号在一种意义上给予基督，在另一种意义上给予人。}\par

天主的名称对信徒而言足以表示祂天主性的荣耀，但加上\Quote{在万有之上，永远受赞颂的天主}，便排除了对它的亵渎和曲解，以免某些心怀恶意的人为了贬低祂绝对的天主性，而引用一个事实，即天主有时因恩宠在神圣计划中暂时性地赐予人天主的称号，从而通过不恰当的比较应用于天主，例如天主对梅瑟说：\Quote{我使你在法老面前像神一样，}\BibleRef{出 7:1} 或在此处：\Quote{我曾说你们是神，}在这些地方，这显然是一个因俯就而赐予的称号。因为正如经上所说\Quote{我曾说}，它不是一个显示权能的名称，而更多是说话者所赐予的一个称号。但那段经文，即\Quote{我使你在法老面前像神一样，}也表明了赐予者的权能，而不是受封者的天主性。因为当它说\Quote{我赐予了，}它当然指明了施予的天主的权能，而不是在受封者身上的天主性。但当论及我们的天主和主耶稣基督时，说\Quote{祂是在万有之上，永远受赞颂的天主，}这一事实立刻由这些话语证明，而话语的意义也由所赐的名称显示：因为对于天主圣子而言，天主的名称并非表示因恩宠而得的名分，而是祂真实的本性。\par

% structure: p0008 source=h2 target=subsection reason=relative-heading-rank; base=h2

\subsection{第三章}

\Emph{祂解释宗徒的这句话：\Quote{我们从今以后不再按肉身认识谁}，等等。}\par

同样，这位宗徒说：\BibleQuote{从此以后，我们不按肉身认识谁了；纵使我们曾按肉身认识基督，但如今不再这样认识祂了。}\BibleRef{格后 5:16} 圣经中的一切著作彼此之间是何等美妙地一致，在每一部分也是如此；即使它们不在措辞的\Emph{形式}上对应，却在意义和实质上相符。正如他所说：\Quote{纵使我们曾按肉身认识基督，但如今不再这样认识祂了。}因为当前这段经文印证了上面引用的那节，其中说：\Quote{按肉身说，基督也是从他们来的，祂是在万有之上，永远受赞美的天主。}因为那里写道：\Quote{按肉身说，基督也是从他们来的；}而这里说：\Quote{纵使我们曾按肉身认识基督。}那里说：\Quote{祂是在万有之上，永远受赞美的天主；}而这里说：\Quote{但如今不再按肉身认识基督了。}措辞看起来不同，但其力量和意义相同。因为他在那里宣称按肉身出生的基督是万有之上的天主，也正是他在这里宣称不再按肉身认识的同一位。理由显然在此：即他曾认识按肉身出生的那一位，却承认祂是永远的天主；因此他说不再按肉身认识祂，因为祂是在万有之上，永远受赞美的天主；那里的短语\Quote{祂是在万有之上，永远受赞美的天主}对应于这里的\Quote{我们不再按肉身认识基督}；而这里的\Quote{我们不再按肉身认识基督}暗示了\Quote{祂是永远受赞美的天主}。因此，宗徒的教导仿佛上升到更高的高度，尽管它自始至终一致，却以更明确的陈述支持着完美信德的奥秘，说：\Quote{纵使我们曾按肉身认识基督，但如今不再这样认识祂了}，即，正如从前我们认识祂既是人又是天主，如今却只认识祂是天主。因为当肉身的脆弱终结时，我们在祂内不再认识任何事物，除了天主性的德能，因为祂内的一切都是神圣威能的德能，而人性软弱的缺陷已不复存在。因此，在这段经文中，他已彻底阐述了降生成人的全部奥秘和祂的圆满天主性。因为当他说\Quote{纵使我们曾按肉身认识基督}时，他指的是天主在肉身中诞生的奥秘。但加上\Quote{但如今不再这样认识祂了}，他显明了软弱被除去后的德能。因此，那按肉身的知识涉及祂的人性，而那不再如此的认识涉及祂天主性的光荣。因为说\Quote{我们曾按肉身认识基督}意味着只要那被认识的事物存在。如今我们不再认识它，因为它已不复存在。因为肉身的本性已转化为属灵的本体；先前属于人性的，已全部成为天主的。因此，我们不再按肉身认识基督，因为当身体的软弱被神圣威能吸收时，在那神圣的身体中，没有什么留下可以从中认识肉身的软弱。因此，先前属于双重本体的一切，已归属于单一的德能。因为毫无疑问，那因人性软弱而被钉死的基督，完全因祂天主性的光荣而活着。\par

% structure: p0011 source=h2 target=subsection reason=relative-heading-rank; base=h2

\subsection{第四章。}

\Emph{他从\BibleBook{迦拉达}{书}引出一段经文，表明基督内的肉身软弱被祂的天主性所吸收。}\par

事实上，宗徒在他的全部著作中都声明了这一点，并在写给迦拉达人时美妙地说：\BibleQuote{我保禄作宗徒，不是由于人，也不是借着人，而是由于耶稣基督和圣父。}\BibleRef{迦 1:1} 你看他在与先前和现在的经文时是多么前后一致。因为那里他说：\Quote{如今我们不再按肉身认识基督。}这里他说：\Quote{不是由于人，也不是借着人，而是由于耶稣基督。}显然，他在此处的教导与先前的经文相同。因为当他说他不是被人派遣时，他暗示的是：\Quote{我们曾按肉身认识基督}；所以我不是\Quote{被派遣}的\Quote{由于人}，而是\Quote{由于基督}；因为如果我是由基督派遣的，我就不是由人派遣的，而是由天主派遣的。因为在祂内，天主性完全为自己保留，再没有人的名称的余地。因此，当他说过他是\Quote{不是由于人，也不是借着人，而是由于耶稣基督}派遣时，他正确地补充说：\Quote{和圣父}，从而表明他是由圣父和圣子派遣的；在其中，由于神圣而不可言喻的受生的奥秘，有两位（一位是生者，一位是受生者），但派遣的天主只有一个单一的德能。因此，在说他是由圣父和圣子派遣时，他显示两位的人数，但也教导他们的德能在派遣中是同一的。\par

% structure: p0014 source=h2 target=subsection reason=relative-heading-rank; base=h2

\subsection{第五章。}

\Emph{正如削去基督的天主性是亵渎，同样，否认祂是真人也属亵渎。}\par

但祂说 \Quote{藉耶稣基督，和那使祂从死者中复活的天主圣父。} 这位著名的、可钦敬的导师，知道我们的主耶稣基督必须被宣认为真人，也是真天主，总是如此宣示祂内的天主性光荣，以至于不丧失对降生成人的宣认：藉着真实的降生成人，明确排斥了马西昂的\OriginalTerm{马西昂幻象论}{phantasm of Marcion}，并藉着天主性，排斥了厄俾翁派的\OriginalTerm{厄俾翁派贫乏}{poverty of the Ebionite}，免得因这些邪恶亵渎中的任何一个，我们的主耶稣基督被相信要么是完全无天主之人，要么是无人的天主。因此，宗徒恰当地宣明祂是由天主圣子及天主圣父所派遣，随即加上对主降生成人的宣认，说：\Quote{那使祂从死者中复活的：}清楚教导那从死者中复活的是降生成人的天主的真实身体；与这话一致：\Quote{纵使我们曾按肉体认识基督，}又卓越地加上：\Quote{但如今我们不再这样认识祂了。}因为他说他认识祂内按肉体的事，即祂从死者中复活了；但不再按肉体认识祂，因为当肉体的软弱终结时，他认识祂只存在于天主的德能中。他确是我们主的天主性必须被宣讲的忠实而充分的见证人，他最初蒙召时被从天而来的击打所击倒，不仅心里相信那从死者中复活的主耶稣基督的光荣，而且藉着他肉眼的证据切实确立了其真理。\par

% structure: p0017 source=h2 target=subsection reason=relative-heading-rank; base=h2

\subsection{第六章。}

\Emph{祂从基督向正迫害教会的宗徒显现时，所赐的神视，指出祂内两性的存在。}\par

\Emph{因此}，当他在阿格黎帕王及世上的审判官面前辩论时，他说：\Quote{我带着司祭长的权柄和许可往大马士革去，中午时分，大王，我在路上看见一道光，从天上发来，比太阳还亮，照射在我和同行的人周围。我们都跌倒在地，我听见有声音用希伯来语向我说：\InnerQuote{扫禄，扫禄，你为什么迫害我？你用脚踢刺棒是难的。}我说：\InnerQuote{主，你是谁？}主对我说：\InnerQuote{我是纳匝肋的耶稣，你正在迫害他。}}\BibleRef{宗 26:12} 你看，宗徒说得多么真实：他不再按肉体认识那曾在如此辉煌和威严中看见的主。因为当他俯伏在地，看见那无法忍受的神圣之光的光辉时，随即有声音说：\Quote{扫禄，扫禄，你为什么迫害我？} 当他问那是谁时，主回答说，明确指明祂的位格：\Quote{我是纳匝肋的耶稣，你正在迫害他。} 那么，异端者，我问你，我提醒你：你相信宗徒为自己说的话，还是不相信？或者你认为那不重要，你相信主为祂自己说的话，还是不相信？如果你相信，那就了结了：因为你必然相信我们所相信的。因为我们，如同宗徒，即使曾按肉体认识基督，但如今不再这样认识祂。\Emph{我们}不侮辱基督。\Emph{我们}不把肉身与天主性分开；凡在基督内的，我们都相信是在天主内的。所以，如果你相信我们所相信的，你就必须承认同样的信仰奥迹。但如果你与我们不同，如果你拒绝相信教会、宗徒，甚至天主关于祂自己的见证，那么在这宗徒所见的异象中，请指出哪些是肉身，哪些是天主。因为我在此不能将两者分开。我看见不可言喻的光，我看见不可形容的辉煌，我看见人性的软弱不能承受的光辉，超出凡人眼睛所能承受的，天主的荣耀以不可思议的光照耀。这里哪有分裂和分离的余地？在声音中我们听见耶稣，在威严中我们看见天主。我们怎能不相信在同一（位格的）实体中天主和耶稣存在？但我还愿意就此事与你多说几句。请告诉我，我求你，如果你在如今对公教信仰的迫害中，那曾向无知时的宗徒显现的同一异象向你显现，如果你在毫无防备时那光辉环绕你，那无限之光的荣耀在惊恐和慌乱中打击你，你俯伏在身体和灵魂的黑暗中，内心的无限和不可描述恐惧加剧，——请告诉我，我恳求你；当即刻死亡的恐惧压迫着你，从上而来的荣耀的恐怖压垮你，而且你在心灵的混乱中也听到那些话（你的罪如此应得）：\Quote{扫禄，扫禄，你为什么迫害我？} 当你问是谁时，天上回答：\Quote{我是纳匝肋的耶稣，你正在迫害他，} 你会说什么？\Quote{我不知道，我还不完全相信。我要自己再思考一下，我认为你是谁，从天上说话，以你的天主性的光辉充满我：我能听见你的声音，却不能承受你的光彩。我必须考虑这事，我是否应该相信你：你是基督或是天主。如果你只是天主，是否在基督内？如果你只是基督，是否在天主内？我希望仔细遵守这个区分，彻底考虑我们应当相信你是什么，以及评判你是什么。因为我不希望我的任何职分被浪费。就好像我把你看作人，却又向你献上某些神圣的尊崇。} 如果你那时躺在地上，正如宗徒保禄那时躺着，被神圣之光的辉煌压倒，奄奄一息，你也许会这样说，用尽所有这些愚蠢的唠叨。但我们如何解释另一件事：宗徒选择了另一条路；当他跌倒，颤抖而且半死时，他不认为他应该再隐瞒他的信仰，或迟疑；他已经被不可言说的论据所教导，知道那他曾无知地以为是人的那位，是天主。他没有隐瞒信仰，他没有延迟。他不再通过怀疑和不信延长他的错误想法，而是当他从天上听见他主耶稣的名字时，他用卑微如仆人的声音、颤抖如被鞭打的声音、充满热情如归化者的声音回答说：\Quote{主，我该作什么？} 所以立刻，由于他迅速而热切的信心，他被赐予永不离开他忠信地相信的那位；而且他所全心投靠的那位，应该自己进入他的心中：正如宗徒自己论到自己说：\BibleQuote{你们寻求基督在我内说话的证明吗？}\BibleRef{格后 13:3}\par

% structure: p0020 source=h2 target=subsection reason=relative-heading-rank; base=h2

\subsection{第七章}

\Emph{他再次藉宗徒的其他段落证明基督是天主。}\par

我要你告诉我，你这异端者，在这段经文中，那位——如宗徒告诉我们——在他内说话的是人还是天主？如果他是人，另一个人的身体怎能在他心里说话？如果他是天主，那么基督就不是人而是天主；因为既然基督在宗徒内说话，而只有天主能在他内说话，所以是神圣的基督在他内说话。因此你看到这里无话可说，不能在天主与基督之间作任何分开或分离：因为完满的天主性在基督内，基督也完全在天主内。这里不能容许两者有任何分开或割裂。只有一个简单、虔诚、纯正的宣认：即朝拜、爱慕、尊崇基督为天主。但你是否想要更充分、更透彻地明白天主与基督之间不能分开，我们必须认为天主与基督完全是一体？请听宗徒对格林多人说的话：\BibleQuote{因为我们众人都必须出现在基督的审判座前，好使每人因身体所行的事，或善或恶，领受相应的报应。}\BibleRef{格后 5:10}但在另一处，他写信给罗马人说：\BibleQuote{我们众人都要站在天主的审判座前，因为经上记载：\InnerQuote{我指着我自己起誓——主说——万物都向我屈膝，万口都要承认归于天主。}}\BibleRef{罗 14:10}你看，天主的审判座与基督的审判座是同一的；你当毫无疑虑地明白基督就是天主；当你看到天主与基督的本质完全不可分时，也要承认位格是不能割裂的。否则，因为宗徒在一封书信中说我们要出现在基督的审判座前，在另一封中说出现在天主的审判座前，你就臆造两个审判座，幻想有些人由基督审判，另一些人由天主审判。但这是愚蠢、狂妄，比疯子的胡言还要疯狂。你要承认万有的主，宇宙的天主，在基督的审判座中承认天主的审判座。要爱生命，爱你的救恩，爱那创造你的祂。要惧怕那审判你的祂。因为无论你愿意与否，你都必须出现在基督的审判座前；尽管你认为天主的审判座与基督的不同，你仍将来到基督的审判座前，并且通过无法反驳的明证发现，天主的审判座的确与基督的审判座相同，并且在基督——圣子——内，有圣子的一切荣耀和天主圣父的权能。\BibleQuote{因为圣父不审判任何人，而把审判全权交给了圣子，为使众人尊敬圣子如尊敬圣父。}\BibleRef{若 5:22}因为凡否认圣父的，也否认圣子。\BibleQuote{凡否认圣子的，也没有圣父；那承认圣子的，也有圣父。}\BibleRef{若一 2:23}所以你们应当知道，圣父与圣子的荣耀是不可分的，他们的尊威也是不可分的，并且不能尊崇圣子而不尊崇圣父，也不能尊崇圣父而不尊崇圣子。但没有人能在基督——天主的独生子——之外尊崇天主和圣子。因为人若不在基督的圣神内，不可能拥有那应受尊崇的天主圣神，正如宗徒所说：\BibleQuote{你们并不属于血肉，而是属于圣神，只要天主的圣神住在你们内。谁如果没有基督的圣神，就不属于祂。}\BibleRef{罗 8:9}又说：\Quote{谁能控告天主的选民？天主是使他们成义的。谁能定他们的罪？是那死过、更好说已复活的基督耶稣。}你现在看到，即使违背你的意愿，天主的圣神与基督的圣神之间，或天主的审判与基督的审判之间，绝对没有任何区别。那么选择你愿意的——二者必居其一——要么在信德中承认基督是天主，要么承认天主在基督内，而你要被定罪。\par

% structure: p0023 source=h2 target=subsection reason=relative-heading-rank; base=h2

\subsection{第八章。}

\Emph{当宣认基督的天主性时，我们不应默不作声地略过那十字架的宣认。}\par

但让我们看看还有什么下文。在写给格林多教会的信中，我们上面提到的那位，所有教会的导师，即\Person{保禄}，这样说道：\BibleQuote{犹太人要求神迹，希腊人寻求智慧，而我们所宣讲的，却是被钉在十字架上的基督；这为犹太人固然是绊脚石，为外邦人是愚妄，但为那些蒙召的，不拘是犹太人或希腊人，基督却是天主的德能和天主的智慧。}\BibleRef{格前 1:22}哦，最有能力的信仰导师，他甚至在这段经文里，在教导教会时，认为仅说基督是天主还不够，还必须加上祂被钉十字架，为的是借着公开而稳固地教导信仰，把他称为被钉十字架的那位宣示为天主的智慧。他随后没有使用任何诡辩或迂回之辞，当他宣讲主的福音时，也毫不以提及基督的十字架为耻。尽管对犹太人来说是绊脚石，对外邦人是愚妄，因为他们听闻天主诞生、天主具有形躯、天主受苦、天主被钉十字架，但他并没有因犹太人冒犯的邪恶而削弱他虔诚言词的力度；也没有因别人的不信和愚妄而减少他信德的活力；而是公开、持续、大胆地宣布：那位由母亲所生、被人杀害、长枪刺透、十字架伸展的——正是\BibleQuote{天主的德能和智慧，对犹太人来说是绊脚石，对外邦人是愚妄。}但那些对某些人是绊脚石和愚妄的，对另一些人却是天主的德能和智慧。因为正如人不同，他们的思想也不同：一个缺乏健全理解、不能认识真善的人，在不信中愚昧地否认的，一个有智慧的信德却能在灵魂深处感到那是神圣而赋予生命的。\par

% structure: p0026 source=h2 target=subsection reason=relative-heading-rank; base=h2

\subsection{第九章。}

\Emph{宗徒的宣讲如何因承认被钉的基督是天主，而遭犹太人和外邦人拒绝。}\par

那末，请你告诉我，你这个异端者，你是一切人的仇敌，但更是你自己的仇敌——对我们主\Person{耶稣基督}的十字架，你如同犹太人视为绊脚石，又如同外邦人视为愚妄；你以犹太人的失足拒绝真实救恩的奥秘，并以外邦人的顽固显现愚昧——为何宗徒\Person{保禄}的宣讲在外邦人看来是愚妄，在犹太人看来是绊脚石呢？的确，如果他像你所主张的那样教导基督只是一个凡人，那他的宣讲决不会冒犯人。因为谁会认为他的诞生、受难、十字架和死亡是难以置信或难以接受的呢？如果\Person{保禄}说的只是一个凡人的基督遭受了凡人在世上常受的苦难，那他的宣讲又有什么新奇或特异之处呢？但正是这一点，外邦人的愚妄不能接受，犹太人的不信予以拒绝，即：宗徒宣称那位像你一样被幻想为凡人的基督，是天主。正是这邪恶之人的思想所拒绝、不信之人的耳朵所不能忍受的：即应在人\Person{耶稣基督}中宣告天主的诞生、主张天主的受难、宣讲天主的十字架。这就是困难所在！这就是难以置信的！因为人的耳朵从未听说过发生在天主性上的事，这对他们来说是不可信的。因此，你的宣讲和教导是相当安全的：你的宣讲永远不会成为外邦人的愚妄，也不会成为犹太人的绊脚石。你永远不会与\Person{伯多禄}一同被犹太人和外邦人钉死在十字架上，也不会与\Person{雅各伯}一同被石头砸死，也不会与\Person{保禄}一同被斩首。因为你的宣讲中没有任何东西会冒犯他们。你坚持一个凡人出生，一个凡人受难。你无需担心他们会用迫害来困扰你，因为你的宣讲正在帮助他们。\par

% structure: p0029 source=h2 target=subsection reason=relative-heading-rank; base=h2

\subsection{第十章。}

\Emph{宗徒如何坚持基督是天主的德能和天主的智慧。}\par

但是，让我们在这个主题上再看一些东西。那么，基督按照宗徒的教导，是\OriginalTerm{天主的德能}{power of God}和\OriginalTerm{天主的智慧}{wisdom of God}。你对此有何话可说？你怎能脱身？你没有地方可以逃避。基督是天主的智慧和天主的德能。我说，祂就是犹太人攻击、外邦人嘲弄、你自己也与他们一同迫害的那一位——我说，祂为外邦人是愚妄、为犹太人是绊脚石、也为你们二者，我说，祂就是\OriginalTerm{天主的德能}{power of God}和\OriginalTerm{天主的智慧}{wisdom of God}。你能做什么？你难道要掩耳不听吗？当宗徒讲道时，犹太人也如此行。你尽管做你想做的，基督在天上，在天主内，与祂同在，在祂内，在高天之上，祂也曾在此低处；你不能再与犹太人一起迫害祂。但你只能做一件事。你在信仰中迫害祂，你在教会中迫害祂，你用邪恶信念的武器迫害祂，你用虚假教义的刀剑迫害祂。也许你比往日的犹太人做得更多。你现在迫害基督，而在那之后，曾迫害祂的人已经相信了。但也许你认为因为不能再下手抓祂，罪就更轻了。我告诉你，那种迫害毫不减轻，毫不减轻对祂的伤害，即罪人在祂的追随者身上迫害祂。但提及主的十字架使你反感。它同样常常惹怒犹太人。你战栗地听到天主受难：外邦人也在他们的错误中嘲弄这点。那么我问你，你们在这一点上有什么不同？既然你们在这悖逆上意见一致。但就我而言，我不仅不淡化这神圣十字架的宣讲，这主受难的宣讲，反而尽我所能地强调它。因为我要宣告：被钉十字架的那位不仅比任何事物都大的\OriginalTerm{天主的德能}{power of God}和\OriginalTerm{天主的智慧}{wisdom of God}，而且实际是\OriginalTerm{绝对天主性与光荣的主}{Lord of absolute Divinity and glory}。而且更是如此，因为我的这个断言是天主的道理，如宗徒所说：\BibleQuote{我们在成全的人中讲智慧，但不是这世代的智慧，也不是这世代将要消灭的执政者的智慧；我们讲的乃是天主奥秘中隐藏的智慧，就是天主在万世以前预定使我们得荣耀的，这智慧是这世代执政者没有一个知道的；因为他们若知道，就决不会把光荣的主钉在十字架上了。正如经上所记：眼所未见，耳所未闻，人心所未想到的，就是天主为爱祂的人所预备的。}\BibleRef{格前 1:6} 你看到宗徒的论述在多么小的范围内包含了多么伟大的事。他说他讲智慧，但智慧只是那些成全的人所能认识的，而此世的聪明人不能认识。因为他说这是天主的智慧，隐藏在天主奥秘中，在万世以前预定为圣徒得荣耀；因此只有那些尝到天主滋味的人才知道，而此世的执政者完全不知道。但他加上了理由，以确定他所提的两点，说：\BibleQuote{因为他们若知道，就决不会把光荣的主钉在十字架上了。正如经上记载：眼所未见，耳所未闻，人心所未想到的，就是天主为爱祂的人所预备的。} 那么你看，天主的智慧隐藏在天主奥秘中，在万世以前预定的，被那些钉死\OriginalTerm{光荣的主}{Lord of glory}的人所不知道，而被那些接受的人所知道。他说天主的智慧隐藏在天主奥秘中，这说得对，因为人的眼睛从未见过，耳朵从未听过，心也从未想到过这一点：即\OriginalTerm{光荣的主}{Lord of glory}应从童贞女出生，降生肉体，受尽各种刑罚和可耻的苦难。但关于这些天主恩赐，既然它们隐藏在天主奥秘中，没有人能凭自己理解它们，所以当它们被启示时，抓住它们的人是有福的。因此，所有未能抓住它们的人必须被算在此世的执政者中，而抓住他们的人算在天主的智慧人中。那么，否认天主降生在肉体内的人就没有抓住它；因此你也没有抓住它，因为你否认这点。但你尽管做你想做的，尽可能不虔地否认，我们却更相信宗徒。但我何必说宗徒呢？我们更相信天主。因为通过宗徒我们相信祂，我们知道祂曾通过宗徒说话。天主圣言说，\OriginalTerm{光荣的主}{Lord of glory}被此世的执政者钉死在十字架上。你否认。那些钉死祂的人也否认他们钉死的是天主。那么，承认祂的人就与承认祂的宗徒同份。你必定与祂的迫害者同份。那么，对此还能回答什么？宗徒说\OriginalTerm{光荣的主}{Lord of glory}被钉在十字架上。你若有本事，就改变这点。现在，如果你愿意，将耶稣与天主分开。至少你不能否认基督被犹太人钉在十字架上。但被钉的是\OriginalTerm{光荣的主}{Lord of glory}。因此，你要么否认基督被钉在十字架上，要么必须承认天主被钉在上面。\par

% structure: p0032 source=h2 target=subsection reason=relative-heading-rank; base=h2

\subsection{第十一章}

\Emph{他通过福音的证明支持同样的道理。}\par

但也许你难以接受的是，我一直主要使用宗徒保禄一人的见证。对我而言，他足够了，是天主所拣选的，我也不耻于将天主愿意他为普世导师的人作为我信德的见证。但为了满足你的愿望，也许你以为我没有其他证据可用，请听玛尔大在福音中宣告的人类救恩和永福的完美奥秘。因为她说了什么？\BibleQuote{主，我确信祢是基督、永生天主之子，来到这世界上的那位。}\BibleRef{若 11:27} 从一位妇人学习真信仰。学习永生的宣认。然而你有一个极好的安慰：你无需羞于由她来教导救恩的奥秘，因为天主没有拒绝接受她的见证。\par

% structure: p0035 source=h2 target=subsection reason=relative-heading-rank; base=h2

\subsection{第十二章}

\Emph{通过有福的伯多禄的著名宣认，他证明基督是天主。}\par

但是，倘若你更愿采纳一位更伟大之人的权威（尽管你不该轻视任何一位男性或女性的权威，因为对奥迹的宣认本身就赋予他们分量——无论一个人的地位如何，或他的处境多么卑微，其信德的价值都不会因此减损），那么，让我们不要去请教什么初学者或未受教的学童，也不要请教一位其信德或许还显得粗浅的妇女；而要去请教那群门徒中最伟大的门徒、众教师中最伟大的教师，他曾主持并治理罗马教会，在司铎职中居首位，正如他在信德中居首位一样。那么，请告诉我们，告诉我们，我们恳求你，啊，\Person{伯多禄}，宗徒之长，请告诉我们，众教会该如何相信天主。因为你理当教导我们，正如主教导了你；你理当为我们开启那扇你已得到钥匙的门。将所有试图推翻这天主居所的人拒之门外；也拒绝那些企图从秘密洞口和非法途径进来的人：因为显然，除非那被授予众教会的钥匙由你向人揭示，否则无人能进入天国之门。那么，请告诉我们，我们该如何相信耶稣基督、宣认我们共同的主。你必定会毫不犹豫地回答：\Quote{你们为何就该如何宣认主来咨询我呢？你们面前不是已有我对祂的宣认吗？你们读一读福音，得到了我的宣认，也就不再需要我本人了。不，你们得到我的宣认时，也就得到了我本人；因为虽然我离开我的宣认就没有分量，但那宣认本身却增添了我这人的分量。}\newline{} 那么，请告诉我们，啊，圣史，请告诉我们那宣认：请告诉我们那首席宗徒的信德：他宣认耶稣仅仅是一个人，还是天主？他说祂内只有血肉，还是宣告祂是天主子？当时主耶稣基督问门徒们以为祂是谁、宣认祂是谁时，宗徒之长伯多禄代表众人回答——因为一人的回答与所有人的信德一致。而由他先作回答本是合宜的，好使回答的顺序与荣誉的等级相应：且使他在宣认上超越众人，正如他在年龄上超越众人一样。那么他说了什么？\Quote{祢是，}他说，\Quote{基督，永生天主之子。}\BibleRef{玛 16:16} 异端者啊，我不得不利用一个简单明了的问题来驳斥你。请告诉我，请问，伯多禄那样回答时，对方是谁？你无法否认那就是基督。那么我问你，你称基督为什么？人还是天主？当然毫无疑问是人：因为你整个异端的根源就在这里——你否认基督是天主子。同样，你也说玛利亚是\OriginalTerm{基督之母}{Christotocos}，而不是\OriginalTerm{天主之母}{Theotocos}，因为她是基督的母亲，而非天主的母亲。因此你主张基督仅仅是一个人，而不是天主，所以祂是人子，而非天主子。那么，伯多禄对此如何回答？\Quote{祢是，}他说，\Quote{基督，永生天主之子。}你所宣告的仅仅是人之子的那个基督，他见证其为天主子。你愿意我们相信谁？你还是伯多禄？我想你不至于无耻到敢于将自己的意见凌驾于宗徒之长之上吧。然而还有什么你不敢尝试的呢？如果你能否认天主，你又怎能不藐视宗徒呢？\Quote{那么祢是，}他说，\Quote{基督，永生天主之子。}这里面有什么令人困惑或隐晦的吗？这只是简单明了的宣认：他宣告基督是天主子。也许你会否认这话说过：但圣史见证它们说过。或者你说宗徒撒谎了？但指控宗徒撒谎是可怕的谎言。或者你会坚持认为这些话是针对另一位基督说的？但这是一种新奇而怪诞的捏造。那么你还有什么可说的呢？确实只剩一件事：既然所写的内容被人诵读，所诵读的内容是真实的，你最终应当被迫（因为你不能断言其虚假）停止抨击真理。\par

% structure: p0038 source=h2 target=subsection reason=relative-heading-rank; base=h2

\subsection{第十三章}

\Emph{真福伯多禄的宣认从基督本人得到对其真理的见证。}\par

但是，既然我已采用了首席宗徒的见证，他在其中公开宣认了主耶稣基督为天主，让我们看看祂（即他所宣认的那位）是如何认可他的宣认的；因为比宗徒的话更有价值的是，天主亲自称赞了他的话。当时宗徒说了\Quote{祢是基督，永生天主之子}，我们的主救主的回答是什么？他说：\Quote{\Person{西满巴约纳}，你是有福的，因为不是血肉启示了你，而是我在天之父的圣神。}如果你不喜欢用宗徒的见证，那就用天主的见证吧。因为天主称赞了所说的话，从而把祂自己的权威加给了宗徒的话，这样，虽然话出自宗徒之口，然而认可这话的天主却使之成为祂自己的。他说：\Quote{\Person{西满巴约纳}，你是有福的，因为不是血肉启示了你，而是我在天之父的圣神。}因此，在宗徒的话中，你有了圣神、临在的圣子以及天主圣父的见证。你还能想要什么，或者还有什么可与之相比？圣子称赞，圣父临在，圣神启示。宗徒的话因此给出了整个天主性的见证：因为这话必具有那推动它所言及者的权威。他说：\Quote{那么\Person{西满巴约纳}，你是有福的，因为不是血肉启示了你，而是我在天之父的圣神。}如果血肉没有启示给伯多禄这话或激励他，那么你最终必须看看是谁在启示你。如果天主之神教导了他，使他宣认基督是天主，那么你若是能否定这一点，你就看出你是受魔鬼之神教导的。\par

% structure: p0041 source=h2 target=subsection reason=relative-heading-rank; base=h2

\subsection{第十四章}

\Emph{真福伯多禄的宣认如何是普世教会的信德。}\par

然而，在主的那句话之后，接着还有什么话，祂藉此称赞伯多禄？祂说：\Quote{我也告诉你，你是伯多禄，在这磐石上，我要建立我的教会。}你看见了吗？伯多禄的话就是教会的信德。因此，凡不持守教会信德的，当然是在教会之外。主又说：\Quote{我要将天国的钥匙交给你。}这信德堪得天国；这信德领受了天国的钥匙。看，等待你的是什么。你若否认这钥匙的信德，就不能进入这钥匙所属的门。祂又加上说：\Quote{地狱的门也不能胜过你。}地狱的门就是异端者的信念，或更确切地说，错谬的信念。因为地狱与天堂相隔多远，否认基督是天主的人与宣认基督是天主的人就相隔多远。祂继续说：\Quote{凡你在地上捆绑的，在天上也要被捆绑；凡你在地上释放的，在天上也要被释放。}宗徒的完美信德，仿佛被赋予了神性的权能，使他在凡地上捆绑或释放的，在天上也得以捆绑或释放。至于你，你反对宗徒的信德，既然你已经看见自己在世上已被捆绑，剩下的只是让你知道你在天上也被捆绑。但要详细叙述这些事，实在太多，即使简明扼要地说，也会成为一个冗长乏味的故事。\par

% structure: p0044 source=h2 target=subsection reason=relative-heading-rank; base=h2

\subsection{第十五章}

圣多默在主复活后也宣认了与伯多禄相同的信德。\par

然而，我还要为你加上一位宗徒的另一个见证，使你能看到，苦难之后所发生的事与先前的事如何相符。当主在门关闭时出现在众门徒中间，并愿向宗徒们证明祂身体的真实性时，当宗徒多默触摸了祂的肉体，探了祂的肋旁，检查了祂的伤口时——在他确信所显给的身体的真实性后，他宣告了什么？他说：\Quote{我的主，我的天主！}\BibleRef{若 20:28}他可曾像你所说的那样，说那是一个凡人而不是天主？是基督而不是神性？他确实触摸了他主的身体，并回答说他就是天主。他是否在人与天主之间作了任何分离？或者，他是否用你的说法，称那肉体为\OriginalTerm{天主之母}{Theotocos}，即那接受了神性的？或者，他是否按照你亵渎的方式，宣称他所触摸的那一位，不应因其自身被尊敬，而应因那接受在他里面的那一位？但也许天主的宗徒不知道你那微妙的分离，也未曾体验过你那判断的精微区分，因为他是一个粗鄙的乡下人，不懂辩证法，也不懂哲学辩论的方法；对他来说，主的教训已十分充足，他除了从主的教导中所学的之外，一无所知！所以，他的话含有天上的教义；他的信德是神圣的教训。他从未学会像你那样将主与祂的身体分离，也不知道如何将天主与祂自己撕开。他是圣洁、正直、公义的：充满实际的纯朴、纯全的信德和洁净的知识：拥有单纯的理解与明智的结合，完全远离一切邪恶的智慧，加上完美的纯朴：不知任何败坏，远离一切异端的乖谬，并因亲身体验了神圣教训的力量，他持守了他所学习的一切。因此，他——照你所想的，那乡下人和无知的人——以一句简短的回答堵住了你的口，并以他的寥寥数语摧毁了你的立场。那么，当宗徒多默走近去触摸他的天主时，他触摸了什么？无疑是基督，毫无可疑。但他喊了什么？他说：\Quote{我的主，我的天主！}现在，你若能，就将基督与天主分开，并改变这句话，如果你能做到的话。利用一切辩证技巧——世上的一切聪明，以及那在于词语微妙之处的愚拙智慧。四处转动，缩回你的角。用尽机巧和艺术，做你能做的一切。说你喜欢的话，做你喜欢的事；你不可能逃脱这一点，而不承认宗徒所触摸的是天主。事实上，如果这是可能的话，你也许还想篡改福音故事的叙述，使我们不再读到宗徒多默触摸了主的身体，或他把基督称为主和天主。但要篡改写在天主的福音中的话，是绝对不可能的。因为\BibleQuote{天地要过去，但天主的话决不会过去。}\BibleRef{玛 24:35}看，如今那当时作证的宗徒多默，仍然向你宣告：\Quote{我所触摸的耶稣是天主。我触摸了祂的四肢，那是天主。我没有感受到非形体的东西，也没有触摸到不可触摸的东西：我不是用手触摸了一个神体，以至于人们可相信我只针对它说\InnerQuote{这是天主}。因为\InnerQuote{神体}，正如我主自己所说，\BibleQuote{没有血肉和骨头}。\BibleRef{路 24:39}我触摸了我主的身体。我触摸了血肉和骨头。我将手指探入钉痕中：并对我所触摸的基督我主宣告：\BibleQuote{我的主，我的天主！}\BibleRef{若 20:28}因为我不知如何将基督与天主分离，也不能在耶稣与天主之间插入亵渎的区分，或将我主与祂自身撕裂。远离我吧，凡持不同意见的，或说不同话的。我不知道基督不是天主。这信德我与我同伴的宗徒们一同持守；我传给了众教会；我向外邦人宣讲；我现在也向你宣告：基督是天主，基督是天主。健全的理性不猜想别的；健全的信德不说的别的。神性不能与自身分离。既然基督的一切都是天主，在天主内除了天主，再也找不到别的。}\par

% structure: p0047 source=h2 target=subsection reason=relative-heading-rank; base=h2

\subsection{第十六章}

他引出天主圣父为圣子神性所作的见证。\par

异端者啊，你现在怎么说？这些信仰的明证——也即你一切不信的明证——够了吗？或者你还要再加一些？但在先知和宗徒之后，还能加什么呢？除非——如同犹太人曾要求的那样——你也要一个从天上来的征兆？你若要求，我们必给你从前给他们的同样回答：\BibleQuote{一个邪恶淫乱的世代求征兆，除了先知约纳的征兆外，再没有征兆给它。}\BibleRef{玛 16:4} 其实，这个征兆对你已足够，如同对钉死祂的犹太人一样，使你仅凭此就学会信靠主天主——连那些迫害祂的人最终也因此相信了。但既然我们提到了天上的征兆，我要给你一个天上的征兆：一个连魔鬼也从未否认过的征兆；它们虽看见耶稣的形体，却被真理所迫，呼喊祂是天主——祂确实是。那么，福音作者论及主耶稣基督说了什么？\BibleQuote{祂受了洗，}他说，\BibleQuote{立刻从水里上来。忽然天为祂开了，祂看见天主圣神如同鸽子降下，来到祂身上；又有声音由天上说：这是我的爱子，我所喜悦的。}\BibleRef{玛 3:16} 对此你怎么说，异端者？你不喜欢所说的话，还是说话者的位格？话语的含义无需解释，说话者的价值也无需言语称赞。是圣父天主说的。祂所说的话十分清楚。你总不能无耻到亵渎地断言：关于天主独生子，圣父天主是不可信的？\BibleQuote{这是，}祂说，\BibleQuote{我的爱子，我所喜悦的。} 但也许你会坚持说这是疯狂，这话是对圣言说的，不是对基督说的。那么告诉我，谁受了洗？是圣言，还是基督？是肉身，还是圣神？你总不能否认是基督。那么，那从人和天主所生、因圣神降临于童贞女并至高者能力的庇荫而受孕的，既是人子又是天主子——你不能否认，就是他受了洗。如果受洗的是他，那么被指的也是他，因为受洗者正是被指的那位。\BibleQuote{这是，}祂说，\BibleQuote{我的爱子，我所喜悦的。} 还有什么能比这更明确、更清晰呢？基督受了洗。基督从水里上来。基督受洗时天开了。为了基督，鸽子降临在基督身上，圣神以形体出现。圣父对基督说话。你若敢否认这话是针对基督说的，那你就得坚持基督没有受洗、圣神没有降临、圣父没有说话。但真理本身催逼着你，压着你，即使你不愿承认，也无法否认。福音作者怎么说？\BibleQuote{祂受了洗，立刻从水里上来。}谁受了洗？肯定是基督。\BibleQuote{忽然，}他说，\BibleQuote{天为祂开了。}为谁？当然是为受洗者。肯定是基督。\BibleQuote{祂看见天主圣神如同鸽子降下，来到祂身上。}谁看见了？确实是基督。圣神降在谁身上？肯定是基督。\BibleQuote{又有声音由天上说——}关于谁？确实是基督：因为接着是什么？\BibleQuote{这是我的爱子，我所喜悦的。} 为了表明这一切是因谁而发生，随后有声音说：\BibleQuote{这是我的爱子，}如同说：这是为祂发生这一切的那位。因为这是我的儿子：为了祂，天开了；为了祂，我的圣神来了；为了祂，我的声音被听见。因为这是我的儿子。那么，当祂说\BibleQuote{这是我的儿子}时，祂指的是谁？当然是鸽子所触的那位。鸽子触了谁？确实是基督。因此基督是天主子。我想，我的许诺实现了。异端者啊，你现在看见从天上给你的征兆了吗？不只是一个，而是许多特别的征兆？因为有一个是天的开启，另一个是圣神的降临，第三个是圣父的声音。这些都最清楚地表明基督是天主：天的开启显示祂是天主，圣神降临在祂身上支持祂的天主性，圣父的宣告确认了它。因为天不会为它的主以外的人开启；圣神也不会以形体降在除天主子以外的人身上；圣父也不会宣称祂是儿子，除非祂真是——尤其有神圣诞生的这些标志，不仅证实正确信仰的真理，也排除了有罪和错误信仰的邪恶。因为圣父以神圣言语不可言喻的威严明确而直接地说出\BibleQuote{这是我的儿子}后，又加上了下面的话——即\BibleQuote{我的爱子，我所喜悦的。} 正如祂已藉先知宣告祂是全能的天主和伟大的天主，所以祂在这里说\BibleQuote{我的爱子，我所喜悦的}时，又在自己儿子的名上加添了爱子的称号——祂所喜悦的那位：好让这些称号的附加表明天主性本性的特殊属性，并使这从未发生在任何人身上的事特别归于天主子的光荣。因此，如同在我们的主耶稣基督身上发生了这些特别而唯一的事：即天开了，在众人面前圣父天主藉鸽子的来临和临在以一种方式触摸了祂，几乎用手指指着祂说\BibleQuote{这是我的儿子；}同样，在祂身上也是特别而唯一的：即祂是特别被爱的，并被特别称为圣父所喜悦的，好让这些特别伴随之物标示祂本性的特殊意义，并且祂名字的特殊性支持独生子的特殊地位——这地位已由先前所赐的征兆的尊荣所证实。但这里到了这卷书的结尾。因为圣父天主的这个言语，人的任何言语都不能增添，也不能相比。对我们而言，圣父天主自己就是关于我们的主耶稣基督足够令人满意的见证，当祂说\BibleQuote{这是我的儿子}时。你若认为圣父天主的这些话可能被反驳，那你就被迫反对祂——祂以最明确的宣告使祂被全世界承认为自己的儿子。\par

\SourceRecord{Translated by C.S. Gibson.；Nicene and Post-Nicene Fathers, Second Series；Vol. 11.；Edited by Philip Schaff and Henry Wace.；Buffalo, NY: Christian Literature Publishing Co.,；1894.；<http://www.newadvent.org/fathers/35093.htm>.}

% structure: p0001 source=h1 target=section reason=page-title

\section{\WorkTitle{论降生成人}（卷四）}

% structure: p0002 source=h2 target=subsection reason=relative-heading-rank; base=h2

\subsection{第一章}

\Emph{基督在降生成人以前是从永远就存在的天主。}\par

既然我们已经用最确实、最宝贵的见证者完成了三卷书，他们的真理不仅有人的证据证实，也有神圣的证据证实，这些就已足够用神圣的权威来证明我们的论点，何况论点本身的神圣权威对此已足矣。然而，由于整部圣经充满了这些证据，而且哪里有如此多的见证者，哪里就有如此多的意见可以提出——不，其实圣经本身以同一神圣的口吻，仿佛在作证——我们认为最好还是再添加一些其他证据，并非出于需要确认，而是因为我们拥有丰富的材料；这样，一些在辩护上可能不必要的东西，在装饰上或许也有用。因此，在前几卷中，我们不仅凭先知和宗徒的证据，也凭福音作者和天使的证据，证明了我们的主耶稣基督在肉身中的天主性，那么现在让我们表明：那位在肉身上出生者，甚至在降生成人以前就是天主；这样你们就能通过圣经证据的和谐一致来明白，你们应当相信，祂在肉身出生时既是天主又是人，祂在出生前只是天主，并且那位由童贞玛利亚在肉身中生出后仍是天主的，在从童贞玛利亚出生前，就是天主圣言。那么首先从全世界的导师——宗徒那里学习一个道理：那位无始的、天主、天主子，在世界的末了，即时期的圆满，成了人子。因为他说：\BibleQuote{但时期一满，天主就派遣了自己的圣子，由女人所生，生于法律之下。}\BibleRef{迦 4:4}那么请告诉我，在主耶稣基督由祂的母亲玛利亚出生以前，天主有圣子还是没有？你不能否认祂有，因为从来没有儿子没有父亲，也没有父亲没有儿子：因为正如儿子的称呼是相对于父亲，父亲的称呼也是相对于儿子。\par

% structure: p0005 source=h2 target=subsection reason=relative-heading-rank; base=h2

\subsection{第二章}

\Emph{他从上述的话推论：童贞玛利亚所生的儿子是先行存在的，且比她本人更大。}\par

那么你看，当宗徒说天主派遣了祂的圣子时，天主派遣的是祂自己的圣子——用宗徒的原话就是\Quote{祂自己的圣子}。因为，既然天主派遣的是祂自己的圣子，那么祂派遣的不是别人的圣子；而且，如果被派遣者不存在，祂根本不可能派遣祂。所以，他说，祂派遣了\Quote{祂自己的圣子，由女人所生}。因此，因为祂派遣了祂，祂派遣了一位存在者；又因为祂派遣的是祂自己的，那么肯定不是别人的，而是祂自己的。那么，你们从世俗的琐碎中得出的那个论证又有什么意义呢？从来没有人给一位已经存在者生育。难道主在玛利亚之前没有先行存在吗？天主子不是在人之女之前就已存在吗？总之，难道天主本身不是在人之先就已存在吗？——因为确实，没有一个人不是来自天主的。那么你看，我不仅说玛利亚生了一位在她之前已经存在者，我不仅说一位在她之前已经存在者，而且说一位是她存在的创造者；这样，在生育她的创造者时，她成了给予她存在的那位的母亲。因为天主为自己带来出生，与人一样简单；祂安排自己由人类所生，如同一个人出生一样容易。因为天主的权能并不限于祂自己的位格，仿佛对所有人都是允许的事，在祂自己身上却是不允许的；仿佛那在神圣本性中作为天主能做一切事的那一位，在祂自己的位格中却不能成为在人之中的天主。所以，放下并拒绝你们那些从世俗事物中得出的愚蠢、软弱和迟钝的论证，我们只应相信直接明了的证据和赤裸裸的真理，并调整我们的信德，只依靠天主所派遣的那些见证者，并且可以说，天主亲自在他们的位格中宣讲。因为，在关于认识祂自己的事情上，相信祂是正当的，因为我们对祂所知的一切，都来自祂自己；因为天主不可能被人认识，除非祂亲自把对祂自己的认识赐给我们。因此，我们应当相信我们对祂所知的一切，因为我们对祂所知的一切都来自祂；因为如果我们不相信我们的知识所从来的那一位，结果就是我们什么都不会知道，因为我们拒绝相信我们通过祂获得知识的那一位。\par

% structure: p0008 source=h2 target=subsection reason=relative-heading-rank; base=h2

\subsection{第三章}

\Emph{他由\WorkTitle{罗马书}证明基督永恒的天主性。}\par

由上述证词清楚可见，天主派遣了自己的圣子，那永恒的天主圣子成了人子。现在让我们看看同一宗徒是否在其他地方也给出了同类证词，好使那本身已足够清楚的事实，借着双重证词的光照更加显明。同一宗徒说：\Quote{天主派遣了自己的圣子，带着罪恶肉身的模样。}\BibleRef{罗 8:3}你看，宗徒绝非偶然或随意使用这些话，因为他重复了已经说过的话——诚然，在他内住满了神圣的谋划和言语，绝无偶然或轻率。那么他说了什么？\Quote{天主派遣了自己的圣子，带着罪恶肉身的模样。}他又一次说并重复：\Quote{天主派遣了自己的圣子。}啊，卓越而杰出的导师！因为他知道，其中包含了公教信仰的全部奥秘，为使人相信主降生在肉身中，天主子被派遣到这世界，他便一再宣告同一信息，说：\Quote{天主派遣了自己的圣子。}我们无需惊讶，那位专为宣讲天主的来临而被派遣的人如此宣告，因为甚至在法律之前，法律本身的颁布者也宣告了这事，说：\Quote{我恳求祢，上主，请另派一位祢所愿派遣的人。}或者，按希伯来文本更清楚的说法：\Quote{我恳求祢，上主，请派遣祢所愿派遣的人。}显然，圣先知在自己内感受到对全人类的渴望，他以全人类的声音向天主圣父祈祷，恳求祂尽快派遣那将由圣父为所有人的救赎和救恩而派遣的一位，他说：\Quote{我恳求祢，上主，请派遣祢所愿派遣的人。}因此他说：\Quote{天主派遣了自己的圣子，带着罪恶肉身的模样。}当他宣称祂被派遣在肉身中时，极好地排除了肉身的罪恶：因为他说\Quote{天主派遣了自己的圣子，带着罪恶肉身的模样}，为使我们知道，虽然肉身真实地被取了，却没有真实的罪恶，就身体而言，我们应理解其真实性；就罪恶而言，仅有罪恶的模样。因为虽然所有肉身都有罪，但祂却有毫无罪恶的肉身，并在自身内具有罪恶肉身的模样，祂在肉身中，却免于真实的罪恶，因为祂是无罪的：因此他说：\Quote{天主派遣了自己的圣子，带着罪恶肉身的模样。}\par

% structure: p0011 source=h2 target=subsection reason=relative-heading-rank; base=h2

\subsection{第四章}

\Emph{他提出其他证词支持同一观点。}\par

若你愿知道宗徒如何出色地宣讲了这一点，请听这话如何被放在他口中，仿佛出自天主亲口，正如主说：\Quote{因为天主派遣圣子到世界上来，不是为审判世界，而是为叫世界藉着祂得救。}\BibleRef{若 3:17}看，正如你所见，主亲自肯定祂是受天主圣父派遣来拯救人类。但若你认为还需更清楚地表明，天主派遣了哪一位圣子来拯救人类——尽管天主的独生子只有一个，当天主被说成派遣祂的圣子时，显然是指派遣祂的独生子——请听达味先知极其清楚地指出那位为人类救恩而被派遣者。他说：\Quote{祂派遣了祂的圣言，治愈了他们。}你能曲解这话，把它指向肉身，仿佛可以说，天主派遣了一个纯人来治愈人类吗？你当然不能，因为达味先知和全部圣经都会向你呼喊说：\Quote{祂派遣了祂的圣言，治愈了他们。}你看，圣言被派遣来治愈人类，因为虽然治愈是通过基督赐予，但天主的圣言在基督内，并通过基督治愈万物：因此，既然基督与圣言在降生成人的奥迹中结合，基督与天主的圣言在两种性体内成为同一的天主子。当若望宗徒急于清楚说明这一点时，他说：\Quote{天主派遣了祂的圣子作世人的救主。}\BibleRef{若一 4:14}你看见他如何将天主与人结合在一个不可分割的合一中吗？因为由玛利亚所生的基督毫无疑问被称为救主，如经上说：\Quote{今天为你们诞生了一位救主，就是主基督。}\BibleRef{路 2:11}\par

% structure: p0014 source=h2 target=subsection reason=relative-heading-rank; base=h2

\subsection{第五章}

\Emph{论如何在基督内两性位格合一的德能下，圣言被正确地称为救主，或降生成人的人，以及天主子。}\par

由此可见，藉着天主的圣言与人结合的奥秘，那被派来拯救人的圣言可被称为救主，而那生于肉身的救主，藉着与圣言的结合，可被称为天主子；因此，由于天主与人结合，不加区分地使用这两种称号，凡属于天主和人的一切，都可整体地称为天主。因此，同一宗徒巧妙地补充说：\BibleQuote{谁若相信耶稣是天主子，天主就住在他内，天主的爱在他内得以成全。} \BibleRef{若一 4:12} 他告诉我们，那相信耶稣是天主子的人，\Emph{他}相信，并宣告\Emph{他}充满天主的爱。但他证明，天主的圣言就是天主子，从而要我们完全明了：天主的独生圣言和耶稣基督天主子是同一的位格。但你要更充分地被告知：尽管基督按肉身确实是从人而生的人，然而由于这奥秘的不可言喻的合一（藉此人与天主结合），基督与圣言之间没有任何分离吗？请听主的福音，更确切地说，听主亲自论及自己：\Quote{这就是永生：}他说，\Quote{使他们认识你，唯一的真天主，和你所派遣的耶稣基督。} \BibleRef{若 17:3} 你上面听到，天主的圣言被派遣来医治人类；这里你被告知，那被派遣的就是耶稣基督。如果你能的话，请把这分开——尽管你看到基督与圣言的合一如此伟大，以至于不仅是基督与圣言结合，而且因这实际的[位格]合一，基督甚至可以说就是圣言。\par

% structure: p0017 source=h2 target=subsection reason=relative-heading-rank; base=h2

\subsection{第六章.}

\Emph{在基督内只有一个\OriginalTerm{位格（自立体）}{Hypostasis}（即，位格自我）。}\par

不过你或许认为阐明这一点是小事：不是因为这事不够清楚，而是因为不信的蒙昧总使清明之事也显得晦暗。那么请听宗徒如何用寥寥数语总结了主（位格）统一的整个奥秘。他说：\BibleQuote{我们唯一的主耶稣基督}，\BibleQuote{万物都是藉着祂而有的}。\BibleRef{格前 8:6} 哦，善哉耶稣，祢的话语何等有力！因为当祢的子女论及祢时，这些话便是祢的。请看宗徒这句话的寥寥数语包含了何等丰富的内涵。他说：\BibleQuote{唯一的主}，\BibleQuote{耶稣基督，万物都是藉着祂而有的}。他是否用了任何迂回之辞来宣示这伟大奥秘的真理？或者，他是否把我们应当把握的道理大加铺陈？他说：\BibleQuote{我们唯一的主}，\BibleQuote{耶稣基督，万物都是藉着祂而有的}。他以简明扼要的言辞教导了这伟大奥秘的玄机，凭此确信，他认识到在关乎天主的事上，他的陈述无需长篇大论，因天主性使他的言语具有信服力。因为事实的展示足以证实所言之物，只要证明基于说话者的权威。所以，他说：\BibleQuote{唯一的主耶稣基督，万物都是藉着祂而有的}。请注意，你读到关于圣父圣言的同一事，也正是你读到关于基督的。因为福音告诉我们：\BibleQuote{万物是藉着祂造的；凡受造的，没有一样不是藉着祂造的}。\BibleRef{若 1:3} 宗徒说：\BibleQuote{万物都是藉着基督而有的}；福音说：\BibleQuote{万物都是藉着圣言而有的}。这些神圣的宣告互相矛盾吗？绝不。但藉着基督——宗徒说万物藉之受造的那位，以及藉着圣言——福音作者记述万物藉之受造的那位，我们应当理解为同一位。请你听一听天主圣言、祂自己就是天主，祂如何论及自己。祂说：\BibleQuote{没有人升过天，除了那从天上降下来的人子，即那在天上的人子}。\BibleRef{若 3:13} 祂又说：\BibleQuote{你们若看见人子升到祂先前所在之处}。\BibleRef{若 6:63} 祂说人子在天上；祂断言人子从天上降下。这是什么意思？你们为何喃喃自语？如果你们能，就否认吧。但你们要问所言之事的理由吗？然而我不给你们理由。天主说了这话。天主对我讲了这话：祂的圣言就是最好的理由。我抛开辩论和讨论。单是发言者的位格就足以使我相信。我不可争论所言之事的可信度，也不可探讨它。既然我不应当怀疑天主所说的是真的，为何要质疑天主所说的是否真实？祂说：\BibleQuote{没有人升过天，除了那从天上降下来的人子，即那在天上的人子}。当然，圣父的圣言永远在天上；祂如何断言人子曾在天上呢？那么你们应当明白，祂表明那永远是圣子的一位也是人子：当祂断言那不久前才显现为人子的一位永远在天上时，即指明了这一点。更指向这一点的是另一段经文，其中祂见证同一位人子，即祂所说的从天上降下的天主圣言，甚至在祂在地上说话的时候，仍在天上。因为祂说：\BibleQuote{没有人升过天，除了那从天上降下来的人子，即那在天上的人子}。请问，说话的是谁？确实是基督。但祂说话时在哪里？确实在地上。祂怎能断言祂出生时从天上降下，而说话时又在天上，或者说祂是同一位人子？当然，除了天主，没有人能从天降下，而当祂在地上说话时，除非凭藉天主的无限本性，否则绝不可能同时在天上。那么，你们终于要思考并注意：人子与天主圣言是同一位；因为祂是人子，既然祂确实由人而生；祂也是天主圣言，既然在地上说话的那一位永远居住在天上。因此，当祂真实地称自己为人子时，这是指祂的人性诞生；而祂从不离开天上，则是指祂天主性的无限属性。因此，宗徒的教导与这些神圣的话语美妙地一致：（他说：\BibleQuote{那降下的，就是那升到诸天之上，为要充满万有的}，\BibleRef{弗 4:10}）当他断言那降下的就是那升上的时。但是，除了天主圣言，无人能从天降下：祂当然\BibleQuote{具有天主的形体，却空虚自己，取了奴仆的形体，成了人的样式；形态上完全像人，祂贬抑自己，听命至死，且死在十字架上}。\BibleRef{斐 2:6} 如此，天主圣言从天降下；但人子升上。祂说同一位既升上又降下。如此你们看见人子与天主圣言是同一位。\par

% structure: p0020 source=h2 target=subsection reason=relative-heading-rank; base=h2

\subsection{第七章}

\Emph{他回到先前的主题，以对涅斯多留异端表明：那些属于天主性的事，如同出自一位天主性位格者，却被归于那人；反之，那些属于人性的事，如同出自一位人性位格者，却被归于天主，因为在基督内只有一个单一的位格自体。}\par

因此，我们遵循圣言的引导，现在可以毫无畏惧、毫不犹豫地说：人子从天降下，光荣的主被钉在十字架上；因为藉着降生成人的奥迹，天主子成了人子，光荣的主在（人子的性体）内被钉在十字架上。还有何必要多说呢？仔细详述会花费太长时间：因为我若试图考察并解释一切与此相关的事，时间将不足。因为凡愿意这样做的人，必须研读并通读整部圣经。有什么不与此相关呢？因为全部圣经都是为此而写。那么，我们必须——尽可言之——简要粗略地讲述一些事，以列举而非解释，并牺牲一些以保全其余；为此，确实最好迅速掠过某些要点，免得不得不几乎沉默地略过一切。救主在福音中说：\BibleQuote{人子来，是为寻找并拯救迷失的人。}\BibleRef{路 19:10} 宗徒也说：\BibleQuote{这话是确实的，值得完全接纳：耶稣基督来到世上，为拯救罪人，在罪人中我是首恶。}\BibleRef{弟前 1:15} 但福音作者若望也说：\BibleQuote{祂来到自己的地方，自己的人却没有接受祂。}\BibleRef{若 1:11} 你看到，圣经在一处说人子，在另一处说耶稣基督，又在另一处说天主的圣言来到世上。因此我们必须坚持，差异在于称号而非事实，在不同名称的外表下只有同一种能力（或位格）。因为虽然有时我们被告知人子来到世上，有时又是天主子来到世上，但两个名称所指的乃是同一位位格。\par

% structure: p0023 source=h2 target=subsection reason=relative-heading-rank; base=h2

\subsection{第八章}

这种称号的互换并不妨碍祂的天主的能力。\par

因为的确，当福音作者说，世界藉着祂而被造的那一位来到世上，并且祂，作为创造世界的天主，成了人子，那么使用什么特定的称号并无差别，因为无论在哪种情况下，所指的都是天主。祂的自谦和意愿并不妨碍祂的天主性，反而证明了祂的天主性，因为凡祂所愿意的，都成就了。因此，正因祂愿意，祂来到世上；正因祂愿意，祂生而为人；正因祂愿意，祂被称为人子。正如有多少词语，就有多少属于天主的能力。祂内名字的多样性并不减损祂能力的效果。无论给祂什么名字，在所有情况下所指的都是同一位位格。虽然祂的称号外表可能有些差异，但所有这些名称所指的只是一个天主位格（\OriginalTerm{尊威}{Majestas}）。\par

% structure: p0026 source=h2 target=subsection reason=relative-heading-rank; base=h2

\subsection{第九章}

祂以旧先知们的权威来证实这一陈述。\par

但既然到目前为止，我们尤其使用了福音作者和宗徒们较新的见证，现在让我们提出古先知们的证言，时而将新事与旧事相混，使众人看见圣经几乎异口同声地宣告基督将以祂自己完整的身体降生成人。因此，那位声名远扬的著名先知，既以其证言又以上主恩赐满被富足，就是在出生前即被祝圣的耶肋米亚\BibleRef{耶 1:5}，说：\Quote{这是我们的主，除祂以外，没有别的可被算作与祂相比。祂发现了知识的一切道路，并将它赐给祂的仆人雅各伯和祂的爱人以色列。此后祂在地上显现，与世人交谈。}\Quote{这是，}他说，\Quote{我们的天主。}你看先知如何以手示意天主，以手指点祂。\Quote{这是，}他说，\Quote{我们的天主。}那么告诉我，先知用这些标记和记号指出谁是天主？难道不是圣父吗？有何必要指出祂，所有人都相信他们认识的那位？因为那时犹太人并非不认识天主，他们生活在天主的法律之下。但先知显然旨在使他们认识天主子是天主。所以先知极佳地说，那发现了一切知识，即赐予法律的那位，将在地上显现，即降生成人，以便犹太人既然不怀疑赐予法律的那位是天主，他们就会认识那将降生成人的是天主，尤其因为他们听见他们相信为天主赐予法律的那位，要借着取人性居住在人间，正如祂自己借着先知应许自己的来临：\Quote{因为我自己说话的那位，看，我在这里。}\BibleRef{依 52:6}\Quote{于是，}圣经说，\Quote{在祂面前不可算有别的神。}先知美妙地预见了虚假教导，因而排除了异端的曲解。\Quote{在祂面前不可算有别的神。}因为祂是独自生为天主，出自天主：祂一声令下，宇宙就得以完成；祂的意愿是事物的起始；祂的统治是世界的构造；祂说话，万物便生成；祂命令，万物便被创造。所以唯独祂曾对圣祖们说话，居住在先知内，由圣神受孕，由童贞玛利亚诞生，出现在世界，生活在人间，将我们过犯的字据钉在十字架上，在自己内凯旋\BibleRef{哥 2:14}，借祂的死亡消灭了与我们为敌和敌对的权势；赐给众人复活的信仰，并借着祂身体的荣耀结束了人肉身的败坏。你看，所有这些只属于主耶稣基督：因此在祂面前不可算有别的神，因为唯独祂是在这光荣和独特真福中由天主所生的天主。这就是先知的教导所指向的目标：即祂被所有人认识为天主父的独生子，并且当他们听说在子面前没有别的算为神时，他们可以宣认在父和子两个位格中只有一位天主。\Quote{此后，}他说，\Quote{祂在地上显现，与世人交谈。}你看这多么清楚地指向主的来临和诞生。因为圣父——我们读到祂只能在子内被看见——并没有在地上显现，也没有降生成人，也没有与世人交谈？绝对没有。那么你看，这一切都是指着天主子说的。因为先知说天主将在地上显现，而除了子以外没有别位在地上显现，显然先知唯独指着祂说的，后来事实证明了这一点。因为当祂说天主将被看见时，除非指着后来确实被看见的那位，否则这话不能真实。就此打住。现在我们转向另一点。\Quote{埃及的劳苦，}依撒意亚先知说，\Quote{与雇士和色巴的商人，高大的男子，要归于你，成为你的仆人。他们要带着锁链跟随你，要朝拜你，向你恳求：因为在你内有天主，除你以外没有别的神。因为你是我们的天主，我们素不认识你，以色列的救主天主。}\BibleRef{依 14:14}圣经总是何等奇妙地一致！因为前面提及的先知说：\Quote{这是我们的天主}，而这位说：\Quote{你是我们的天主。}在一位中有天主性的教导，在另一位中有人的宣认。一位显出师傅教导的特性，另一位显出百姓宣认的特性。因为现在请思量：耶肋米亚先知在教会中每日教导，如上所述，论及主耶稣基督说：\Quote{这是我们的天主}，整个教会除了以另一位先知对主耶稣所说的\Quote{你是我们的天主}来回应，还能做什么呢？所以，过去无知的提及能与现在的承认很好地结合，以百姓的话说：\Quote{你是我们的天主，我们素不认识你。}因为从前沉迷于魔鬼的迷信而不认识天主的人，现在归依了信仰，能说：\Quote{你是我们的天主，我们素不认识你。}\par

% structure: p0029 source=h2 target=subsection reason=relative-heading-rank; base=h2

\subsection{第十章。}

\Emph{他由犹太化犹太人的亵渎以及归依基督信仰者的宣认，证明基督的天主性。}\par

但是，如果你更愿意从犹太人的代表那里得到证明，那么请考虑犹太民族，在他们可悲的无知和邪恶的迫害之后，他们悔改了，并到处承认天主，看看他们能否正当地说：\Quote{你是我们的天主，我们曾不认识你。}但我要添加另一件事，以向你证明这不单从那些承认祂的犹太人，也从那些否认祂的犹太人。因为请问那些仍处于不信状态的犹太人，他们是否认识相信天主。他们必然会承认他们既认识又相信祂。但另一方面，请问他们是否相信天主子。他们会立刻否认并开始亵渎祂。那么你看到先知论及此人所言，正是犹太人一直无知、至今仍不认识的那位，而非他们自以为相信和承认的那位。因此，那些曾在无知中从犹太教出来而进入信德的人，可以完全正当地说：\Quote{你是我们的天主，我们曾不认识你。}因为那些在无知之后相信的人，正当地说他们不认识那位他们至今仍未相信、并且努力不去认识的那位。因为很明显，那些在先前无知之后承认祂的人，说他们先前不认识祂，而他们至今仍在无知中否认的那位。\par

% structure: p0032 source=h2 target=subsection reason=relative-heading-rank; base=h2

\subsection{第十一章}

\Emph{他回到依撒意亚的预言。}\par

他说：\Quote{埃及的劳力和雇士的货物，以及身材高大的色巴人，都要来归于你。}无人怀疑，在这些不同民族的名称中，预示着将信仰的外邦人的来临。但你不能否认，外邦人已经归于基督，因为自从基督教的名号兴起以来，他们已归于主耶稣基督，不仅是在信仰上，实际上也在名称上。因为他们被称为他们实际上所是的，信德的工作成了他们被称呼的标志。他说：\Quote{他们要来归于你，成为你的；他们要带着锁链跟随你。}正如有强制之链，同样也有爱之链，如主所说：\Quote{我用爱的锁链牵引他们。}\BibleRef{欧 11:4}的确，这些锁链是伟大的，是无可言喻的爱之链，那些被它们捆绑的人以他们的镣铐为乐。你想知道这是真的吗？请听宗徒保禄如何因他的锁链而欢欣鼓舞，他说：\Quote{所以，我在这主里的囚犯劝勉你们。}\BibleRef{弗 4:1}又说：\Quote{我劝勉你们，我这个为基督耶稣而成为囚犯的保禄老人。}你看到他如何以他的锁链的尊荣为乐，他的榜样甚至激励了他人。但毫无疑问，哪里有对主的专一之爱，哪里就有为主而戴的锁链的专一之乐，正如经上记载：\Quote{众信徒都是一心一意。}\BibleRef{宗 4:32}他说：\Quote{他们要朝拜你，向你祈求：因为天主在你内，除你以外没有天主。}宗徒清楚地解释了先知的话，他说：\Quote{天主在基督内使世界与自己和好。}\BibleRef{格后 5:19}他说：\Quote{在你内，有天主，除你以外没有天主。}当先知说\Quote{天主在你内}时，他绝妙地指出了那可见的一位，更指出了那在可见者内的一位，通过在指出二性的同时区分居住者与祂所居住的那位，而不是否认（位格的）合一。\par

% structure: p0035 source=h2 target=subsection reason=relative-heading-rank; base=h2

\subsection{第十二章}

\Emph{救主这名号如何以一层意义归于基督，以另一层意义归于人。}\par

他说：\Quote{你是我们的天主，我们曾不认识你，以色列的救主天主。}尽管圣经已用许多清晰的标记表明了这里所说的是谁，但它最清楚地指向了基督的名字，使用了救主这个名号：因为救主就是基督，正如天使所说：\Quote{今天为你们诞生了一位救主，就是主基督。}\BibleRef{路 2:11}因为众所皆知，在希伯来语中，\Quote{耶稣}意为\Quote{救主}，正如天使向童贞玛利亚所宣布的：\Quote{你要给祂起名叫耶稣，因为祂要把自己的民族从他们的罪恶中拯救出来。}\BibleRef{玛 1:21}并且，为了不让你以为祂被称为救主的意义与别人相同（\Quote{上主给他们兴起了一位拯救者，即克纳兹的儿子敖特尼耳}\BibleRef{民 3:9}，又说：\Quote{上主给他们兴起了一位拯救者，即革辣的儿子厄胡得}），祂补充说：\Quote{因为祂要把自己的民族从他们的罪恶中拯救出来。}但是，没有人有能力将他的民族从罪的囚禁中赎回来——只有祂才能做到，经上论及祂说：\Quote{请看天主的羔羊，除免世罪者。}\BibleRef{若 1:29}因为其他的拯救者拯救的不是他们自己的民族，而是天主的民族，并且不是从他们的罪恶中，而是从他们的仇敌中。\par

% structure: p0038 source=h2 target=subsection reason=relative-heading-rank; base=h2

\subsection{第十三章}

\Emph{他解释依撒意亚先知以谁的名义说：\Quote{你是我们的天主，我们曾不认识你。}}\par

\Quote{那么，你是我们的天主，}他说，\Quote{我们却不认识你，啊，拯救者以色列的天主。}你认为说这话的主要是什么人？这些词特别适合谁的口，是犹太人还是外邦人？若你说是犹太人：确实，犹太人不认识基督，正如经上所记：\Quote{以色列却不认识我，我的百姓也不留意。}\BibleRef{依 1:3} 以及：\Quote{世界是藉着祂造的，世界却不认识祂。祂来到了自己的地方，自己的人却没有接受祂。}\BibleRef{若 1:11} 但若你说是外邦人，显然外邦世界曾委身于偶像，不认识基督，尽管他们也不认识圣父；但如今他们若认识祂，也只是藉着基督。那么你看，无论相信的子民是属于犹太人或外邦人，两者都能真实地为自已说：\Quote{你是我们的天主；我们却不认识你，拯救者以色列的天主。}因为从前敬拜偶像的外邦人不认识天主；而否认主的犹太人不认识天主子。因此两者都真实地论基督说：\Quote{你是我们的天主，我们却不认识你。}因为那些不信天主的人对祂的无知，与那些否认天主子的人是一样的。所以，如果基督是当被相信的，正如真理所宣称、神性所断言、而且基督亲自——祂兼是二者——所宣称的，你们为何在癫狂中悲惨地试图将天主与基督隔开？为何寻求将祂的身体与天主子分离，并试图将天主与祂自己分开？你们正在割裂那合一的，分离那相连的。要相信天主关于天主的圣言：因为你们不可能做出比用口承认天主关于祂自己所教导的更好的对天主神性的宣认。因为你们必须知道，正如先知所说：\Quote{上主本身就是天主，祂发现了全部知识之路；祂曾在地上显现，与人同行。}祂将信德之光带入世界。祂指明了救恩之光。\Quote{因为上主是天主，祂赐给了我们光明。}所以，要相信祂、爱祂、宣认祂。因为正如经上所记：\Quote{一切膝都要向祂跪拜，天上的、地上的和地下的，一切唇舌都要宣认耶稣基督是主，在天主圣父的光荣中。}\BibleRef{斐 2:10} 无论你们愿意与否，你们都不能否认耶稣基督是天主圣父光荣中的主。因为完美宣认的最高美德，就是承认耶稣基督在天主圣父的光荣中永远是主和天主。\par

\SourceRecord{Translated by C.S. Gibson.；Nicene and Post-Nicene Fathers, Second Series；Vol. 11.；Edited by Philip Schaff and Henry Wace.；Buffalo, NY: Christian Literature Publishing Co.,；1894.；<http://www.newadvent.org/fathers/35094.htm>.}

% structure: p0001 source=h1 target=section reason=page-title

\section{论降生成人（第五卷）}

% structure: p0002 source=h2 target=subsection reason=relative-heading-rank; base=h2

\subsection{第一章}

\Emph{他严厉谴责白拉奇异端者的错谬，他们宣称基督只是一个纯然的人。}\par

我们在第一卷中说过，那仿效并追随白拉奇主义（Pelagianism）的异端，竭尽全力并千方百计使人相信，主耶稣基督、天主子，在由童贞玛利亚所生时，仅仅是一个纯然的人；并且后来他踏上德性之路，凭他的圣善虔敬的生活，堪受天主性本身与他结合，由于他生活的这种圣洁，而配得上这份尊荣；如此，他们完全剥夺了他神圣根源的荣耀，只留给他基于功绩而被选的地位。他们的目的与努力在于：将他降至普通人的水平，使其成为一般民众之一，从而断言：所有人都能凭其良善的生活和行为，获得他凭其良善生活所获得的一切。这确实是一种极其危险和致命的断言，它夺走真正属于天主的，却对人作出虚假的承诺；它应在两点上受谴责，因其以邪恶的亵渎攻击天主，并给人以虚假保证的希望。这是一种极其悖逆和邪恶的断言，因为它将不属于人的东西给予人，并夺走属于天主的东西。因此，对于这种危险而致命的邪恶，这个新近兴起的异端，可谓在拨动并重燃其余烬，从古老的灰烬中升起新的火焰，它断言我们的主耶稣基督生为一个纯然的人。所以，我们何须询问其后果是否危险？因为在其根源上，它已完全错误。既然在其起始处已无审查的余地，便无需考察其结局如何。因为，若它（此乃最可怕的罪）同样夺走属于天主的东西，那么询问它是否像之前的异端一样对人作出相同的承诺，又有何意义？因此，当我们看到其开端时，几乎不必追问接下来会怎样，仿佛在否认天主的人身上，还能证明他并非不敬神。\newline{}这个新异端，正如我们已多次声明的，说主耶稣基督由童贞玛利亚所生，仅仅是一个纯然的人；因此玛利亚应被称为 \OriginalTerm{基督之母}{Christotocos}，而非 \OriginalTerm{天主之母}{Theotocos}，因为她只是基督的母亲，而非天主的母亲。此外，在这亵渎的声明上，他们还增加了同样邪恶愚蠢的论证，说：\Quote{没有人能生出比她更早存在的人。}仿佛天主独生子的降生——由先知预言，自创世以来便宣告——可以用人类的理性来对待或衡量。或者，啊，异端者，无论你是谁，诽谤她生育的童贞玛利亚，她凭一己之力促成并完成了所发生的事，以至于在如此重大的事和事件中，可以将人性的弱点拿来作为反对的理由？因此，如果在这伟大的事件中有任何属于人的作为，那就寻找人的论证。但如果所发生的一切皆出自天主的德能，为何当你看到这是天主大能的作为时，却要思考在人所不能的事呢？此事容后再谈。现在让我们继续我们之前稍早开始讨论的主题；好让每个人都知道，你们试图在白拉奇主义的灰烬中煽风点火，并通过吹出新亵渎的口气来复活余烬。\par

% structure: p0005 source=h2 target=subsection reason=relative-heading-rank; base=h2

\subsection{第二章}

\Emph{论聂斯多略的教义与白拉奇异端者的错谬密切相关。}\par

那么，你说基督生为一个纯然的人。但这一点，正如我们在第一卷中清楚表明的，白拉奇的邪恶异端确实主张过：基督生为一个纯然的人。此外，你还加上：万有的主耶稣基督应被称为一种接受了天主的形态（\OriginalTerm{接受了天主者}{\GreekText{Θεοδ}\GreekText{οχος}}），即不是天主，而是天主的接受者；因此你的观点是：他受尊荣，不是因为他本身是天主，而是因为他将天主接纳到自己内。但显然，我之前所说的那个异端也主张过这一点：基督不应因自身是天主而被敬拜，而是因为他凭其良善虔敬的行为赢得了这一点：即天主住在他内。你看，你们正在喷吐白拉奇主义的毒液，嘶嘶作响着白拉奇主义的精神。因此，你们似乎早已被审判，而非现在才要受审判；因为你们的错误是同一的，你们应当相信落在同一的谴责之下。目前暂不提及：你们将主比作皇帝的雕像，并爆发出如此邪恶亵渎的不虔敬，以至于在这癫狂中，你们似乎甚至超过了白拉奇本人，而他在不虔敬上几乎超过了所有其他人。\par

% structure: p0008 source=h2 target=subsection reason=relative-heading-rank; base=h2

\subsection{第三章}

\Emph{白拉奇异端者和聂斯多略派归给基督的这种对天主性的分享，如何为所有圣人所共有。}\par

你于是说基督应被称为一个\OriginalTerm{收纳天主者}{\GreekText{Θεοδόχος}}，即祂之所以受尊敬不是因祂本身是天主，而是因祂收纳了天主内在于祂。这样你就认为祂与所有其他圣人没有区别：因为所有圣人确实有天主内在于他们。因为我们深知天主在圣祖们内，且在先知们内发言。总之我们相信，我且不说宗徒和殉道者，但所有圣人和天主的仆人都有天主之神内在于他们，依此：\Quote{你们是永生天主的殿：正如天主所说，我要住在他们中间。}\BibleRef{格后 6:16}又说：\Quote{你们不知道你们是天主的殿，天主之神住在你们内吗？}\BibleRef{格前 3:16}因此我们都是\OriginalTerm{收纳天主的人}{\GreekText{Θεοδόχοι}}；这样你就说所有圣人仅仅与基督一样，且与天主同等。但离弃如此邪恶可憎的异端吧，竟使造物主与被造物相比，主与祂的仆人相比，天地天主与地上脆弱者相比：且因祂的恩惠而施此不义于祂；即，祂因居于人内而尊荣人，因此被说成只是与人相同。\par

% structure: p0011 source=h2 target=subsection reason=relative-heading-rank; base=h2

\subsection{第四章}

\Emph{基督与圣人之间的区别。}\par

此外，在祂与所有圣人之间存在同一区别，正如住所与居住者之间的区别；因为居住者而非住所的作为，若住所被居住，则建造房屋和占有房屋都取决于他。我说，他可以选择，若他愿意，使之成为住所，并且当他使之成为住所后，就住在其中。\Quote{你们是在寻求一个证据吗？}宗徒说，\Quote{基督在我内说话的证明？}\BibleRef{格后 13:3}又在别处说：\Quote{你们岂不是该知道耶稣基督在你们内，除非你们是经不起考验的？}又说：\Quote{在心内，使基督因信德住在你们心里。}\BibleRef{弗 3:16}你们难道看不出宗徒的教导与你们的亵渎之间有何区别吗？你们说天主住在基督内如同在人内。祂作证说基督亲自住在人内：这当然，如你们所承认，血肉不能做到；所以祂被证明是天主，正是从你们否认祂是天主的事实。因为既然你们不能否认那住在人内的是天主，那么我们必须相信，那我们知道住在人内的，实在就是天主。那么，无论是圣祖、先知、宗徒、殉道者、圣人，他们每人都内有天主，都成为天主的儿子，都是\OriginalTerm{收纳天主的人}{\GreekText{Θεοδόχοι}}，但以非常不同和独特的方式。因为凡信天主的人，都是藉着义子而成为天主的儿子；但只有独生者是本性上的儿子：祂由祂的父所生，不是出于任何物质实体，因为万有和万有的实体都藉着天主的独生子而存在——也不是出于无，因为祂来自父：不是像出生，因为天主内没有虚无或可变之处，而是以不可言喻和不可理解的方式，天主圣父，在祂自己重生之处，生了祂的独生子；如此，从至高的、非受生的、永恒的圣父，产生了至高的、独生的、永恒的圣子。祂在肉身中与在神性中必须是同一位格：在身体中与在荣耀中必须是同一位格，因为当祂即将在肉身中诞生时，祂内在没有分裂或分开，仿佛祂的一部分生了而另一部分未生；或仿佛神性的一部分后来临于祂，而这部分在祂由童贞女诞生时并不在祂内。因为按照宗徒：\Quote{天主圆满的全部神性在基督内以身体居住。}\BibleRef{哥 2:9}不是有时住在祂内，有时不住；也不是后来才有而先前没有：否则我们就陷入佩拉吉乌斯那个邪恶的异端，以致说从某一固定时刻天主住在基督内，然后临于祂；当祂藉其生活和行为赢得这一点时，即神性的能力应住在祂内。这些事属于人，我说，属于人，不属于天主——就像人按人的软弱所能做的，他们应谦卑自己于天主，服从天主，使自己成为天主的居所，并藉其信德和虔诚赢得这一点，使天主作为他们的客人和内住者。因为一个人越配得天主的恩赐，天主的恩宠就越酬报他：一个人越显得相称于祂，就越配得天主，就越享受天主的临在，按照主的应许：\Quote{谁爱我，必遵守我的话；我父也必爱他，我们要到他那里去，并要在他那里作我们的住所。}\BibleRef{若 14:23}但基督的情况完全不同；在祂内，全部神性以身体居住：因为祂内具有全部神性，以致祂将祂的全部丰富赐予众人，而祂——因全部神性住在祂内——祂自己按祂认为各人配得祂临在的程度，住在每位圣人内，并将祂的丰富赐予众人，然而祂自己仍保持全部丰富——即使当祂在肉身中在地上时，祂仍临在于所有圣人的心中，并以祂无限的能力和威严充满天、地、海，乃至整个宇宙；然而祂在祂自己内如此完全，以致整个世界不能容纳祂。\par

% structure: p0014 source=h2 target=subsection reason=relative-heading-rank; base=h2

\subsection{第五章}

\Emph{基督在时间内的诞生之前，先知们常称祂为天主。}\par

祂就是先知所说的那一位：\BibleQuote{因为天主在祢内，除祂以外没有别的神。祢是我们的天主，我们却不认识祢，以色列的天主、救主。}\BibleRef{依 45:14}祂\BibleQuote{以后在地上出现，与世人往来。}\BibleRef{巴 3:37}先知达味也论及祂，并以祂的口气说：\BibleQuote{祢是从我母胎中就作我的天主。}这清楚表明，那既是主又是人的一位从未与天主分离；在童贞母的胎中，天主性的圆满已住在祂内。正如同一先知别处所说：\BibleQuote{真理从地上生出，正义由天俯视。}为叫我们知道，当天主子从天俯视（即降临）时，正义由童贞女的肉身而生，并非身体的幻影，而是真理——因为祂就是真理，正如祂亲自为真理作证：\BibleQuote{我是真理，也是生命。}\BibleRef{若 14:6}因此，正如我们在前几卷中已证明，这真理，即主耶稣基督，在由童贞女所生时就是天主，现在让我们按照本卷之前的计划，证明那将由童贞女所生的一位，早已被宣告为天主。于是先知依撒意亚说：\BibleQuote{你们不要依靠鼻孔内有气息的人，因为他被视为崇高。}或如希伯来原文更精确清晰：\BibleQuote{因为他被看为崇高。}但藉着\BibleQuote{不要依靠}一词——这词劝阻暴力——他卓越地指出了迫害的扰乱。\BibleQuote{不要依靠，}他说，\BibleQuote{鼻孔内有气息的人，因为他被视为崇高。}他岂不是在同一句话中既论及人性的承受，也论及天主性的真实吗？\BibleQuote{不要依靠，}他说，\BibleQuote{鼻孔内有气息的人，因为他被视为崇高。}我请问你，他岂不是明明白白地对迫害主的人说：\BibleQuote{你们不要依靠那受迫害的人}，因为这人是天主；虽然祂显现于卑微的人身中，却仍居于天主荣耀的高位？但藉着\BibleQuote{不要依靠鼻孔内有气息的人}，他以人身最明显的标记卓越地显示了祂的人性，并且无所畏惧、充满信心，仿佛他要同样急切地主张祂人性的真实和天主性的真实——因为这才是真实而大公的信德：相信主耶稣基督拥有真实身体的实体，正如祂拥有真实而完美的天主性。除非你以为，他用\BibleQuote{崇高}一词代替\BibleQuote{天主}这一点能说明什么；然而圣经的习惯是用\BibleQuote{崇高}指\BibleQuote{天主}，如先知说：\BibleQuote{至高者发出声音，大地震动}；\BibleQuote{惟独祢是至高的，超越全地。}依撒意亚也说：\BibleQuote{那崇高而尊贵、永远居住者}：\BibleRef{依 57:15}在这里我们清楚理解，正如他此处没有加上\BibleQuote{天主}的名称而用了\BibleQuote{至高者}，同样地，他也以\BibleQuote{至高者}的名号来指天主。所以，既然先知所说的天主圣言已清楚预告主耶稣基督既是天主又是人，现在让我们看看新约是否与旧约的见证相符合并和谐一致。\par

% structure: p0017 source=h2 target=subsection reason=relative-heading-rank; base=h2

\subsection{第六章}

\Emph{祂用新约的经文阐明同一道理。}\par

\Quote{那，}宗徒\Person{若望}说，\BibleQuote{那从起初就有的，是我们所听见的，是我们亲眼所看见的，是我们所注视过，我们的手所摸过的，论的是生命之圣言：因为生命已显现出来，我们看见了，也作证，且向你们宣报那与圣父同在、且已显示给我们的永生。}\BibleRef{若一 1:1} 你看，旧的见证如何被新的见证所证实，新的宣讲如何支持了古代的预言。\Person{依撒意亚}说：\Quote{你们要停止信任那鼻孔里有气息的人，因为他被视作崇高。}但\Person{若望}说：\BibleQuote{那从起初就有的，是我们亲眼所看见的，是我们所注视过，我们的手所摸过的。}前者说，祂作为人将受犹太人的迫害；后者宣告，祂作为人被人手所触摸。一位预言，他所宣告为人者将是至高者天主；另一位断言，他所显示曾被人手触摸者，从起初就是天主。因此，他们两人极其清楚地将主耶稣基督指明为既是天主又是人；同一位，祂永远是神，后来成为人，因此祂既是神又是人，因为天主自己成了人。然后他说：\BibleQuote{那从起初就有的，是我们所听见的，是我们亲眼所看见的，是我们所注视过，我们的手所摸过的，论的是生命之圣言；生命已显现出来，我们看见了，也作证，且向你们宣报那与圣父同在、且已显示给我们的永生。}你看，那么多的证明和方式，非常不同且众多，那位深受爱戴、如此忠于天主的宗徒，借着它们指示了神圣降生成人的奥迹。首先，他证明那从起初就永远存在者，曾在肉身中被看见。惟恐对不信者而言，他说祂被看见和听见或许显得不足，他进而说祂被摸过，即被他和别人的手触摸和感觉到。的确，他通过显示祂如何取了肉身，令人钦佩地排除了\OriginalTerm{马西翁派人}{Marcionites}的观点和\OriginalTerm{摩尼派人}{Manichees}的错误，使无人以为一个\OriginalTerm{幻影}{phantom}出现在人前，因为一位宗徒已宣告一个\OriginalTerm{真实的身体}{true body}被他触摸过。然后他加上：\BibleQuote{生命之圣言：生命已显现出来}；并说他看见了、宣布了、宣扬了：这样同时尽着信德的职责，使不信者感到惊惧，以便当他宣布他宣讲祂时，使那不愿听从的人认识到他所处的危险。他说：\BibleQuote{我们向你们宣报那与圣父同在、且已显示给我们的永生。}他教导说，那永远与圣父同在的，显现于人；那从起初就永远存在的，被人看见；那无始的生命之圣言，被人手触摸。你看他展开肉身与天主结合的奥迹的方式之多、变化之多、详尽与清晰，以至于无人能在承认两者之一时不承认两者。正如宗徒自己在别处清楚地说：\BibleQuote{因为耶稣基督昨天、今天，直到永远，常是一样的。}\BibleRef{希 13:8}这就是他在上文所说的：\BibleQuote{那从起初就有的，我们的手曾摸过。}这不是说一个神灵按本性可以被触摸，而是说成为血肉的圣言，在某种程度上，在与祂结合的人性中被触摸了。所以耶稣\BibleQuote{昨天、今天是一样的}：即同一位在创世之前、也同一位在肉身中；同一位在过去，也同样在现在，同一位也直到永远，因为祂在万世之前、也在万世之中都是一样的。所有这一切就是主耶稣基督。\par

% structure: p0020 source=h2 target=subsection reason=relative-heading-rank; base=h2

\subsection{第七章}

\Emph{祂再次从基督内两性在一位中的结合表明，那属于天主性的可以正确地归于人，而那属于人性的可以归于天主。}\par

在世界的起源之前的那同一位格，怎会是在最近才诞生的呢？因为最近在人性中诞生的同一位格，就是在万有兴起之前已为天主的祂。因此，基督的名包含了天主的名所包含的一切；因为基督与天主的结合如此紧密，以至无人能在使用基督之名时不藉基督之名提及天主，也无人能在提及天主时不藉天主之名提及基督。并且，由于藉着祂圣洁诞生的荣耀，每一本性的奥迹在祂内结合在一起，凡存在的——我指既有人性的也有神性的——一切都被视为天主。因此，宗徒保禄用信德被揭开的眼睛看见在基督内那不可言喻的荣耀的全部奥迹，便这样说话，邀请那些不知天主美善的万民向天主谢恩：\BibleQuote{感谢父，祂使我们有资格在光明中分享圣徒的福分，祂救了我们脱离黑暗的权势，将我们迁入了祂爱子的国度，在祂内我们藉着祂的血获得救赎，罪过的赦免；祂是不可见的天主的肖像，一切受造物的首生者：因为在天上和地上的一切，看得见的和看不见的，无论是宝座、宰制者还是掌权者，都是在祂内受造的：万有都是藉着祂并为了祂而受造。祂在万有之先，万有都靠祂而存在。祂是身体——教会的头；祂是元始，从死者中首先复活的，好使祂在一切事上居首位。因为父乐意使一切的圆满居住在祂内，并藉着祂使万有与自己和好，藉着祂十字架的血成就和平，无论是地上的还是天上的。} \BibleRef{哥 1:12} 的确，这不需要任何进一步解释的帮助，因为它表达得如此充分和清晰，以致其本身不仅包含信德的实质，而且是对信德的清晰阐述。因为他命令我们感谢圣父：并加上了一个如此感恩的有力理由；即，因为祂使我们有资格与圣徒一同分享，并救我们脱离黑暗的权势，将我们迁入祂爱子的国度，在祂内我们获得了救赎和罪过的赦免：祂是不可见的天主的肖像，一切受造物的首生者；因为在祂内并藉着祂万有受造；祂既是万有的创造者又是统治者。在此之后又是什么呢？\Quote{祂是}，他说，\Quote{身体——教会的头；祂是元始，从死者中首先复活的。} 圣经把复活说成是一种诞生：因为正如诞生是生命的开始，同样，复活生出生命。因此，复活实际上也被称为重生，正如主的话：\BibleQuote{我实在告诉你们：你们这些跟随我的人，在重生的时候，当人子坐在祂荣耀的宝座上时，你们也要坐在十二个宝座上，审判以色列十二个支派。} \BibleRef{玛 19:28} 因此，他称祂为从死者中首先复活的，而他先前曾宣告祂是不可见的圣子和天主的肖像。但谁是不可见的天主的肖像呢？岂不是独生子、天主的圣言吗？我们怎能说那被称为不可见天主的肖像和圣言的祂从死者中复活了呢？随后又是什么呢？\Quote{好使祂在一切事上居首位：因为父乐意使一切的圆满居住在祂内，并藉着祂使万有与自己和好，藉着祂十字架的血成就和平，无论是地上的还是天上的。} 当然，万有的创造者不需要在万有中居首位？那创造万有的祂，不需要在祂所造的万有中居首位？我们怎能说，关于圣言，当祂自己在万有的起源之前就是天主的独生子和圣言，并在祂内怀着不可见的圣父，因此首先在祂内有一切的圆满，以便祂自己成为万有的圆满时，却使天主乐意让一切的圆满居住在祂这位从死者中首先复活的那位里面？接着呢？\Quote{藉着祂十字架的血使万有和平，无论是地上的还是天上的。} 当然，当他称祂为从死者中首先复活的那位时，他已尽可能清楚地表明他在谈论谁。因为万有岂是藉着圣言或圣神的血和好并带来和平的？绝对不是。因为不能受苦的本性不可能遭遇任何形式的苦难，而且除了人之外，没有任何人的血可以流出，也没有任何人能死；然而，在接下来的经文中被说成死了的同一位格，在上面却被称作不可见的天主的肖像。这怎么可能呢？因为宗徒们竭尽一切可能防止有人认为基督内有任何分裂，或者天主子与人子结合后，可能通过荒谬的诠释被弄成两个位格，从而那在自身内只是一个的祂，可能因我们错误和邪恶的观念而被弄成在一个本性中有两个位格。因此，宗徒的宣讲最卓越、最令人钦佩地从天主的独生子过渡到与人子结合的天主子，以便教义的阐释符合所发生之事的实际进程。因此，他以不间断的联系继续，并建造了一座桥梁，使得你可以没有任何空隙或分离地在时间的终点找到我们在世界开始时读到的那位；并且，你不可因承认某种分裂和错误的分离而想象天主子在肉体内是一个位格，在圣神内是另一个位格；因为宗徒的教导已通过祂在身体内诞生的奥迹将天主和人如此紧密地联系在一起，以致表明那就是同一位格在十字架上使万有与自己和好，祂在世界奠基之前已被宣告为不可见的天主的肖像。\par

% structure: p0023 source=h2 target=subsection reason=relative-heading-rank; base=h2

\subsection{第八章}

\Emph{祂藉主的权威确认了宗徒的判断。}\par

虽然这是宗徒的话，但却是主自己的道理。因为同一位借着祂的宗徒对基督徒说这话，祂自己曾在福音中向犹太人说过非常相似的话，祂说：\BibleQuote{但现在你们想要杀我，一个将我从天主那里听到的真理告诉你们的人；因为我不是凭自己来的，而是祂派遣了我。}\BibleRef{若 8:40}祂清楚地表明祂既是天主又是人：是人，因祂说祂是一个人；是天主，因祂肯定自己被派遣。因为祂必须与祂所来自的那一位同在；祂来自那一位，来自祂所说的那派遣者。因此，当犹太人对祂说：\BibleQuote{你还没有五十岁，就见过亚伯拉罕吗？}祂用完全适合祂永恒与光荣的话回答说：\BibleQuote{我实实在在告诉你们：在亚伯拉罕存在以前，我是。}那么我问，你认为这是谁的话？当然是基督的，毫无怀疑。但这位刚刚出生不久的人，怎能说祂在亚伯拉罕之前呢？只因着天主圣言，祂与圣言完全合一，好使众人理解基督与天主结合的紧密：因为凡天主在基督内说的，天主性的合一都完全归为自己。但祂意识到自己的永恒，因此当祂在身体中时，祂用从前在圣神中对梅瑟说过的话回答犹太人。因为这里祂说：\BibleQuote{在亚伯拉罕存在以前，我是。}但对梅瑟祂说：\BibleQuote{我是自有者。}\BibleRef{出 3:14}祂确实以语言的奇妙庄严宣告了祂天主性本质的永恒，因为没有什么比说祂永远是更配得上天主的。因为\BibleQuote{是}不允许过去有开始或未来有终结。所以这非常清楚地论及永恒天主的本性，因为它精确描述了祂的永恒。主耶稣基督自己在谈论亚伯拉罕时，也通过所用词语的差异表明了这一点，祂说：\BibleQuote{在亚伯拉罕存在以前，我是。}祂谈到亚伯拉罕时说：\BibleQuote{在他成为存在以前，}谈到自己时说：\BibleQuote{我是}，因为成为存在属于暂时的事物，\Emph{是}属于永恒。因此祂将\BibleQuote{成为存在}归给人类的短暂性，而将\BibleQuote{是}归给自己的本性。这一切都在基督内，借着人性与天主性结合的奥迹，祂曾是那永远\Emph{是}者，因此能说祂早已\Emph{是}。\par

% structure: p0026 source=h2 target=subsection reason=relative-heading-rank; base=h2

\subsection{第九章}

\Emph{既然那些从梅瑟时代以来向以色列子民展示的奇妙作为都归于基督，所以祂一定在时间中的诞生之前已经存在很久了。}\par

宗徒愿意将这道理清楚明白地公诸于众时，就这样说：\Quote{耶稣既从埃及地救出百姓，后来就把那些不信的灭绝了。}但在别处我们也读到：\Quote{我们也不要试探基督，像他们中有些人试探过而被蛇所灭。}\BibleRef{格前10:9}宗徒之长伯多禄也说：\Quote{现在你们为什么试探天主，要把一个轭放在门徒的颈上，这轭我们的祖先和我们都不能负？但是，我们相信，我们得救，是藉着主耶稣基督的恩宠，正和他们一样。}\BibleRef{宗15:10}我们确知，天主的百姓从埃及被救出，乾步行过浩大的水域，并在广漠的旷野中被保存，皆出自唯独天主；正如经上所载：\Quote{上主独自领导了他们，没有外来的神与他们同在。}\BibleRef{申32:12}宗徒如何能在如此众多而清楚的经文中宣告，犹太百姓是由\Person{耶稣}从埃及救出，并且\Person{基督}当时在旷野中被犹太人所试探，说：\Quote{我们也不要试探基督，像他们中有些人试探过而被蛇所灭}？此外，蒙福的宗徒\Person{伯多禄}论及所有活在旧约律法下的圣人们，说他们是由我们主\Person{耶稣基督}的恩宠得救。那么，如果你能，就脱身吧，无论你是谁——你这以虚浮的口舌和亵渎的精神发怒的人，以为\Person{亚当}和\Person{基督}之间根本没有区别；你否认祂在生于童贞女之前就是天主，你要清楚证明，在祂的身体出现之前，祂不是天主。因为看，宗徒说百姓是由\Person{耶稣}从埃及地救出；基督在旷野中被不信者试探；我们的祖先，即圣祖和先知，是由我们主\Person{耶稣基督}的恩宠得救。你若能否认，就否认吧。你若能否认我们所读到的一切，我也不会惊讶，因为你已经否认了我们所信的一切。要知道，那时在\Person{基督}内就是天主带领百姓出埃及，在\Person{基督}内就是天主被那试探的百姓所试探，在\Person{基督}内就是天主藉着祂丰厚的恩宠拯救了所有义人：因为（降生成人）奥迹的一体性，天主和基督这两个术语彼此贯通，以致天主所做的，我们可说基督做了；后来基督所受的，我们可说天主受了。所以，当先知说：\Quote{你中间不可有别的神，也不可敬拜别的神}时，他是以同一意义和同一精神宣布的，正如宗徒说基督是带领以色列百姓出埃及的那位；为表明那由童贞女诞生为人的那位，甚至藉奥迹的一体性仍在天主内。否则，除非我们相信这事，我们要么与异端者一同相信基督不是天主，要么违反先知的教导，认为祂是一位新的神。但愿天主的公教百姓远避这样的想法：似乎与先知相异或与异端者同流；以免那本当蒙福的百姓陷入咒诅，并被指控将希望寄托于人。因为凡宣称主\Person{耶稣基督}在诞生时只是一个凡人，无论他信或不信，都双重地处于咒诅之下。因为若他信，经上说：\Quote{凡信赖世人的人，是可咒骂的。}\BibleRef{耶17:5}但若他不信，他仍然被咒诅，因为虽然他不信人，却全然否认了天主。\par

% structure: p0029 source=h2 target=subsection reason=relative-heading-rank; base=h2

\subsection{第十章}

\Emph{他解释宣认耶稣和瓦解耶稣各是什么意思。}\par

原来，那为天主所爱的若望，从主亲自启示他所预见的，并因此论及那在他内发言的一位说：\Quote{凡神}，他说，\Quote{凡宣认耶稣在肉身内降来的，便是出于天主；凡分裂耶稣的，便不是出于天主：这便是指那假基督的神，你们已听说他要求，如今他已是在世界上了。}啊，天主何等奇妙而独特的良善！祂如同最审慎、最熟练的医生，预先预告了将临到祂教会身上的疾病，并在预示灾祸时，同时提供了救治之法：使所有人一见到邪恶临近，便立即尽可能逃离那已知的迫在眉睫之事。因此圣若望说：\BibleQuote{凡分裂耶稣的神，不是出于天主；这便是假基督的神。}啊，你这异端者，你认出他了吗？你认出这话明明就是指着你说的吗？因为除了那不宣认祂是天主的人之外，无人这样分裂耶稣。既然教会的全部信德和全部敬拜都在于此，即宣认耶稣是真天主，那么还有谁比那否认祂内我们众人所敬拜的一切存在的人，更能分裂祂的光荣和敬拜呢？所以我恳求你，千万小心，免得有人竟称你为假基督。你以为我是在辱骂诅咒吗？我所说的并非我自己的意见：看，福音使徒说：\BibleQuote{凡分裂耶稣的，不是出于天主；这便是假基督。}你若不分裂耶稣，不否认天主，便无人可称你为假基督。但若你否认，又为何责怪他人称你为假基督？就在你否认之时，我断定你已亲口说了。你愿知道这是否真实吗？请告诉我，当耶稣由贞女诞生时，你认为祂是人是天主？若只是天主，你当然分裂耶稣，因为你否认在祂内人性与天主性相结合。但若你说祂是人，你照样分裂祂，因为你亵渎地说（如你所愿）诞生的只是一个纯粹的人。也许你认为自己并未分裂耶稣，你否认祂是天主，即便你并不否认人与天主一同诞生，你也必然分裂祂。但或许你希望用实例使这更清楚。你要从两方面看实例。摩尼教徒在教会之外，他们宣称耶稣只是天主；厄彼翁派则说祂只是人。两者都否认并分裂耶稣：一说祂只是人，另一说祂只是天主。尽管他们的意见彼此对立，但这不同意见的亵渎却大致相同，除非能在邪恶的严重程度上加以区分：你断言祂只是人的亵渎，比说祂只是天主的更坏；因为虽然两者都错，但剥夺主的属于天主的性质比属于人性的性质更为侮辱。因此，唯有这才是公教与真信仰：即相信主耶稣基督既是天主，也是人；既是人，也是天主。\BibleQuote{凡分裂耶稣的，不是出于天主。}但分裂祂，就是试图撕裂在耶稣内合一的事物，割裂那唯一不可分的事物。在耶稣内合一而唯一的是什么呢？当然是人性与天主性。因此，谁若分割并撕裂这两者，就是分裂耶稣。否则，若他不撕裂割开，便不分裂耶稣；但若他撕裂，就确实分裂了祂。\par

% structure: p0032 source=h2 target=subsection reason=relative-heading-rank; base=h2

\subsection{第十一章}

\Emph{主降生成人的奥秘清楚表明基督的天主性。}\par

因此，对于每一个陷入这种疯狂亵渎的人，主耶稣在福音中亲自向法利塞人重复祂所说的话，并宣告说：\Quote{天主所结合的，人不可拆散。}\BibleRef{玛 19:6} 因为虽然这话最初由天主说出时似乎是针对另一件事，但天主的深奥智慧——祂所谈论的不只是血肉之事，更是属灵之事——愿意这话确实适用于那个主题，但更适用于这个主题。因为当当时的犹太人像你们一样相信耶稣只是一个没有神性的人时，主被问及关于婚姻结合的问题，在祂的教导中，祂不仅提及了那事，也提及了这事：虽然被咨询的是次要之事，祂的回答却适用于更大更深的事，当祂说：\Quote{天主所结合的，人不可拆散，}即不要拆散天主在我身上所结合的一切。不要让人类的邪恶拆散天主的神圣荣耀在我身上所结合的一切。但若你希望更充分地知道确实如此，请听宗徒谈论救主当时正在教导的这些主题本身，因为他作为天主所差派的教师，为了使他的听众——那些心智软弱的人——能够接受他的教导，解释了天主在奥秘中所宣告的这些主题本身。因为当他讨论血肉结合（救主在福音中被问及此事）的主题时，他重复了祂所强调的旧约中的那些段落，目的是让他们看到，既然他使用同样的权威，他就是在解释同样的主题。此外，为了使他的论证没有任何缺失，他加上了血肉结合的提及，并插入了丈夫和妻子的名称，劝勉他们彼此相爱：\Quote{丈夫们，你们应该爱妻子，如同基督爱了教会一样。} 又说：\Quote{同样，丈夫也应当爱妻子如同爱自己的身体；爱自己妻子的，就是爱自己。因为从来没有人恨过自己的血肉，而是培养抚育它，如同基督对待教会一样，因为我们都是祂身体的肢体。}\BibleRef{弗 5:25} 你看，他如何通过加上对基督和教会的提及，将对丈夫和妻子的提及引领众人离开血肉的理解，而领会属灵的意义。因为当他说完这一切后，他加上了主在福音中所引用的那些段落，说：\Quote{为此，人要离开父亲和母亲，依附自己的妻子，两人成为一体。} 此后，他特别强调地加上：\Quote{这是一个伟大的奥秘。} 他确实完全切断了并排除了任何血肉的解释，说这是一个神圣的奥秘。他在这之后又加上了什么？\Quote{但我是指基督和教会说的。} 也就是说：\Quote{但那是一个伟大的奥秘。但我是指基督和教会说的，}即，因为也许在目前所有人不能领会那更深的事，他们至少可以领会这一件，它与那事并不矛盾，也不相异，因为两者都指向基督。但因为不能领会那些更深的真理，至少让他们接受这些较容易的，好使他们通过抓住表面的事作为开始，得以进入更深的事，并且通过获得一些较为简单的事，为将来理解更深的事开辟道路。\par

% structure: p0035 source=h2 target=subsection reason=relative-heading-rank; base=h2

\subsection{第十二章}

\Emph{他更充分地解释了在以丈夫和妻子为名之下所象征的奥秘是什么。}\par

那么，那以男人和妻子之名所象征的伟大奥秘是什么呢？让我们问问宗徒本人吧，他在别处教导同一事物时，用了同样有力的词句，说：\Quote{确实，虔敬的奥秘是伟大的，它彰显于肉身，在圣神内被成义，被天使看见，被传于外邦人，在世上被相信，被接于荣耀中。}那么，那彰显于肉身的伟大奥秘是什么呢？显然是天主由肉身而生，天主以形体显现；祂公开地被接于荣耀中，正如祂公开地彰显于肉身一样。这就是那伟大的奥秘，关于它，他说：\Quote{因此，人要离开父亲和母亲，依附于妻子，二人成为一体。}那么，那二人成为一体的，是谁呢？天主和灵魂，因为在与天主结合的那个人类的肉身中，同时有天主和灵魂在场，正如主自己所说：\Quote{没有人能夺去我的生命（\OriginalTerm{灵魂}{anima}），是我自己舍弃的。我有权舍弃，也有权再取回。} \BibleRef{若 10:18} 因此，你看到在此三者：天主、肉身和灵魂。说话者是祂天主；祂在其中说话的是肉身；祂所谈论的是灵魂。那么，祂就是先知所说的那个人：\Quote{兄弟不能赎回，人也不能赎回}？那人，如经上所说，\Quote{升到了祂原先所在之处}，\BibleRef{若 6:62} 并且我们读到：\Quote{除了从天降下、仍在天上的人子，没有人升过天。} \BibleRef{若 3:13} 因此，我说，祂离开了父亲和母亲，即生祂的天主和\Quote{那作我们众人母亲的耶路撒冷}，\BibleRef{迦 4:26} 并依附于人类的肉身，如同依附于妻子。所以，在父亲方面，祂明确地说人要离开\Emph{他的}父亲，但在母亲方面，祂没有说\Quote{他的}，而只是说\Quote{母亲}：因为她不单是祂的母亲，更是所有信徒——即我们众人——的母亲。祂与妻子结合，因为正如男人和妻子成为一体，天主性的光荣与人的肉身也合一，二者——即天主和灵魂——成为一体。因为正如那肉身有天主居于其中，同样也有灵魂与天主同居于其中。这就是那伟大的奥秘，为了探究它，我们对宗徒的钦佩召唤我们，而天主自己的劝勉也吩咐我们：它并非与基督和教会无关，正如他说：\Quote{但我是指基督和教会说的。}因为教会的肉身是基督的肉身，而在基督的肉身中，有天主和灵魂同在：因此，同一位者临在于基督和教会中，因为我们在基督肉身中所信的奥秘，也由信德包含在教会内。\par

% structure: p0038 source=h2 target=subsection reason=relative-heading-rank; base=h2

\subsection{第十三章}

\Emph{古圣祖们渴望见到那奥秘启示的愿望。}\par

这个奥秘，既彰显于肉身，显现于世界，并被传于外邦人，许多古圣人在灵性中预见它时，曾渴望在肉身中见到它。因为\Quote{主说：我实在告诉你们，有许多先知和义人，渴望看你们所看见的，却没有看见；渴望听你们所听见的，却没有听见。} \BibleRef{玛 13:17} 因此，依撒意亚先知说：\Quote{愿你，上主，撕裂诸天，降临人间，} \BibleRef{依 64:1} 达味也说：\Quote{上主，求你使诸天低垂，降临人间。}梅瑟也说：\Quote{求你把你自己显示给我，让我能清楚地看见你。} \BibleRef{出 33:13} 从来没有人比梅瑟更接近从云中说话的天主，以及祂荣耀的真实临在，他领受了法律。如果说没有人比他更清晰地看见天主，那他为何还要祈求更清晰的显现，说：\Quote{求你把你自己显示给我，让我能清楚地看见你}？仅仅是因为他祈求那件事发生，正如宗徒用几乎相同的话告诉我们确实发生的：即主公开地彰显于肉身，公开地显现于世界，公开地被接于荣耀中；最终，圣人们得以用他们肉体的眼睛看见他们曾以灵性之眼预见的一切。\par

% structure: p0041 source=h2 target=subsection reason=relative-heading-rank; base=h2

\subsection{第十四章}

\Emph{他驳斥异端者邪恶亵渎的观念，即他们以为天主像在器具或雕像中那样住在基督内并藉祂说话。}\par

\Emph{否则}，如异端者所说，天主在主耶稣基督内如同在雕像或器具中，即祂仿佛住在一个人里面并藉一个人说话，而并非祂自己作为天主，在自己的身体内居住和说话：而实际上，祂早已这样住在圣人们里面，并藉圣人们说话。在我上面提到的那祈求祂来临的人中，祂也曾这样居住和说话。那么，如果他们在寻求之前已经领受了的，那他们祈求祂来临还有什么必要呢？或者，为什么他们渴望用眼睛看见内心已经拥有的东西，尤其是，一个人内在拥有同样的东西比在外面看见更好？如果天主住在基督内如同住在所有圣人们里面一样，那么为什么所有圣人们渴望看见基督而不是他们自己？如果他们只在耶稣基督身上看到他们自己所拥有的同一事物，那他们为什么不是更愿意在自己内拥有它，而不是在别人身上看见它？但你错了，可怜疯癫的人，\Quote{不明白，如宗徒所说，自己所讲论和所坚持的一切}：\BibleRef{弟前 1:7} 因为所有先知和圣人，都按各自所能担当的，从天主领受了圣神的一部分。但在基督内，全部天主性曾并仍以形体住在祂内。因此，他们都远不及祂的圆满，而从祂的圆满中领受了一些：因为他们得以充满，乃是基督的恩赐：因为如果没有祂作为万有的圆满，他们就都必然是空虚的。\par

% structure: p0044 source=h2 target=subsection reason=relative-heading-rank; base=h2

\subsection{第十五章}

\Emph{圣人们为默西亚来临的祈祷包含什么；以及他们那种渴望的性质是什么。}\par

这就是众圣人渴求的：他们为此祈祷。按照他们内心和心智的智慧程度，他们切望亲眼看见这事。因此，\Person{依撒意亚}先知说：\Quote{愿你撕裂诸天，降临下来。}见：\BibleRef{依 64:1} \Person{哈巴谷}也在宣告他人所渴求的同一件事，说：\Quote{当岁月临近，祢将显现自己；时届来临，祢将被彰显：天主将从\Place{特曼}而来，}或\Quote{天主将从南方而来。} \Person{达味}也说：\Quote{天主必会亲自来临：}又说：\Quote{祢坐在\OriginalTerm{革鲁宾}{Cherubim}之上者，请显现。}有些人宣告祂向世界显示的来临，另有人为此祈祷。有些人以不同形式，但都怀着同等的渴望：他们多少明白他们所祈求的是何等伟大的事，即那居于天主内、继续存于天主的形态和怀抱中的天主，竟会\Quote{空虚自己，}见：\BibleRef{斐 2:7}并取了仆人的形态，使自己忍受受难的一切痛苦和侮辱，并为祂的善而受惩罚，而最艰难、最可耻的是，祂竟死在那些祂将为之而死的人手中。所有圣人因此多少明白这事——我说多少，因为这事何其伟大，无人能全懂——他们用一致的声音，（可以这样说）经相互同意，都祈求天主的来临：因为他们确实知道，全人类的希望就在于此，所有人的救恩都系于此，因为除了那本身不受锁链束缚者，无人能解开囚犯；除了那本身无罪者，无人能释放罪人；因为任何人若不在某方面本身自由，就无法使别人自由。因此，当死亡临到众人，众人都缺乏生命，好使他们死于\Person{亚当}，而活于\Person{基督}。因为虽然有许多圣人，许多蒙选者，甚至天主的朋友，但他们中没有一人能凭自己得救，除非他们因主的来临和祂的救赎而得救。\par

\SourceRecord{Translated by C.S. Gibson.；Nicene and Post-Nicene Fathers, Second Series；Vol. 11.；Edited by Philip Schaff and Henry Wace.；Buffalo, NY: Christian Literature Publishing Co.,；1894.；<http://www.newadvent.org/fathers/35095.htm>.}

% structure: p0001 source=h1 target=section reason=page-title

\section{\WorkTitle{论降生成人（第六卷）}}

% structure: p0002 source=h2 target=subsection reason=relative-heading-rank; base=h2

\subsection{第一章}

\Emph{从五个大麦饼和两条鱼饱饫群众的奇迹，显示天主大能的威荣。}\par

我们在福音中读到，当主命人将五个饼带到祂面前时，不可胜数的天主子民藉此得以饱饫。但此事\Emph{如何}成就，却是无法解释、无法理解或想象的。天主大能的力量是如此伟大而不可理解，以致我们虽确知\Emph{事实}，却无法理解事实的\Emph{方式}。因为首先须领会，如此少的饼如何足以——我不说使他们吃饱，甚至只是分给他们摆放出来——何况人数比饼多出数千倍，而人群的组数几乎超过全部饼可能分成的碎块。因此，充足的供应是主的话的造化。工作在进行中增长。所见虽少，所赐予他们的却多到无法计算。故此，没有猜测、人的思辨或想象的余地。在这种情况下，唯一当行的是，如同忠信智慧之人，我们应当承认：无论天主所行之事何等伟大而不可理解，即使完全超出我们的领悟，也必须认识到在天主没有不可能的事。关于这些天主大能不可言喻的事迹，我们将按主题需要在后文更充分地讨论，因为它恰与祂神圣诞生的奥妙奇迹相符。\par

% structure: p0005 source=h2 target=subsection reason=relative-heading-rank; base=h2

\subsection{第二章}

\Emph{作者将七的数字（由五个饼和两条鱼组成）的奥秘应用于自己的著作。}\par

\Emph{同时}我们既已提及五个饼，我以为将它与我们已经写成的五卷书作一比较，并非不当。因为它们数目相等，性质也不无相似。因为正如饼是大麦的，这些书（就我的能力而言）也可公平地称为\Quote{大麦的，}尽管它们因引用圣经而丰富，在卑微的外表下含有赋予生命的宝藏。即便在这一点上，它们也与那些饼不无相似：它们看似贫乏，却证明充满祝福；同样，这些书虽就我的能力而言毫无价值，但因其中混合的神圣内容而宝贵；外形上因我的言辞而似大麦般卑微，内里却因主亲口的见证而有生命之粮的滋味。剩下的，是愿它们效法祂的榜样，藉神圣恩宠的恩赐，从无数种子中提供赋予生命的食粮。正如那些饼给吃的人带来体力，愿这些书给阅读的人带来属灵的力量。然而，正如当时主——这恩赐来自祂，正如那恩赐来自祂——藉那食物使众人得以饱饫而不至在路上昏厥，同样现在祂也能藉此使人饱饫而不致（从信仰中）迷失。然而，仍因为在那里，当不可胜数的天主子民以强大的恩赐被喂养时，虽然可吃的甚少，我们读到那五个饼之外又加了两条鱼，故此，我们这些渴望给所有跟随的天主子民提供属灵宴席滋养的人，也适宜相应于五个饼加上两卷书，相应于两条鱼加上两卷书：祈祷并恳求祢，主啊，愿祢垂顾我们的努力和祈祷，并赐予我们虔诚事业以顺利结局。既然我们出于爱情和服从，愿望我们的卷数能与饼和鱼的数目相符，求祢将祢祝福的德能赐予它们；并且，正如祢以福音的数目祝福我们这微小的工作，同样愿祢以福音的果实充满这数目，并赐予这一切成为祢教会全体子民——无论年龄性别——的神圣救恩之粮。若有人被那毒蛇的致命气息所影响，在不健康的灵魂和精神状态下，因薄弱的气质感染了瘟疫般的疾病，求祢赐予他们健康的全部力量和信仰的完全健全，好使祢藉这些著作赐予他们祢恩赐的救恩护理——正如福音中那食物被祢完全圣化，以致饥饿的灵魂因食用而坚强——同样，愿祢命令衰弱的灵魂藉此得医治。\par

% structure: p0008 source=h2 target=subsection reason=relative-heading-rank; base=h2

\subsection{第三章}

\Emph{他以安提约基雅会议的证词驳斥对手。}\par

\Emph{因此}，既然我们已在先前各卷书中，藉着神圣见证的份量，似乎已对那位否认天主的异端者作了完备的答复，现在让我们来探讨安提约基亚信经的信德及其价值。因为此人既在此信经中受洗并重生，就应被自己的宣认所驳倒，并（可谓）被自己的武器压垮；这方法正是：既然他已被圣经的证据定罪，如今也当由他自己口中的证据定罪。当他已清楚而明确地自证其罪时，便无需再用别的来反对他了。安提约基亚信经的文本及其信德如下：\Quote{我信唯一而真实的天主，全能的圣父，一切有形与无形之物的创造者。又信唯一的主耶稣基督，祂的独生子，一切受造物的首生者，在万世以前由祂所生，而非受造：出自真天主的真天主，与圣父同体：藉着祂，诸世界得以建立，万物得以造成。祂为了我们降临，由童贞玛利亚诞生，在本丢·彼拉多治下被钉十字架并埋葬：第三日按圣经复活了：升了天，还要再来审判生者死者，}等等。在这给予所有教会信德的信经中，我想知道：你宁愿追随人的权威，还是天主的权威？虽然我不愿苛刻或不仁慈地对你，而是让你自由选择你喜欢的任一选项；因此，若你接受其一，我便可否认另一。我问：我允许什么？我会强逼你接受两者之一，即使违背你的意愿。因为，你若愿意，就应出于自由意志明白信经中必有二者之一；若不情愿，则必须被迫看见它。因为，如你所知，信经（Symbolum）的名称来源于\Quote{集合}。在希腊文中称为\OriginalTerm{集合}{\GreekText{σύμβολος}}，拉丁文则称为\Quote{集合}。然而它是集合，因为当主的宗徒们将整个公教律法的信德汇集在一起时，所有那些以极丰富细节散布于整部圣经中的事理，都被汇集在信经无与伦比的简洁之中，正如宗徒的话所说：\Quote{完成祂的话，并在公义中缩短；因为主将在地上实行简短的话。\BibleRef{罗 9:28}} 这即是主实行的\Quote{简短的话}，祂以寥寥数语汇集了祂两约的信德，并以数句简短的条款包含了全部圣经的要旨，用祂自己的事培育祂自己的人，并以最简明扼要的公式赐予全部法律的效力。借此，祂如同最慈爱的父亲，顾念某些子女的疏忽与无知，使任何心智虽简单无知者，对如此易记之事也不至有任何困难。\par

% structure: p0011 source=h2 target=subsection reason=relative-heading-rank; base=h2

\subsection{第四章。}

\Emph{信经如何兼有神圣与人为的权威。}\par

你看，信经有天主的权威：因为\Quote{主将在地上实行简短的话}。但或许你还需要人的权威；这也不缺，因为天主藉着人实行了它。正如祂藉着圣祖、尤其是祂自己的先知，形成了整部圣经，同样，祂藉着自己的宗徒与司铎形成了信经。凡祂在这些（圣经）中以丰富而充足的素材所详述的，祂在此藉着自己的仆人，以最完备而扼要的方式囊括了。因此，信经毫无欠缺；因为它由天主的圣经藉着天主的宗徒而形成，它拥有它可能拥有的一切权威，无论人的或天主的。虽然如此，那由人所作成的，仍应被视为天主的工程；因为我们不当视之为藉以作成之人的工程，而应视之为实际造作者（天主）的工程。\Quote{我信，}于是信经说，\Quote{唯一真实的天主，全能的圣父，一切有形与无形之物的创造者；又信唯一的主耶稣基督，祂的独生子，一切受造物的首生者；在万世以前由祂所生，而非受造；出自真天主的真天主，与圣父同体；藉着祂，诸世界得以建立，万物得以造成；祂为了我们降临，由童贞玛利亚诞生；在本丢·彼拉多治下被钉十字架，并埋葬。第三日按圣经复活；升了天；还要再来审判生者死者，}等等。\par

% structure: p0014 source=h2 target=subsection reason=relative-heading-rank; base=h2

\subsection{第五章。}

\Emph{他以精选的论证攻击对手，并表明我们应持守由祖先所领受的宗教。}\par

如果你曾是亚略异端或撒伯流异端的维护者，并且不使用你自己的信经，我仍会用圣经的权威来驳斥你；我会用法律本身的言语来驳斥你；我会用那在整个世界都被认可的信经的真理来驳斥你。我会说，即使你缺乏理智和理解，你仍然至少应当跟随普世的共识：不应把少数恶人的歪曲见解看得比所有教会的信仰更重；这信仰由基督确立，由宗徒传递，应当被视为天主权威的声音，而天主权威的声音当然拥有天主的话语和心意。那么，如果我用这种方式对待你，又会如何？你会说什么？你会回答什么？我恳求你，岂不是会如此：即你没有被这样培养和教导；你的父母、主人和老师传授给你的是一些不同的东西。你在你父亲教导的聚会场所没有听到这个，在你领受圣洗的教会中也没有听到；最终，传授和教导给你的信经的文本和词语包含不同的内容。你是在它里面领受圣洗并重生的。你会说，你要持守你所领受的，并活在那信经中，你在其中得知你重生了。当你这样说时，我请问，你难道不认为你是在用一个非常坚固的盾牌来对抗真理吗？的确，即使是在恶事中，这也不是一个不合理的辩护，一个为错误提供不坏借口的辩护，如果没有把固执与错误结合在一起的话。因为如果你持守了你从小领受的，我们应当努力修正并改正你现在的错误，而不是严厉惩罚你过去的过错：然而现在，既然你出生在一个公教城市，受公教信仰的教育，并由公教圣洗重生，我怎能把你当作亚略派或撒伯流派来对待呢？我真希望你是！如果你是在错误中长大的，而不是从正确中堕落，我会少一些悲伤：如果你从未接受过信仰，而不是失去了它：如果你是一个老异端者而不是一个新背教者，因为你对整个教会造成的丑闻和伤害会少一些；最后，如果你能以一个平信徒而不是司铎的身份来试探教会，那会是一个不那么苦涩的悲伤和较少伤害的榜样。因此，如上所述，如果你曾是撒伯流主义或亚略主义或任何你喜欢的异端的追随者和维护者，你或许会躲在父母榜样、教师教导、周围人的交往和你信经的信仰下面。我问你，啊，异端者，没有什么不公平，没有什么苛刻。既然你在公教信仰中长大，就做你为错误信仰所会做的事吧。持守你父母的教导。持守教会的信仰：持守信经的真理：持守圣洗的救恩。你是什么奇事——你是什么怪物？你不会为自己做别人为他们的错误所做的事。但我们已说得够远了：出于对我们有联系的城市的爱，我们像顺风一样让悲伤发泄，而当我们急于前进时，却超出了正轨的目标。\par

% structure: p0017 source=h2 target=subsection reason=relative-heading-rank; base=h2

\subsection{第六章}

\Emph{他再次挑战他，要求他宣认\Place{安提约基雅}信经。}\par

那么，哦，你这个异端者，我们上面给出的信经文本，虽然是所有教会的（因为所有人的信仰只是一个），但特别属于\Place{安提约基雅}城和教会，即你被抚养、教导和重生的那个教会。这信经的信仰将你带到了生命之泉、救赎的再生、圣体圣事的恩宠、主的共融：还有更多！唉，这痛苦而可悲的抱怨！甚至到了职务的行使、长老的高位、司铎的尊严。你这个可怜疯狂的家伙，你认为这是轻微或琐碎的事吗？你没看到你所做的吗？你陷入了多深的深渊？失去了信经的信仰，你失去了一切你所是的。因为司铎职和你的救恩的奥迹基于信经的真理。你能否认这一点吗？我说你否认了你的自我。但也许你认为你不能否认自我。让我们查看信经的文本；这样如果你说你所习惯做的，你可能不被驳斥，但如果你说广泛不同和相反的事，你可能不被我看作被驳斥，因为你已经谴责了自己。因为如果你现在坚持与信经中以及你以前自己坚持的不同的东西，你如何能不将自己的惩罚归咎于自己，当你看到别人对你的看法和你自己的一样？\Quote{我信，}信经说，\Quote{在一个天主，全能的圣父，一切有形无形之物的创造者；并在主耶稣基督，祂的独生子，一切受造物的首生者；在万世以前由祂所生，而非受造。}你最好先回答这个：你承认耶稣基督天主子这些事，还是否认？如果你承认，一切都很好。但如果不是，你现在如何否认你自己以前承认的？那么选择你要的：两件事中必须有一件发生；即，你的那个宣认，如果仍然有效，应该独自释放你，或者如果你否认它，首先谴责你。因为你在信经中说：\Quote{我信唯一的主耶稣基督，祂的独生子，一切受造物的首生者。}如果主耶稣基督是独生子和一切受造物的首生者，那么通过我们自己的宣认，祂确实是天主。因为没有别的谁是独生子和一切受造物的首生者，除了天主的独生子：正如祂是受造物的首生者，祂也是万物的创造主天主。你怎么能说祂在由童贞玛利亚出生时只是一个凡人，而你却承认祂在万世以前就是天主？接下来信经说：\Quote{在万世以前由圣父所生，而非受造。}这信经是你念出的。你通过你的信经说，耶稣基督在万世以前由天主圣父所生，而非受造。信经说了任何关于那些幻影的事吗，那些你现在狂乱的东西？你自己说了关于它们的任何事吗？雕像在哪里？我请问，你的那个乐器在哪里？因为天主禁止这属于别人而不属于你。你在哪里断言主耶稣基督像一座雕像，所以你认为祂应该被敬拜不是因为祂是神，而是因为祂是天主的肖像；并且从光荣的主那里你制作了一个乐器，并且亵渎地说祂应该被敬拜不是为了祂自己，而是为了祂里面如同呼吸和发声的那位？你在信经中说主耶稣基督在万世以前由圣父所生，而非受造：这肯定只属于天主的独生子：祂的出生不是受造，并且祂可以被简单地说成是生而非受造：因为这与万物的本性及其尊荣相悖，即万物的创造者应该被相信是一个受造物；并且祂，一切有开端之物的作者，自己有一个开端，因为万物从祂开始。所以我们说祂是生而非受造：因为祂的生育是独特的，不是普通的受造。并且既然祂是天主，由天主所生，那被生的天主性必须具有那生者之威严所拥有的一切完备。\par

% structure: p0020 source=h2 target=subsection reason=relative-heading-rank; base=h2

\subsection{第七章。}

\Emph{他继续从\WorkTitle{安提约基雅信经}引出相同的论证线。}\par

但信经接着说道：\Quote{真天主，出自真天主；与圣父同一性体；万物都是藉着他造的，一切受造之物都是藉着他而有的。}当你宣认这一切时，要记得你是对主耶稣基督宣认的。因为你在信经中读到：你信主耶稣基督，天主的独生子，是首生于一切受造之物；在此之后及其他条文中说：\Quote{真天主，出自真天主，与圣父同一性体；万物都是藉着他造的。}那么，同一主体怎能既是天主又不是天主；既是天主又是雕像；既是天主又是工具呢？你这异端者，这些在任何一个主体中都不和谐，也不相合，以致你可以随意称他为天主，又随意认为同一主体是受造物。你在信经中说了\Quote{真天主}。现在你说：\Quote{只是一个凡人。}这些怎能相合协调，使得同一主体既是最伟大的能力，又是极致的软弱；既是至高荣耀，又仅是必死的呢？这些不能在同一主内相遇。这样，你将他在敬拜和贬抑上分开，一方面你可以随意尊崇他，另一方面可以随意损害他。当你领受真实救恩的圣事时，你在信经中宣认：\Quote{主耶稣基督，真天主，出自真天主，与圣父同一性体，世界的创造者，万物的造主。}你如今在何处！唉！你昔日的你何在？你的信德何在？你的宣认何在？你怎样败落，成了怪物和奇胎？是何等愚昧，何等狂妄使你毁灭？你将全能的天主化为了无生命的物质和受造的虚无。你的信德确实随着时间、年龄和司铎职分而增长。你年老时比昔年幼时更糟；如今身为老手比身为新手更糟；身为主教比作初学生更糟；你从未在开始成为教师之后学习过。\par

% structure: p0023 source=h2 target=subsection reason=relative-heading-rank; base=h2

\subsection{第八章}

\Emph{基督如何可以说是降来并由贞女所生。}\par

但是让我们看看接下来的部分。信经说：\Quote{主耶稣基督，真天主，出自真天主，与圣父同一性体；万物都是藉着他造的，一切受造之物都是藉着他而有的；}紧接着便最紧密地连上下文说：\Quote{祂为了我们降来，并由童贞玛利亚诞生。}那么，这位是真天主，与圣父同一性体，是万物的造主，我说，祂来到世界并由童贞玛利亚诞生，正如宗徒保禄所说：\BibleQuote{但时期一满，天主就派遣了自己的儿子，生于女人，生于法律之下。}\BibleRef{迦 4:4}你看信经的奥迹与圣经是如何相符的。宗徒宣称天主子是\Quote{由圣父派遣来的：}信经肯定祂\Quote{降来。}因为我们的信德理应宣认祂已\Quote{降来}，就是宗徒教导我们被派遣的那一位。然后宗徒说：\Quote{生于女人：}信经说：\Quote{生于玛利亚。}因此你看，藉着信经说话的是圣经本身，信经承认自己是从圣经而来的。但当宗徒说\Quote{生于女人}时，他恰当地用\Quote{生}来指\Quote{诞生}，按圣经的习惯，其中\Quote{生}代替\Quote{诞生}：就如这句：\Quote{代替你的父亲们，给你生下了儿子：}或者这句：\Quote{亚巴郎未有之前，我就是。}\BibleRef{若 8:58}这里我们清楚看到祂的意思是\Quote{在他诞生之前，我就是：}用\Quote{生}一词来指他的诞生，因为凡不需要被造的，都具有受造的真实性。于是它说：\Quote{祂为了我们降来并由童贞玛利亚诞生。}如果只是一个凡人由玛利亚诞生，怎能说祂\Quote{降来}呢？因为除了那有能力降来的，没有人\Quote{降来}。但是对于一个还没有存在的，他怎能有能力降来呢？可见，藉着\Quote{降来}这个词，表明那降来的已经存在：因为只有那已经存在的，才有能力降来，因为从他已存在的事实，才能有降来的可能。然而，一个凡人固然在受孕之前并不存在，因此他没有能力降来。显然，是天主降来了：祂属于既是\Emph{存在}又是\Emph{降来}的每一个情况。因为祂确实\Emph{降来}，因为祂\Emph{存在}；祂永远\Emph{存在}，因为祂能永远\Emph{降来}。\par

% structure: p0026 source=h2 target=subsection reason=relative-heading-rank; base=h2

\subsection{第九章}

\Emph{他又藉其异端者自己的宣认，定此人的死罪。}\par

但是，当事实足够清楚时，我们何必在言辞上争论？何必从信经的措辞中寻求对这问题的裁决，而信经本身就已经处理了这问题？让我们重复信经的宣认，也重复你自己的宣认（因为它是你的，正如它是信经的，因为你通过宣认它而使其成为你的），好让你看到你不仅背离了信经，也背离了你自己。信经说：\Quote{我信}，然后信经说：\Quote{在唯一真天主，全能圣父，一切有形与无形之物的创造者；并在主耶稣基督，祂的独生子，一切受造物的首生者；在万世之前由祂所生，而非受造；真天主，出自真天主；与圣父同一性体；万有都是藉祂而成，一切被造之物都是藉祂而造的。祂为了我们降临，并由童贞玛利亚降生为人。}信经接着说：\Quote{为了我们}，我们的主耶稣基督\Quote{降临并由童贞玛利亚降生，在\Person{般雀·比拉多}手下被钉十字架；被埋葬，并照圣经所载复活。}众教会不以宣认这些为耻；宗徒们不以宣讲这些为耻。你自己——我说的是你——如今你的每一句话都是亵渎，你如今否认一切，你从前并没有否认这些真理：天主降生，天主受难，天主复活。然后呢？你坠落到何处？你变成了什么？你沦落到了什么地步？你在说什么？你在呕吐什么？正如某人所说，就连疯狂的\Person{俄瑞斯忒斯}自己也会发誓这是疯子的言语。因为你在说什么？\Quote{那么，由\OriginalTerm{基督之母}{Christotocos}所生的天主子是谁？例如，如果我们说\InnerQuote{我信天主圣言，天主的独生子，由祂的圣父所生，与圣父同一性体，祂降下来并被埋葬}，我们的耳朵岂不被这声音所震惊？天主死了？}又说：\Quote{你说，在万世之前所生的那一位，有可能第二次降生并成为天主吗？}如果这一切都不可能，那么众教会的信经为何说它们确实发生了？你自己为何又说它们确实发生了？因为让我们比较你如今所说的和你从前所说的。你曾说过：\Quote{我信全能天主圣父；并信耶稣基督祂的子，真天主，出自真天主；与圣父同一性体；祂为了我们降临并由童贞玛利亚降生；在\Person{般雀·比拉多}手下被钉十字架；并被埋葬。}但如今你说的是什么？\Quote{如果我们说\InnerQuote{我信天主圣言，天主的独生子，由祂的圣父所生，与圣父同一性体，祂降下来并被埋葬}，我们的耳朵岂不被这声音所震惊？}你的话语的苦毒和亵渎，或许会驱使我们以狂怒凶猛的攻击来回应；但我们在某种程度上必须勒住我们虔诚悲伤的缰绳。\par

% structure: p0029 source=h2 target=subsection reason=relative-heading-rank; base=h2

\subsection{第十章}

\Emph{他谴责他，因为虽然他背弃了公教信仰，却仍然胆敢在教会内教导、献祭和作出裁决。}\par

那么，我恳求你，对你自己，我说。请告诉我，我祈求，若任何犹太人或外邦人否定了天主教信仰的信经，你认为我们应当听他吗？绝对不。若异端者或背教者同样做呢？我们更不应听他，因为一个人放弃他已认识的真理，比从未认识而否认更糟。于是我们看见你内有两人：一个天主教徒和一个背教者：先是天主教徒，后是背教者。你自己决定你认为我们应跟随哪一个：因为你不能强调你内一人的要求而不谴责另一人。那么你说是你从前的自我该受谴责：而你谴责天主教信经、所有的人的宣认和信仰？那又如何？啊，可耻的行为！啊，悲惨的悲痛！你在天主教会中做什么，你这位阻止天主教徒的人？为什么你，已否认了民众的信仰，仍污染民众的聚会：尤其敢站在祭台前，登上讲道台，向天主子民展示你无耻而奸诈的面孔——占住主教宝座，行使司铎职，自封为教师？教基督徒什么？不信基督：否认他们所在的神圣圣殿里那一位是天主。而在这一切之后，啊，愚妄！啊，疯狂！你幻想自己是教师和主教，而（啊，悲惨的盲目）你否认祂的天主性，祂的天主性（我重复）你自称是其司铎的那一位。但我们被悲痛所带走。那么，信经说什么？或者你自己在信经中说了什么？当然，\Quote{主耶稣基督，真天主出自真天主，与圣父同一性体；万有藉祂被造，一切受造之物由祂造成：}而同一为\Quote{我们而来，由童贞玛利亚诞生。}既然你曾说天主由玛利亚诞生，你怎能否认玛利亚是天主之母？既然你曾说天主来了，你怎能否认那来了的是天主？你在信经中说：\Quote{我信耶稣基督、天主子；我信真天主出自真天主，与圣父同一性体；祂为了我们而来，由童贞玛利亚诞生，并在般雀比拉多治下被钉十字架，并被埋葬。}但现在你说：\Quote{若我们说：我信天主圣言，天主的独生子，由圣父所生，与圣父同一性体，祂来了并被埋葬，我们的耳朵岂不被这声音震惊？}你看你如何彻底摧毁并抹煞整个天主教信经的信仰和天主教奥迹？\Quote{啊，罪，啊，畸形，当被逐出，}如某人所说，\Quote{到地极：}因为这话更真实地用于你，让你确实进入那孤寂之地，在那里你无法找到任何人去败坏。那么你以为我们救恩的信仰和教会希望的奥迹对你的耳朵和听力是震惊。而从前你急于受洗时，你听这些奥迹耳朵怎么没受损？当教会教师教导你时，你的耳朵怎么没坏？你那时当然毫无双重震惊地履行了你的职责，你的口和耳都未受损；当你重复从他人所听的，并作为说话者自己听自己说时。那时你的耳朵受的伤在哪里？你听力的震惊在哪里？你为何不反驳并大声反对？但确实，你随己意幻想，当你愿意时，是门徒；当你愿意时，是教会的敌人；当你愿意时，是天主教徒；当你愿意时，是背教者。一位配得上的领袖，确实，无论你依附哪边，都将教会引向你后面；使你的意愿成为我们生活的法律，并按你自身的变化改变人类，好让你不愿成为所有其他人所是的，他们可能成为你想要的！一种光辉的权威，确实，因为你现在不再是以前的你，世界就必须不再是它从前所是的！\par

% structure: p0032 source=h2 target=subsection reason=relative-heading-rank; base=h2

\subsection{第十一章}

\Emph{他驳斥异端者无声的反对，他们欲否认童年时所宣认的信仰。}\par

但也许你说，你重生时是个婴孩，因此当时不能思想或反驳。的确如此：你的幼年\Emph{确实}阻止了你反驳，那时若你是成人，你会因反驳而死。因为倘若在那最忠实而虔诚的基督教会中，司铎将信经交给望教者和见证的民众时，你曾试图在任何一点上住口或反驳，又如何？也许你已被听到，而不会立刻像某种新怪物或畸形物一样被作为瘟疫驱逐出去。不是因为那最热忱而虔诚的天主子民愿意沾染哪怕最坏之人的血：而是因为尤其是在大城市中，被天主之爱激发的民众见到有人反对他们的天主时，无法抑制他们信德的热忱。但就算如此。作为婴孩，倘若如此，你不能反驳并否认信经。为什么当你年长且强壮时你闭口不言？无论如何你长大了，成了人，并被安置在教会的职务中。经过所有这些年，经过所有职务和尊严的阶梯，难道你从未明白你早已教导的信仰？无论如何，你知道你是祂的执事和司铎。若救恩的规范对你是个困难，为什么你承担了你所不喜欢的信仰的尊荣？但确实，你是个有远见而单纯虔诚的人，你希望在两者之间平衡自己，好同时维护你邪恶的亵渎和公教的荣誉！\par

% structure: p0035 source=h2 target=subsection reason=relative-heading-rank; base=h2

\subsection{第十二章}

\Emph{被钉的基督，对那些宣称祂仅仅是人的，是绊脚石和愚妄。}\par

那时，你们的听觉和耳朵受到冲击，因为天主竟然诞生了，天主竟然受苦了。宗徒保禄啊，你的话在哪里？\BibleQuote{但我们所宣讲的，是钉在十字架上的基督：这为犹太人确是绊脚石，为外邦人是愚妄，但为那些蒙召的，不拘是犹太人或希腊人，基督却是天主的德能，天主的智慧。}\BibleRef{格前 1:23}天主的智慧和德能是什么？毫无疑问，就是天主。但他宣讲了被钉在十字架上的基督，作为天主的德能和智慧。那么，如果基督毫无疑问是天主的智慧，那么他无疑就是天主。\BibleQuote{我们}，他继续说，\BibleQuote{所宣讲的，是钉在十字架上的基督：这为犹太人确是绊脚石，为外邦人是愚妄。}所以，主的十字架，对外邦人是愚妄，对犹太人是绊脚石，对你们却两者兼具。的确，没有比不信更大的愚妄，也没有比拒绝聆听更大的绊脚石。那时，他们的耳朵因宣讲和天主的受难而受伤，正如你们现在受伤一样。他们和你们一样，认为这冲击了他们的耳朵。因此，当宗徒宣讲基督为天主时，一听到我们的天主和主耶稣基督的名字，他们就塞住头上的耳朵，正如你们塞住悟性的耳朵一样。你们两人在这件事上的罪过似乎相等，要不是你们的过犯更大：因为他们否认了他，在他的苦难中仍显示人性；而你们却否认了他，因为复活已证明他是天主。所以，他们在地上迫害他，而你们甚至在天上迫害他。不仅如此，这更为残忍和邪恶，因为\Emph{他们}在无知中否认他，\Emph{你们}却是在接受了信德之后否认他；\Emph{他们}不认识主，\Emph{你们}却是在宣认他为天主之后；\Emph{他们}以对法律的热情为掩饰，\Emph{你们}以你们的主教职为掩饰；\Emph{他们}否认了他们以为不相干的那一位，\Emph{你们}却否认了你们身为司铎的那一位。啊，不堪的行为，前所未闻！你们迫害和攻击的，正是你们仍在执行的职务所事奉的那一位。\par

% structure: p0038 source=h2 target=subsection reason=relative-heading-rank; base=h2

\subsection{第十三章。}

\Emph{他回应了那个反对意见，即他们说所生的孩子应当与母亲\OriginalTerm{同一实体}{one substance}。}\par

但确实，在你们的欺骗和亵渎中，你们用了一个重要的论证来否认和攻击主天主，即你们说：\Quote{所生的孩子应当与母亲\OriginalTerm{同一实体}{one substance}。}我不完全承认这一点，并主张在关于天主诞生之事上，这并不成立；因为诞生与其说是那生育者的工作，不如说是她儿子的工作；他按自己的意愿诞生，他使其成就。其次，如果你们说所生的孩子应当与父母\OriginalTerm{同一实体}{one substance}，我肯定主耶稣基督与他的圣父\OriginalTerm{同一实体}{one substance}，也与他的母亲\OriginalTerm{同一实体}{one substance}。因为按位格的不同，他显示了与双亲的相似。按天主性，他与圣父\OriginalTerm{同一实体}{one substance}；按肉体，他与母亲\OriginalTerm{同一实体}{one substance}。并非一位与圣父\OriginalTerm{同一实体}{one substance}，另一位与母亲\OriginalTerm{同一实体}{one substance}，而是因为同一主耶稣基督，作为人诞生，同时也是天主，在他内兼具双亲的特性：作为人，他显示与人类母亲的相似；作为天主，他拥有天主圣父的本性。\par

% structure: p0041 source=h2 target=subsection reason=relative-heading-rank; base=h2

\subsection{第十四章。}

\Emph{他将这种错误观点与白拉奇派的教导相比较。}\par

\Emph{否则}，如果由玛利亚诞生的基督与出自天主的基督不是同一位，你们无疑就制造了两个基督，正如白拉奇那可憎的谬误一样；这谬误声称童贞玛利亚所生的是一个纯粹的人，并说他是人类的教师而非救赎者；因为他带给人的不是生命的救赎，而只是生活的榜样，即人通过跟随他，做同样的事，从而达到相似的状态。因此，你们的亵渎只有一个来源，谬误的根也是同一的。他们主张由玛利亚所生的是一个纯粹的人；你们也主张同样的事。他们将人子与天主之子分开；你们也这样做。他们说救主在受洗时成为基督；你们说他在受洗时成为天主的宫殿。他们不否认他在受难后成为天主；你们甚至在他升天后仍否认。因此，你们的乖谬只有一点超过了他们：因为他们似乎在地上亵渎主，而你们甚至在天上。我们不否认你们超越了你们所效仿的人。他们最终不再否认天主；你们却永不停止。虽然他们的宣认绝不能视为真实的，因为他们只允许救主在受难后享有天主性的光荣，而否认他在此之前是天主，只在此之后宣认；因为在我看来，那在有关天主的事上否认某一部分的人，就是完全否认了他；那不承认他永远存在的人，就是永远否认了他。正如你们一样，即使你们现在承认在诸天中由童贞玛利亚诞生的主耶稣基督是天主，你们也没有真正宣认他，除非你们承认他永远是天主。但事实上，你们根本不想在任何方面改变或改变你们的意见。因为你们断言，你们称为纯粹人所生的他，目前仍不是天主。啊，新奇而奇妙的亵渎！虽然你们与异端者一同断言他是人，却不同异端者一同宣认他是天主！\par

% structure: p0044 source=h2 target=subsection reason=relative-heading-rank; base=h2

\subsection{第十五章。}

\Emph{他指出那些支持这种错误教导的人承认两个基督。}\par

然而，我已开始说，既然你确实得出两位基督，此事必须加以说明和澄清。请问，我祈求你，你这将基督与天主子分离的人，你怎能信经中宣认基督是由天主所生？因为你说：\Quote{我信天主圣父，并信耶稣基督，祂的圣子。} 那么，此处你有耶稣基督——天主子；但你却说由玛利亚所生的，并非那同一位天主子。因此，有一个属天主的基督，另一个属玛利亚的。按你的看法，便有两位基督。因为，虽在信经中你并不否认基督，但你说玛利亚的基督与你信经中所宣认的那位不同。但或许你说基督并非由天主所生：那么你怎能在信经中说：\Quote{我信耶稣基督，天主子？} 那么你或是否认信经，或是承认基督是天主子。但若你在信经中承认基督是天主子，你也必须承认那同一位基督、天主子，是出自玛利亚。或者，你若造出一个出自玛利亚的另一个基督，你确实就做出亵渎的主张，说有两个基督。\par

% structure: p0047 source=h2 target=subsection reason=relative-heading-rank; base=h2

\subsection{第十六章。}

\Emph{他进一步说明这教导破坏了天主圣三的宣认。}\par

但即使你的固执和狡诈并未被信经的这信仰所约束，难道，我问你，你没有被诉诸理性和真理之光的论证所压倒吗？请告诉我，无论你是谁，哦，你这异端者——至少我们相信并宣认一个天主圣三：圣父、圣子、圣神。关于圣父和圣神的光荣，没有问题。你在诽谤圣子，因为你说由玛利亚所生的，并非同一位格，与那由天主圣父所生的位格相同。那么请告诉我：既然你并不否认天主的独生子是由天主所生，那么你究竟说由玛利亚所生的是谁？你说\Quote{一个纯粹的人}，按照祂自己所说的：\BibleQuote{那由肉生的，就是肉。}见：\BibleRef{若 3:6} 但是，那不仅按人类生育之律而生的一位，不能被称为纯粹的人。因为天使说：\BibleQuote{那在她内受孕的，是出于圣神。}见：\BibleRef{玛 1:20} 这一点甚至你也不敢否认，尽管你几乎否认了所有救恩的奥迹。既然祂是由圣神所生，不能被称为纯粹的人，因为祂是因天主的灵感而受孕；如果祂不是那一位——如宗徒所说：\Quote{祂空虚了自己，取了奴仆的形体}，并且\Quote{圣言成了血肉}，并且\Quote{祂贬抑自己，服从至死}，并且\Quote{祂为我们成了贫穷，祂本是富有的}——那么请告诉我，那由圣神所生、由天主的庇荫而受孕的，祂是谁？你说祂肯定是另一个位格。那么，就有两个位格：即一位，是在天上由天主圣父所生的；另一位，是由玛利亚因天主的灵感而受孕的。这样，你就引入了第四位格，并且（虽然你口头上称祂为纯粹的人）你实际上断言祂并非纯粹的人，因为你承认（虽然并非适当地）祂应受尊荣、敬拜和钦崇。既然那由圣父所生的天主子确实应受敬拜，那由玛利亚因圣神而受孕的也应受敬拜，你就造出了两位应受尊荣和崇敬的位格，你将他们彼此分离，以致各以特殊专属的尊荣敬拜。这样，你看到，因否认圣子并使之与自身分离，你尽你所能摧毁了整个天主性的奥迹。因为你正试图在天主圣三中引入第四位格，但你看到你已经完全否认了整个天主圣三。\par

% structure: p0050 source=h2 target=subsection reason=relative-heading-rank; base=h2

\subsection{第十七章。}

\Emph{那些在天主教信仰的一点上犯错的人，丧失整个信仰，以及信仰的全部价值。}\par

既然是这样，你若否认耶稣基督天主子是唯一的，就否认了一切。因为教会奥迹和天主教信仰的计划是：一个人若否认神圣奥迹的一部分，就不能宣认另一部分。因为所有部分如此紧密相连且结合一起，以致一个不能离开另一个而成立；若一个人否认总体中的一点，他相信所有其他部分也是无用的。因此，你若否认主耶稣基督是天主，结果在否认圣子的同时也否认了圣父。正如圣若望所说：\BibleQuote{那没有子的，也没有父；但那有子的，也有父。}见：\BibleRef{若一 2:23} 因此，否认那被生的一位，也就否认了那生的一位。同样，否认圣子生于血肉，也导致你否认祂生于圣神，因为同一位格，祂先生于圣神，后又生于血肉。你若不信祂生于血肉，结果就是你也不信祂受苦。你若不信祂的苦难，除了否认祂的复活之外，你还能剩下什么？因为对复活者的信仰源于对死亡者的信仰。复活的印证也不能保持，除非对祂死亡的信仰先于它。因此，你若否认祂的苦难和死亡，也就否认了祂从地狱的复活。随之，你当然也否认祂的升天，因为没有复活就没有升天。我们若不信祂复活了，我们也不能信祂升天了：正如宗徒所说：\BibleQuote{那下降的，正是那上升的同一者。}见：\BibleRef{弗 4:10} 这样，就你而言，主耶稣基督既没有从地狱复活，也没有升天，也没有坐在天主圣父的右边，也不会在我们所期待的审判之日来临，也不会审判生者与死者。\par

% structure: p0053 source=h2 target=subsection reason=relative-heading-rank; base=h2

\subsection{第十八章。}

\Emph{他将讲论指向与他辩论的对手，并恳求他恢复理智。和好圣事对于跌倒者的救恩是必要的。}\par

因此，你这可怜、疯狂、顽固的受造物，你看到你已经彻底推翻了信经的全部信德，以及我们望德和奥迹中一切宝贵的东西。然而你仍然敢留在教会内，并自以为是一位司铎，尽管你否认了使你成为司铎的一切。那么，回到正路上来，恢复你从前的思想，回到理智中来，如果你曾经有过任何理智。回到你自身，如果你里面曾经有一个你可以回到的自身。承认你救恩的圣事，你藉以被入门和重生的圣事。它们如今对你并不比那时用处更少；因为它们如今能藉补赎使你重生，正如那时它们通过圣洗池生了你。牢牢持守信经的完整纲要。持守信德的全部真理。相信天主圣父：相信天主圣子：相信一位生者和一位被生者，万有的主，耶稣基督；与圣父同性同体；生于祂的天主性中；生于肉身：双重出生，却只有同一荣耀；祂自己是万有的创造者，由圣父所生，后来又由童贞女所生。\par

% structure: p0056 source=h2 target=subsection reason=relative-heading-rank; base=h2

\subsection{第十九章。}

\Emph{基督在时间中的诞生并未减损祂天主性的荣耀与权能。}\par

因为祂由肉身而来并在肉身中，这涉及祂的诞生，在祂身上并无减损：祂只是被生，而非变得更糟。因为，尽管祂仍处于天主的形象中，却取了仆人的形象，然而祂人性构造的软弱并未影响祂作为天主的本性；反之，当祂天主性的能力保持完整无损时，祂人性肉身中所发生的一切皆是祂人性的进展，而非祂荣耀的减损。因为当天主在人类肉身中诞生时，祂并非以那种方式在人类肉身中诞生以至于祂在自己内不再保持天主性，而是，当天主性保持如前时，天主成了人。因此，当玛尔大用肉眼看见那人时，她藉着灵性的眼光承认祂是天主，说：\Quote{是的，主，我信祢是基督，永生天主之子，来到这世界上。}\BibleRef{若 11:27}同样，伯多禄因着圣神的启示，外表上看见人子，却宣认祂是天主子，说：\Quote{祢是基督，永生天主之子。}\BibleRef{玛 16:16}同样，多默触摸到肉身时，相信他触摸到了天主，说：\Quote{我的主，我的天主！}\BibleRef{若 20:28}因为他们都承认只有一个基督，以免使祂成为两个。因此，你相信祂吧；并且如此相信：万有之主耶稣基督，既是独生者又是首生者，既是万物的创造者又是人类的保存者，并且同一位首先是整个世界的创造者，后来成为人类的救赎者？祂仍与圣父同在并在圣父内，与圣父同性同体，却（如宗徒所说）\Quote{取了仆人的形象，贬抑自己，听命至死，且死在十字架上：}\BibleRef{斐 2:7}并且（如信经所说）\Quote{由童贞玛利亚诞生，在般雀比拉多治下被钉十字架，被埋葬。第三日祂按照圣经复活了；升了天；将来还要再来审判生者与死者。}这就是我们的信德；这就是我们的救恩：相信我们的天主和主耶稣基督在万有之前和万有之后乃是同一位。因为经上记载：\Quote{耶稣基督昨天、今天、直到永远，常是一样。}\BibleRef{希 13:8}因为\Quote{昨天}意指一切过去的时间，在其中，祂在起初之前由圣父所生。\Quote{今天}涵盖这世界的时间，在其中祂再次由童贞女所生，受难，并复活。但\Quote{直到永远}这一表达指整个无尽的将来永恒。\par

% structure: p0059 source=h2 target=subsection reason=relative-heading-rank; base=h2

\subsection{第二十章。}

\Emph{他由前述证明，我们并非说天主在万世以前是必死的或有肉身，尽管基督——从永远就是天主并在时间中成为人——只是一个位格。}\par

但也许你会说：如果我承认在时间的终结由童贞女所生的那一位，与在万有之前由天主圣父所生的那一位是同一位，我便暗示在世界开始以前天主已在肉身之中，如同我说祂后来成为了人，而祂一直是天主；因此我会说后来出生的那个人一直存在。我不愿你被这种盲目的无知和模糊的误解所困扰，以致以为我在主张由玛利亚所生的人性在万有开始之前就已存在，或断言天主在世界开始以前始终以形体存在。我不说，我重复，我不说那人性在出生之前就已经在天主内：而是说天主后来在人性中出生。因为由童贞女肉身所生的那肉身并非一直存在；但一直存在的天主，藉童贞女肉身的血肉来到人的肉身中。因为\Quote{圣言成了血肉}，并非将血肉与祂一同显明；而是在天主性的荣耀中将自己与人的肉身结合。因为请告诉我，圣言何时何地成了血肉，或祂何时取了仆人的形状而空虚了自己，或祂何时成了贫穷的，尽管祂是富有的？难道不是在童贞女的圣胎中吗？在那里，在祂降生成人时，天主的圣言被说成成了血肉；在祂诞生时，祂真正取了仆人的形状；并且当祂在人性中钉在十字架上时，祂成了贫穷的，并在肉身的苦难中成为贫穷的，尽管祂在天主性的荣耀中是富有的？否则，如果如你所说，在某段较晚时期，天主性进入祂内，如同进入一位先知或圣人那样，那么那些祂乐意居住的人里面，\Quote{圣言也成了血肉}；在每一个人里面，祂空虚了自己并取了仆人的形状。因此，在基督内便没有新事或独特之处。祂的受孕、诞生和死亡都没有任何特殊或神奇之处。\par

% structure: p0062 source=h2 target=subsection reason=relative-heading-rank; base=h2

\subsection{第二十一章。}

\Emph{圣经的权威教导基督从永远就已存在。}\par

然而，为了回到我们先前所说的，尽管这一切都如我们所陈述的那样：我们如何读到耶稣基督（你们断言祂只是一个凡人）甚至在祂由童贞女诞生之前就已存在，以及祂如何被先知和宗徒宣称为天主，甚至在世界之前？正如保禄所说：\Quote{只有一个主耶稣基督，万有都藉着祂而有。}\BibleRef{格前 8:6} 他在别处又说：\Quote{因为在天上和地上的一切受造物，无论看得见的还是看不见的，都是在基督内受造的。}\BibleRef{哥 1:16} 那由人及天主权威所制定的信经也说：\Quote{我信天主圣父，并信主耶稣基督，祂的圣子。} 在其他条款之后：\Quote{真天主出自真天主；藉着祂，世界被创造，万有得以形成。} 又说：\Quote{祂为了我们而来，由童贞玛利亚诞生，被钉十字架，并被埋葬。}\par

% structure: p0065 source=h2 target=subsection reason=relative-heading-rank; base=h2

\subsection{第二十二章}

\Emph{位格结合\OriginalTerm{位格结合}{hypostatic union}使我们能够将属于基督肉身的属性归于天主。}\par

那么，基督（你们称之为一个凡人）如何被圣经宣称是无起始的天主，如果我们自己承认主的人性在祂由童贞女成孕诞生之前并不存在？我们又如何读到人和天主如此紧密的结合，以至于看起来人永远与天主同永恒，而后天主又与人同受苦？而我们无法相信人没有起始，或天主能够受苦。这正是我们在先前著作中所确立的：即，天主与人性的结合——也就是与祂自己的身体结合——不容许人们在思想中将人与天主分离。祂也不允许任何人认为人子是一个位格，天主子是另一个位格。反而，在全部圣经中，祂将主的人性结合并如同纳入天主性中，以至于无人能在时间上将人与天主分离，在祂的苦难中将天主与人分离。因为如果你在时间中看待祂，你会发现人子永远与天主子同在。如果你留意祂的苦难，你会发现天主子永远与人子同在，并且基督——人子和天主子——是如此的一体不可分，以至于按照圣经的说法，人不能从时间上与天主分离，天主也不能在祂的苦难中与人分离。由此而来：\Quote{没有人升过天，除了那从天上降下来的人子，祂仍在天上。}\BibleRef{若 3:13} 在此，天主子在地上说话时，证实人子在天上；并证实那同一位人子，祂说将要升天，先前已从天上降下。又有：\Quote{如果你们看到人子升到祂原先所在之处，又怎样呢？}\BibleRef{若 6:63} 宗徒在思考时间中所发生的事时，也说万有都是藉着基督造的。因为他说：\Quote{只有一个主耶稣基督，万有都藉着祂而有。}\BibleRef{格前 8:6} 但当论及祂的苦难时，他指出荣耀的主被钉十字架。\Quote{因为，}他说，\Quote{如果他们认识，就绝不会把荣耀的主钉在十字架上。} 同样，信经论及那唯一和首生的主耶稣基督，\Quote{真天主出自真天主，与圣父同性同体，万物的创造者}，肯定祂由童贞女诞生、被钉十字架并后来被埋葬。这样，正如将天主子与人之子结合在一个身体中，并将天主与人联合，以至于在时间上或在苦难中都不能分离，因为主耶稣基督被显示为同一的位格，既作为天主而永恒，又作为人而忍受祂的苦难；尽管我们不能说人没有起始或天主可受难，但在主耶稣基督的独一位格中，我们可以说人是永恒的，以及天主是死亡的。你看见，基督意指整个位格，这名字代表两种性体，因为人和天主都被生，所以它包含了整个位格，以至于当这个名称被使用时，我们看到没有任何部分被遗漏。在童贞女诞生之前，人性并不像天主性那样拥有从过去来的永恒；但因为天主性在童贞女的胎中与人性结合，所以当我们使用基督这个名字时，就不能不提及另一方而谈及另一方。\par

% structure: p0068 source=h2 target=subsection reason=relative-heading-rank; base=h2

\subsection{第二十三章}

\Emph{提喻\OriginalTerm{提喻}{Synecdoche}（以部分代表整体）这一修辞手法在圣经中非常常见。}\par

那么，凡你对主耶稣基督所说的一切，都是就整个位格而言；提到圣子，你就提到人子；提到人子，你就提到圣子：这是通过语法修辞手法\OriginalTerm{提喻法}{synecdoche}，即由局部理解整体，或局部代指整体。\newline{}圣经的确证明了这一点，因为主在其中常使用这种修辞手法，这样教导关于他人的事，也愿意我们以同样方式理解关于祂自己的事。因为有时日、事物、人、时期在圣经中以别的方式表示。例如，天主宣告以色列要服事埃及人四百年，并向亚巴郎说：\BibleQuote{你要知道，你的后裔必会在异地作客，他们要奴役他们、压迫他们四百年。}\BibleRef{创 15:13}\newline{}然而，若你考虑天主说话之后的总时间，则超过四百年；但若只计算他们受奴役的时间，则不足四百年。事实上，在给出这个时期时，除非你这样理解，我们便不得不认为天主的话说了谎（愿这种思想远离基督徒的头脑！）。但因从天主发言之时起，他们生活的整个时期共计超过四百年，而他们的奴役持续了不到四百年，所以你必须理解，这里是以部分代整体，或以整体代部分。还有一种类似的方式来表示日和夜，即当意指某一昼夜中的一天时，每一时期都由该单一时期的一部分来表示。确实，用这种方式，主苦难的时间难题就得以澄清：因为主曾预言，如同约纳先知的预像，人子要在在地中心三天三夜，\BibleRef{玛 12:40}\newline{}而在祂被钉死的第六日之后，祂只在地狱待了一天两夜，我们如何证明天主的话的真实性呢？当然是靠提喻法，即因为祂被钉死的那一天属于前一夜，而祂复活的那一夜属于次日；因此，当加上属于那一日的前一夜，以及属于那一夜的后一日，我们便看到整个时期毫无缺失，它是由各部分组成的。圣经中充满了这类表达方式，但若全部列举则过于冗长。例如，圣咏说：\BibleQuote{人算什么，你竟顾念他？}这里我们从局部理解整体，因为只提到一个人，却意指全人类。同样，阿哈布犯罪时，我们被告知百姓犯罪了。在那里，虽然提及全体，却要从整体理解局部。主的先驱若翰也说：\BibleQuote{在我以后要来的那个人，因他原在我以前，就成了在我以前的。}\BibleRef{若 1:15}\newline{}那么，祂如何意指祂要在他之后来，却显明祂在他之前呢？因为若这是指以后将出生的人，他如何在他之前？但若这是指圣言，又怎能说\Quote{有\Emph{一个人}在我以后来？}除非在唯一的主耶稣基督身上，既显出人性之在后，又显出神性之在先。因此，同一主既在他之前，又在他之后：按肉身，祂在时间上后于若翰；按神性，祂在万有之先。所以当他提到那个人时，他同时指人性与圣言，因为主耶稣基督，天主之子，既完全具有人性又完全具有神性，在祂内提及其中一种本性时，便指称了整个位格。此外还有什么必要呢？我想，我若试图收集或讲述有关此主题的一切可说之事，时日必不敷用。我们已说的已经足够，至少在这一部分主题上，既足以阐明信经，又符合我们论证的需要，也符合本书的篇幅。\par

\SourceRecord{Translated by C.S. Gibson.；Nicene and Post-Nicene Fathers, Second Series；Vol. 11.；Edited by Philip Schaff and Henry Wace.；Buffalo, NY: Christian Literature Publishing Co.,；1894.；<http://www.newadvent.org/fathers/35096.htm>.}

% structure: p0001 source=h1 target=section reason=page-title

\section{\WorkTitle{降生成人}（第七卷）}

% structure: p0002 source=h2 target=subsection reason=relative-heading-rank; base=h2

\subsection{第一章}

\Emph{当他准备回应对手的诽谤时，他恳求天主恩宠的帮助，传授一个祈祷文，供那些与异端者辩论者使用。}\par

正如那些逃脱海上危险的人，仍在害怕港口的沙滩或岸边的礁石一样，我的情况也是如此——由于我把异端者的诽谤留到最后——尽管我已达到了自定工作的终点，却开始害怕我曾渴望到达的结束。但是，如先知所说：\BibleQuote{主是我的助佑；我不惧怕人能对我做什么。} 因此，我们不怕狡猾的异端者在我们面前挖掘的陷阱，也不怕布满可怕荆棘的道路。因为他们使我们的道路变得困难，但没有封闭它，我们面前有清除它们的劳苦，而非无法做到的恐惧。因为当我们软弱地行走在正道上时，他们挡在路上，惊吓行人而非伤害他们，我们的工作与事务更多在于清除它们，而非因困难而恐惧。因此，我们将手放在这致命巨蛇的恐怖头上，渴望抓住缠绕在其巨大躯体蜷曲盘绕中的所有肢体，我们一再向祢祈祷，主耶稣啊，我们一向向祢祈祷，求祢开口赐给我们言语，\BibleQuote{为了摧毁堡垒，摧毁阴谋，以及一切高举起来反对天主知识的自高之事，并将一切理智掳去服从祢。} \BibleRef{格后 10:4} 因为那开始被祢俘获的人，才是真正自由的。求祢临于此祢的工作，以及那些为祢而奋斗、超越其力量的祢的属下。求祢使我们击碎这新蛇张开的嘴，以及它充满致命毒液的脖颈，祢使信者的脚无恙地踏在蛇和蝎子上，行走在毒蛇和蜥蜴上，践踏狮子和毒龙。并赐予，藉着坚定无邪的无畏胆量，\BibleQuote{吸奶的孩童可以在毒蛇洞上嬉戏，断奶的幼儿可以将手伸入毒蛇的穴中。}\BibleRef{依 11:8} 求祢也赐予我们，使我们无恙地将手伸入这可怕且最邪恶的毒蛇的穴中；若它在任何洞穴中，即人的心中，有藏身处或休息处，或在彼产卵，或留下其泥泞路径的痕迹，求祢从它们中除去这最有害毒蛇的一切污秽和致命污染。除去他们的亵渎带给他们的不洁，并用\BibleQuote{祢神圣洁净的簸箕}\BibleRef{拉 3:2} 净化那些陷入恶臭泥沼的灵魂，使\BibleQuote{贼窝}成为\BibleQuote{祈祷之所：} \BibleRef{玛 21:13} 并使那些如今——如经上所写——是刺猬、怪物、野山羊和各样奇异生物栖息之所的地方，在那里祢圣神的恩赐，即信德和圣洁之美，得以照耀。正如祢曾摧毁偶像崇拜，驱逐偶像，并将魔鬼的庙宇转变为美德的神殿，让闪烁的光芒照入蛇和蝎子的洞穴，将错误和羞耻的洞穴变为美丽和辉煌的家园，同样求祢将祢的慈悲和真理之光倾注在所有被异端顽固的黑暗蒙蔽双眼的人身上，使他们最终能以清澈无遮的目光，瞻仰祢降生成人的伟大而赋予生命的奥迹，从而认识祢作为\OriginalTerm{真人}{Very man}，由那神圣的童贞女之胎所生，并承认祢永远是\OriginalTerm{真天主}{Very God}。\par

% structure: p0005 source=h2 target=subsection reason=relative-heading-rank; base=h2

\subsection{第二章}

\Emph{他回应从这些话而来的反对意见：\Quote{没有人生育过在她之前就已存在的人。}}\par

在我开始讲述前几卷中未曾提及的事之前，我认为应当先尝试履行我已作出的承诺，以便在完全兑现诺言之后，可以更自由地谈论尚未涉及的内容。因此，那条新蛇为摧毁圣诞的信德，向天主的教会嘶喊说：\Quote{没有人会生出比自己更年长的人。}首先，我认为你既不知道你在说什么，也不知道这话从何而来。因为如果你知道或明白这话从何而来，就绝不会按照人的幻想看待天主独生子的诞生，也不会试图用纯粹人的论断来判定那位不由男子受孕而生者的事，更不会以人的不可能之事来反驳天主的全能，你若知道在天主前没有不可能的事。你说，没有人会生出比自己更年长的人。那么，我请你告诉我，你是指哪些情形而言？你以为你能为哪类受造物的本性制定规则？你以为你能为人类、走兽、飞鸟或牲畜定法律吗？只有这些（以及类似的事物）才能做出这样的断言。因为没有一种受造物能生出比自己更年长的；已经产生的事物不可能重新回到它那里，再被新生一次。所以没有人能生出比自己更年长的，如同没有人能生出比自己更年长的一样；因为生育的时机只能出现在有可能受孕的地方。那么，你难道认为在论及全能天主的诞生时，也须顾及与地上受造物出生相同的考虑吗？你竟以人的处境之本性作为困难，加于那本性自身的创造者吗？你看，正如我上述所说，你不知道你谈论的是谁或关乎谁，因为你将受造物与造物主相比；为了衡量天主的能力，你从那些若非来自天主就根本不会存在的事物中举例。天主是按照自己的意愿、在自己所愿的时候、从自己所愿的那一位而来的。时间、人物、人的方式、受造物的惯例，对祂都没有阻碍；因为受造物的法则不能阻挡祂，祂是万有的创造者。凡祂愿意可能的事，都已在祂手中，因为愿意的能力属于祂。你想知道天主的全能扩展到多远、有多大吗？我相信主在祂的受造物身上甚至能做到你不相信祂在自己身上能做到的事。因为现今所有生出比自己年幼的活物，只要天主发命，就能生出比自己年长得多的活物。因为甚至食物和饮料，如果天主愿意，也能化为胎儿和后代；甚至从万物的起始一直流动、一切活物所使用的水，只要天主发命，就能在母腹中形成身体，并生出。因为谁能给天主的作为设定界限，或限制天主眷顾？或者，谁（借用圣经的话）能向祂说：\Quote{你做什么？}如果你否认天主能做一切事，那么当天主降生时，你也要否认比\Person{玛利亚}更年长的一位能由她而生。但如果在天主没有不可能的事，你为什么以不可能之事来反驳祂的降临？你知道对祂而言没有一件事是不可能的。\par

% structure: p0008 source=h2 target=subsection reason=relative-heading-rank; base=h2

\subsection{第三章}

\Emph{他驳斥了关于出生者必须与生育者同体的异议。}\par

你异端的第二个亵渎性的诽谤或诽谤性的亵渎，是你说所生者必须与生育者同体。这与前一条并无太大差别，因为它在措辞上而非事实和实际上有所不同。因为当我们论及天主的诞生时，你认为从玛利亚不能生出更有能力者，正如你前面认为不能生出更年长者一样。因此，你可以认为对此的答复与对你先前所说的话相同：或者你可以认为，对你现在所提出的这一论断的答复也适用于那一条。那么你说，所生者必须与生育者同体。如果这是指地上的受造物，那确实是如此。但如果是指天主的诞生，为什么在祂的诞生这事上你要考虑自然的先例？因为受造物服从于造物主，而非造物主服从于祂的受造物。但你是否愿意更充分地了解，你的这些诽谤不仅是邪恶的，而且是愚蠢的，是一个完全看不到天主全能者的空谈？请告诉我，我求你，你这位认为同类只能由同类产生的人：在古时旷野中喂养以色列子民的那无法解释的鹌鹑群是从何而来的？因为我们无处读到它们先前是由母鸟所生，而是它们被带来并突然出现。再者，那四十年间降在希伯来人营中的天粮又是从何而来？玛纳生玛纳吗？但这些是关于古代奇迹的。那么更近的奇迹呢？主耶稣基督用几个饼和几条小鱼喂饱了跟随祂的无数群众，而且不止一次。他们满足的原因不在于食物：因为一种隐秘无形的原因使饥饿的人满足，尤其是当他们吃饱后，剩下的食物比他们饥饿时摆在面前的还要多得多。这一切是如何发生的：当吃的人满足时，食物本身通过一种奇异的增长而倍增？我们读到在加里肋亚，水变成了酒。请告诉我，那本性为水的东西如何产生了与自身性质完全不同的东西？尤其是当（这正适用于主的诞生）它是从较低劣之物产生更高贵之物时？那么请告诉我，如何从仅仅的水中产生出丰盛而优质的酒？如何倒出来是一种东西，灌进去的是另一种东西？那水泉本身的性质能如此改变从中取出的水成为美酒吗？还是容器的特性或仆人的勤勉造成了这一切？肯定两者都不是。那么，事实的\Emph{方式}为何不能由心的思想所理解，而事实的\Emph{真理}却为良心所坚定持守？在福音中，泥土被放在一个盲人的眼睛上，当洗去时，眼睛就产生了。水有能力生出眼睛吗？泥土有能力创造光明吗？当然不能，特别是水对盲人毫无用处，泥土反而会阻碍看得见之人的视线。那么，一件本身性质有害的东西如何成为恢复健康的手段？通常对有视力的人有害的东西，那时却成了医治的工具？你说这是天主的大能所致，是天主的医治所促成，我们所说的一切都是藉着神圣的全能而成就的；这全能能从非常规的材料中塑造新事物，能从对立之物中制造出有用的东西，能将属于不可能和不可行领域的事物转变为可能和现实的行为。\par

% structure: p0011 source=h2 target=subsection reason=relative-heading-rank; base=h2

\subsection{第四章}

\Emph{天主如何在祂在时间中的诞生以及在其他一切事上显示了祂的全能。}\par

那么，关于主实际的\OriginalTerm{诞生}{nativity}，你应当宣认与一切其他事物相同的真理。要相信天主在祂愿意的时候降生，因为你并不否认祂能成就祂所愿意的事；除非你或许认为，那属于祂、为一切其他事物的能力，在祂自身方面有所欠缺，并且祂的\OriginalTerm{全能}{Omnipotence}虽出自祂并渗透万有，却不足以促成祂自己的诞生。关于主的诞生，你向我提出这样的异议：没有人能生出时间上在先者；关于全能天主所受的诞生，你说被生者应当与生者具有\OriginalTerm{同一性体}{one substance}；仿佛你是在处理人的法律，如同对待任何普通人，你可以向他提出不可能性的异议，因你将他纳入世间事物的软弱中。你说，对所有人而言，有共同的出生条件，只有一条生育的法则；并且，天主禁止对所有人生效的事，不可能唯独发生在全人类中的一个人身上。你并不明白你在说谁；也未看清你所谈论的是谁；因为祂是一切条件的创造者，是一切本性的法则本身，凡人力所能及与不能及者，都藉祂而存在：因为祂确实为二者设定了界限，即人的能力应扩展到多远，以及软弱不应超越的极限。那么，在拥有一切能力和可能性的祂身上，你多么狂妄地将人类的不可能性作为异议提出！如果你以世间的软弱来评价主的位格，并以人的规则来衡量天主的全能，那么你必然无法找到任何似乎与天主相配的事物，来涉及祂身体的苦难。因为如果玛利亚能生出时间上先于她的天主，这对你来说似乎不合理，那么天主被人钉在十字架上，又怎会合理呢？然而，这位被钉十字架的天主亲自预言说：\BibleQuote{人岂能困扰天主呢？你们却困扰了我。} \BibleRef{拉 3:8} 那么，如果我们不能认为主由童贞女所生，因为被生者先于生祂者，我们怎能相信天主有血呢？然而，曾对厄弗所的长老们说：\BibleQuote{你们要牧养天主的教会，就是祂用自己的血所取得的。} \BibleRef{宗 20:28} 最后，我们怎能认为\OriginalTerm{生命之主}{Author of life}祂自己被剥夺了生命？然而伯多禄说：\BibleQuote{你们杀了生命的创始者。} \BibleRef{宗 3:15} 没有人在地上而能同时在天上；然而主自己却说：\BibleQuote{人子，祂是在天上的}？ \BibleRef{若 3:13} 那么，如果你认为天主的诞生不是由童贞女，因为被生者必须与生者有同一性体，你又怎能相信不同的事物能从不同的本性产生呢？因此，照你看来，风没有突然带来鹌鹑，吗哪没有落下，水没有变成酒，几千人没有用几只饼吃饱，盲人也没有在用泥抹眼后复明。但如果这一切看似不可思议且违反本性，除非我们相信它们是因天主而行，你为什么在祂的诞生这件事上否认你在祂的作为中所承认的呢？或者，祂既为了人的救助和益处而不拒绝行事，难道反倒不能为自己的诞生和来临贡献什么吗？\par

% structure: p0014 source=h2 target=subsection reason=relative-heading-rank; base=h2

\subsection{第五章}

\Emph{祂以取自本性本身的证明来显示，祂的对手所设定的法则，即被生者应当与生者有同一性体，在许多情形下不成立。}\par

若继续就此话题详述，难免冗长且近乎幼稚。但为驳斥你的愚昧与狂妄——你坚持认为所生者应与生养者同一实体，即无物能产生与自己本性不同之物——我将引用一些尘世实例，以说服你：许多受造物确是从不同本性之物中产生的。这并非因为在此等事上可以或应当作任何比较，而是为使你不再怀疑：在神圣的诞生中，正如你亲眼所见在脆弱的尘世之物中发生的那样，同样可能发生。蜜蜂虽是最微小的受造物，却如此灵巧机敏，据记载它们可以从完全不同本性之物中产生并涌现。因它们是具有奇妙智慧的生物，不仅具有感觉，更有预见，它们从采集的植物花朵中产生。你认为还能举出并引用什么更重大的实例呢？从无生命物中产生出生命物；从无感觉物中产生出有感觉物。有何工匠、建筑师在那里？谁塑造了它们的身体？谁将灵魂吹入其身？谁赋予它们清晰的声音以彼此交谈？谁形成了并安排了它们足部的和谐、口器的精巧、翅膀的整洁？它们的能力、愤怒、预见、运动、平静、和谐、差异、战争、和平、安排、巧计、事务、治理——所有这些它们与人类共有的东西——是从谁的教育或恩赐得来的？是从谁的植入或教导学得的？它们是通过生育获得的吗？还是在母腹中或从母体的血肉中学得的？它们从未在母腹中，也未曾经历生育。仅仅是它们所采集的花朵被带进蜂房，并由此通过奇妙的巧计，蜜蜂便涌现出来。因此，母体的子宫并未将任何东西传给后代：蜜蜂也不是由蜜蜂产生的。它们只是蜜蜂的制造者，而非作者。从植物的花朵中产生了活物。植物与动物之间有何亲缘关系？我想你已看出谁是这些事物的创造者。去吧，现在去探究：主能否在祂自己的诞生中实现你在这些最微小的受造物身上所见的事。也许此后无需再添加什么。但为支持论证，我们仍可添加一些对证明论点可能并非必要的事。我们看到天空突然昏暗，遍地充满蝗虫。请指出它们的种子——它们的出生——它们的母亲。因为如你所见，它们从那里而来，从那里获得出生。请在所有这一切中断言：所生者必须与生养者同一实体。而在这些断言中，你将被证明如同你在否认主的诞生时那般狂妄，那般愚蠢。还有呢？连\Emph{你}都认为我们必须继续下去吗？但我们仍将再添一事。毫无疑问，鸡身蛇尾怪是从埃及人称为朱鹭的鸟的卵中产生的。鸟与蛇之间有何同族或亲缘关系？为何所生之物不与生养之物同一实体？然而，生养者并非这一切的作者，所生者也不明白这一切：它们是由隐秘的原因和某种不可解释的、复杂的自然法则所产生，这法则产生了它们。而你却拿尘世观念中的琐碎论断来反对祂的诞生，同时你却不能解释那些由祂的命令和统治所产生的事物之起源；祂的意志成就一切，祂的权能导致一切；无物能反对或抵抗祂；而祂的意志足以成就一切可能之事。\par

% structure: p0017 source=h2 target=subsection reason=relative-heading-rank; base=h2

\subsection{第六章}

\Emph{他驳斥了\Person{聂斯多略}的另一个论据，他在该论据中试图证明基督在每一方面都与\Person{亚当}相似。}\par

但既然我们无法（尽管我们更愿如此）忽略它们，现在是揭露你们其余更精微、更具欺骗性的亵渎言论的时候了，至少使它们不至欺骗无知的人们。在你们一篇有害的论文中，你们主张并说：\Quote{既然人是天主本性的肖像，魔鬼把它拉下来并打碎了，天主为自己的肖像而悲伤，如同一位皇帝为自己的雕像而悲伤，便修复打碎的肖像；且从童贞玛利亚无生地形成一个本性，与无生而生的亚当相似；并藉着人高举人的本性：因为死亡藉着人而来，死人的复活也藉着人而来。}他们告诉我们，某些下毒者有一种习惯，将蜜糖与毒药混合在他们备好的杯子里；这样有害的成份就可以被甜味掩盖：当一个人被蜜糖的甜美所迷住时，就可能被致命的毒药所害。所以，当你们说人是天主本性的肖像，魔鬼把它拉下来并打碎了，天主为自己的肖像悲伤如同一位皇帝为自己的雕像悲伤时，你们（可以说）将杯口涂抹了像蜜糖一样的甜物，好使人们饮尽递给他们的杯，却察觉不到其致命性，只尝到诱人的味道。你们提出天主的名字，为的是在宗教的名义下说出谎言。你们把神圣之物摆在前面，为的是说服人们相信虚假的东西：借着你们对天主的宣认，你们竟然设法否认你们所宣认的那一位。谁看不见你们的目标呢？你们在策划什么？你们确实说天主为自己的肖像悲伤如同皇帝为自己的雕像，修复了打碎的肖像，从童贞玛利亚无生地形成一个本性，与无生而生的亚当相似，并藉着人举起人的本性，因为死亡藉着人而来，死人的复活也藉着人而来。那么，凭着你们所有的热忱，所有的空谈，你们这狡猾的策划者，用你们圆滑的主张，在前面提到天主，却在结论中降低到（仅仅是一个）人：最终你们把他贬低到仅仅是一个人的境地，在谦卑的幌子下，你们已经从他身上夺走了天主的荣耀。你们还说，天主的良善已恢复了魔鬼打碎并毁灭的天主肖像，因为你们说他恢复了打碎的肖像。现在，你们多么狡猾地说他恢复了打碎的肖像，为的是让我们相信，在他内，肖像得到了恢复，却没有什么比这实际肖像更多的东西，而这个肖像的恢复正是要实现的。如此你们提出主只与亚当相同：肖像的修复者并不比这实际可毁坏的肖像更好。最后，在接下来的话里，你们显明了你们的目标和意图，当你们说他从童贞玛利亚无生地形成了一个本性，与无生而生的亚当相似，并藉着人举起人的本性时。你们主张主耶稣基督在各方面都与亚当相似：一个无生，另一个无生；一个仅仅是凡人，另一个仅仅是凡人。因此，你们看，你们小心翼翼地防范并阻止我们认为主耶稣基督在任何方面比亚当更大或更善：因为你们用同一标准来比较两者，以至于你们会认为，如果你们在基督身上添加更多，就减损了亚当的完美。\par

% structure: p0020 source=h2 target=subsection reason=relative-heading-rank; base=h2

\subsection{第七章}

\Emph{异端者通常以圣经的外衣掩盖自己的教义。}\par

你们说：\Quote{因为正如，}你们说，\Quote{死亡藉着人而来，死人的复活也藉着人而来。}你们真的想要藉着宗徒的见证来证明你们错误而邪恶的见解吗？你们将\Quote{特选的器皿}与你们的邪恶思想混在一起，使他蒙受耻辱吗？我的意思是，既然你们无法理解你们救恩的作者，因此宗徒必须被说成是否认天主。然而，如果你们想使用宗徒的见证，为什么只满足于一个，而将其余的都沉默地带过？为什么你们不立刻加上：\Quote{保禄，宗徒，不是由人，也不是藉着人，而是藉着耶稣基督：}\BibleRef{迦 1:1}或这句：\Quote{我们在成全的人中讲智慧：}紧接着：\Quote{这智慧，}他说，\Quote{今世的首领中没有一个人认识的；因为他们若认识了，就不至于将荣耀的主钉在十字架上。}\BibleRef{格前 2:6}或这句：\Quote{因为在天主内，圆满地居住着整个天主性，是有形有体的。}\BibleRef{哥 2:9}以及：\Quote{只有一个主耶稣基督，万物藉着他而存在。}\BibleRef{格前 8:6}难道你们部分同意、部分不同意宗徒，只在因降生成人而称基督为人的范围内接受他，却在他说基督是天主时拒绝他吗？因为保禄并不否认耶稣是人，但他仍宣认那人是天主：并宣告复活的来临藉着人，以这样的方式表明在那个人内，天主复活了。因为看，他是否宣告那复活的是天主，正如他作证说那被钉的是荣耀的主。\par

% structure: p0023 source=h2 target=subsection reason=relative-heading-rank; base=h2

\subsection{第八章}

\Emph{异端者只将天主性的阴影归于基督，并因此主张他应与天主一起受敬拜，但并非作为天主。}\par

但为避免你视主耶稣为全体民众中的一员，你归给祂某种荣耀，藉着将尊荣归祂为圣人，而非作为真人和真天主的神性。因为你说什么？\Quote{天主成就了主的降生成人。让我们与天主一起敬礼\OriginalTerm{承受天主者}{Theodochos}的形态，作为天主性的一种形态，作为不能与神圣联系分离的形象，作为不可见天主的肖像。}上面你称亚当为天主的肖像，这里你称基督为肖像：你说一个是雕像，另一个也是雕像。但我以为，为了天主的荣耀，我们应当感激你，因为你允许承受天主者的形态与天主一同受敬拜：然而你这样做并非尊荣祂，反而亏负祂。因为在此你没有将神性的荣耀归于主耶稣基督，反而否定了它。藉着一种狡猾而邪恶的技巧，你说祂应当与天主同受敬拜，为的是你无需承认祂就是天主；而恰恰在那看似欺诈地将祂与天主联合的陈述中，你实际上将祂与天主分离了。因为当你亵渎地说祂决不应作为天主受崇拜，而应与天主同受敬拜时，你便赐予祂与神性接近的结合，为的是摆脱祂神性的真理。哦，你这最邪恶、最狡诈的天主仇敌，你想要在宣认祂的伪装下犯下否认天主的罪行。你说：让我们敬拜祂，作为不能与神圣旨意分离的形象，作为不可见天主的肖像。那么，我想，主耶稣基督是因其仁慈的行为，在我们当中获得了作为创造者和救赎者的尊荣。如果我们被祂从永恒的毁灭中救赎出来，却称我们的救赎者为肖像，我们确实是在试图以相称的服务和相称的忠诚来回应祂的仁慈与良善，如果我们试图除掉祂为了我们而不惜降卑的荣耀的话。\par

% structure: p0026 source=h2 target=subsection reason=relative-heading-rank; base=h2

\subsection{第九章。}

\Emph{如何那些说基督的诞生是隐秘的人是错误的，因为这甚至向先祖雅各伯清楚显示了。}\par

但我猜想，你藉着\Quote{作为隐秘天主的肖像}这话，用次等的尊荣来原谅对主的贬抑。因你称祂为肖像，便将祂与人的地位相比。在称祂为隐秘天主的肖像时，你减损了祂应得的荣耀。因为\BibleQuote{天主，达味说，必会降临；我们的天主，决不缄默。}祂果真来了，并未曾缄默，祂在亲自开口说话之前，早已藉着地上的和天上的见证宣告了祂的来临：星辰指认祂，贤士朝拜祂，天使传报祂。你还想要什么？祂的声音尚在地上沉默，祂的荣耀却已在天上高呼。那么你还说天主在祂内是隐秘的吗？但这并非先知、先祖乃至全部法律的宣告。因为他们没有说祂将是隐秘的，尽管他们全都预言了祂的来临。你在可悲的盲目中犯错，寻找亵渎的根据却寻不见。你说祂在来临后仍是隐秘的。我坚持说祂甚至在来临前也不是隐秘的。因为那要由童贞女诞生天主的奥秘，难道逃过了那位著名先祖的认知吗？他因与他同在的天主的显现而获得了一个称号，从而从'篡夺者'之名上升为'以色列'之名？当他与那与他摔跤的人角力后，明白了将来降生成人的奥秘，便说：\BibleQuote{我面对面见了天主，我的生命得以保存。}请问，他看见了什么，以致他相信他看见了天主？难道天主在雷电中显现给他吗？还是当天裂开时，神性耀眼的面容向他显示？绝不是：相反，他看见了一个人，却承认了天主。哦，他真配得上他所接受的名字，因他以灵魂之眼而非肉体之眼赢得了天主所赐称号的尊荣！他看见一个人形与他摔跤，便宣告他看见了天主。他确实知道那人形就是天主：因为当时天主显现的形状，正是祂后来要以真实形状来临的形状。虽然，我们何必惊讶如此伟大的先祖毫不犹豫地相信了天主亲自向他如此显明的事，当他说：\BibleQuote{我面对面见了天主，我的生命得以保存。}？天主如何向他显现了如此多的神性临在，以致他能说天主的容颜向他展示了？因为似乎只有一个人向他显现，他实际上在角力中胜了那人。但天主一定藉着预兆安排这事，以便没有人不信天主由人诞生，因为早已在这之前先祖已见过人形的天主。\par

% structure: p0029 source=h2 target=subsection reason=relative-heading-rank; base=h2

\subsection{第十章。}

\Emph{他汇集更多同一事实的见证。}\par

但既然许多见证并不缺乏，我为何在一例上停留如此之久呢？因为甚至那时，当天主将要在肉身中来临的事实怎能逃过人的认知？当先知向全人类公开论及祂说：\BibleQuote{看，你们的天主；}并在别处说：\BibleQuote{看，我们的天主。}以及：\BibleQuote{全能的天主，永远的父，和平的君王；}又说：\BibleQuote{祂的王国没有终结。}此外，当祂已经来到时，祂已来临的事实怎能逃过那些公开宣认祂已来临者的认知？难道伯多禄不知道天主的来临吗？当他说：\BibleQuote{你是基督，永生天主之子。}\BibleRef{玛 16:16}难道玛尔大不知道她在说什么，或她相信的是谁吗？当她说：\BibleQuote{是的，主，我信你是基督，永生天主之子，来到这世界的。}\BibleRef{若 11:27}而所有那些寻求祂治愈疾病、恢复肢体或复活死者的人，他们是从人的软弱还是从天主的全能求这些事呢？\par

% structure: p0032 source=h2 target=subsection reason=relative-heading-rank; base=h2

\subsection{第十一章。}

\Emph{魔鬼如何因诸多理由被迫承认基督是天主。}\par

最后，至于魔鬼本身，当他用各种诱惑的外表和邪恶的伎俩试探祂时，他在无知中怀疑什么，或者想通过试探祂发现什么？是什么如此强烈地促使他，使他在人谦卑的形态下寻找天主？他是否通过先前的证据学到了这一点？或者他曾知道有谁以天主的身份来到人的身体中？绝对没有。但是，正是由于神迹的有力证据，行动的强大结果，以及真理本身的话语，他被驱使去怀疑并查究这件事：因为他已经一次从若翰那里听到：\BibleQuote{请看天主的羔羊，请看除免世罪者。}\BibleRef{若 1:29} 又从同一个人那里听到：\BibleQuote{我需要受祢的洗，而祢却来就我吗？}\BibleRef{玛 3:14} 那从天上降下并停在主头上的鸽子也成了自显的天主清楚明白的证据。那从天主发出的声音，并非谜语或预表，也感动了他，当它说：\BibleQuote{祢是我的爱子，我因祢而喜悦。}\BibleRef{玛 3:17} 尽管他在耶稣身上看到一个外在的人，但他仍在寻找天主子，当他说：\BibleQuote{祢若是天主子，就命这些石头变成饼。}\BibleRef{路 4:3} 对那人的注视是否驱散了魔鬼对祂天主性的怀疑，以至于因为看到一个人，他就相信祂不可能是天主？绝对没有。但他说什么？\BibleQuote{祢若是天主子，就命这些石头变成饼。} 他当然不怀疑他所查究之事存在的可能性。他的焦虑在于其真实性。对于不可能性，他并不确信。\par

% structure: p0035 source=h2 target=subsection reason=relative-heading-rank; base=h2

\subsection{第十二章。}

\Emph{他将魔鬼的这种观念和合理的怀疑与他对手的顽固不变的想法相比较，并表明后者比前者更糟糕，更亵渎。}\par

但他一定知道主耶稣基督由玛利亚所生：他知道祂被裹在襁褓里，放在马槽里；祂的童年是人一生开始时的贫穷状态；祂的婴孩期没有适当的摇篮辅助品；此外，他不怀疑祂有真正的血肉，且是一个真正的人所生。为什么这在他看来不足以使他确信？为什么他相信祂不可能是天主，而他知道祂是真人呢？那么，你这可怜的精神错乱者，学习吧，你这疯子，你这残忍的罪人，学习吧，我恳求你，甚至从魔鬼那里，也要减少你的亵渎。他说：\Quote{祢若是天主子。} 你说：\Quote{祢不是天主子。} 你否认他所问的。除了你，从未有人能在亵渎上超过魔鬼。他所承认在主身上是可能的，你却不相信是可能的。\par

% structure: p0038 source=h2 target=subsection reason=relative-heading-rank; base=h2

\subsection{第十三章。}

\Emph{魔鬼如何始终保持对基督天主性的这种观念（由于他所经历的秘密工作），甚至直到祂的十字架与死亡。}\par

但也许他后来停止了，休息了，当他的试探被战胜时，由于没有结果而搁置了怀疑？不，它始终留在他身上，甚至直到主的十字架，怀疑仍然存在，并因特殊的恐惧而加剧。还需要什么进一步的吗？甚至在那时，他得知如此大的许可被赐予祂的迫害者来反对祂之后，他也没有停止认为祂是天主子。但狡猾的敌人甚至在祂身体的苦难中也看到了天主性的记号，尽管他更希望祂只是一个（平凡）人，却被强迫怀疑祂是天主：因为尽管他宁愿相信他所想要的，却被最确凿的证据驱使他走向他所害怕的。这并不奇怪：因为虽然他看见祂被吐唾沫、鞭打、羞辱、被领到十字架，他却看到神圣的力量在侮辱和冤屈中间仍很丰盛；当圣殿的帐幔裂开，太阳隐藏自己，白昼变暗，一切受造物都感受到苦难的效果：甚至一切不认识天主的受造物，都承认天主性的工作。因此，魔鬼看到这些并战栗，竭力以各种方式认识祂的天主性，甚至在人性的死亡中，借着那些钉祂的人说：\BibleQuote{祂若是天主子，现在就从十字架上下来，我们就相信祂。}\BibleRef{玛 27:42} 他确实感知到，通过祂身体的苦难，我们的主天主正在进行人类救恩的救赎，并且通过它，他正被摧毁和制服，而我们正被救赎和拯救。因此，人类的敌人想要用一切手段和诡计来击败他所知道的一切人正在做的救赎。他说：\Quote{如果，} 他说，\Quote{祂是天主子，现在就从十字架上下来，我们就相信祂：} 目的是让主因话语的责备而感动，并摧毁这奥迹，同时祂报复这冤屈。你看，即使当主挂在十字架上时，也被称为天主子。你看，他们怀疑他们所指的事实。所以，正如我上面所说的，你甚至从祂的迫害者，甚至从魔鬼那里，学习相信天主子。有谁曾达到魔鬼的不信？有谁超越它？\Emph{他} 怀疑祂是天主子，甚至当祂忍受死亡时。\Emph{你} 却否认祂，即使当祂复活了。\Emph{他} 怀疑祂是天主，从他面前隐藏了自己。\Emph{你}，祂已经证明给你，却否认它。\par

% structure: p0041 source=h2 target=subsection reason=relative-heading-rank; base=h2

\subsection{第十四章。}

\Emph{他通过回答从宗徒的话中所引的论证，即\Quote{无父，无母}等\BibleRef{希 7}，来表明异端者如何曲解圣经。}\par

这样，你竟利用圣经来反对天主，并试图把祂自己的见证人引用来反对祂。但这是怎么做的呢？确实，你不仅成了天主的诬告者，也成了这些证据本身的诬告者。事实上，既然你做不到你想做的，你只做你能做的，这并不奇怪：既然你不能使神圣的见证人反对天主，你就做你能做的，歪曲它们。因为你声称：那么保禄说谎了，因为他论及基督说：\Quote{无母，无族谱。}\BibleRef{希 7:3}我问你，你认为保禄说这话是指谁？是指天主圣子和圣言，还是指你所分离出来并亵渎地断言仅是个人的基督？如果是指基督——你坚持他仅是个人——一个人怎能无母而生，且没有母系的族谱呢？但如果是指天主的圣言和圣子——当同一位宗徒，即你自己不敬地想象为你的见证人，在同一处并藉着同一个见证，证明你所断言的无母的那位也是无父的，说：\Quote{无父、无母、无族谱}时，我们又能作何理解？因此，如果你利用宗徒的见证，既然你断言天主子\Quote{无母}，你也必须犯有亵渎地认为祂\Quote{无父}之罪。那么，你看你在这不敬的堕落中陷入了何等境地：你急于坚持你的乖戾和邪恶，以致在说天主子没有母亲的同时，你也必须否认祂有父亲——自世界创始以来，除了疯子，从未有人这样做过。而且，我不知道这是出于更大的邪恶还是愚蠢；因为还有什么比给予祂子的名号却试图保留父的名号更愚蠢、更愚蠢的呢？但你说，我没有保留，我没有否认。那么，是什么疯狂驱使你引用那段话，在那里，当你声称祂没有母亲时，你也似乎必须否认祂有父亲？因为在同一段中，祂被说成是无母和无父，因此，如果可以在那里理解为祂是无母的，那么我们同样也必须相信祂是无父的，正如我们理解祂是无母的那样。但是那种急于否认天主的轻率狂热并没有看到这一点；当它残缺不全地引用完整记载时，它未能看出这无耻而明显的谎言可以通过揭示圣卷的内容而被驳斥。啊，愚蠢的亵渎和疯狂！它既看不到它应当遵循什么，甚至也没有智慧看到它本可读到的内容：仿佛因为它可以摆脱自己的理智，它就能使其他人失去阅读的能力，或者因为失去了心智的眼睛，所有人都会在阅读时失去他们头上的眼睛。那么，异端者，请听你篡改的段落：完整地听，彻底地听，你残缺不全地引用和支离破碎的段落。宗徒想要向每个人阐明天主的双重诞生——为了显示主如何在神性中降生和如何在肉身中降生，他说：\Quote{无父、无母}：因为一个属于神性的诞生，另一个属于肉身的诞生。因为正如祂在神性中受生是\Quote{无母}的，同样祂在肉身中也是\Quote{无父}的：因此，尽管祂并非无父或无母，我们必须相信祂\Quote{无父无母}。因为如果你从祂由圣父所生来看，祂是无母的；如果从祂由母亲所生来看，祂是无父的。因此，在这两种诞生中，祂各有一个亲源；而两者合起来，祂又各缺一个：因为神性的诞生不需要母亲，而祂肉身的诞生，祂自己就足够了，无需父亲。因此宗徒说：\Quote{无母、无族谱。}\par

% structure: p0044 source=h2 target=subsection reason=relative-heading-rank; base=h2

\subsection{第十五章}

\Emph{基督如何能被宗徒说成是无族谱的。}\par

他怎会说主是\Quote{无族谱的}呢？当福音传道者玛窦的福音以救世主的族谱开始，说：\BibleQuote{亚巴郎之子，达味之子耶稣基督的族谱}？\BibleRef{玛 1:1}所以，按照福音传道者，祂有族谱；按照宗徒，祂没有：因为按照福音传道者，祂有母系的族谱；按照宗徒，祂没有，因为祂由圣父所出。因此宗徒说得好：\Quote{无父、无母、无族谱}：并且在他陈述祂无母而受生的地方，他也记载了祂无族谱。因此，关于主的这两种诞生，福音传道者和宗徒的著作彼此一致。因为按照福音传道者，祂在肉身出生时有\Quote{无父}的族谱；而按照宗徒，主在神性中受生时没有\Quote{无母}的族谱，正如依撒意亚说：\BibleQuote{但谁能述说祂的世代呢？}\BibleRef{依 53:8}\par

% structure: p0047 source=h2 target=subsection reason=relative-heading-rank; base=h2

\subsection{第十六章}

\Emph{他表明，如同魔鬼在试探基督时一样，异端者篡改并歪曲圣经。}\par

那么，异端者，你为何没有如此引用你所读过的整段经文呢？所以你看，宗徒规定主是\Quote{没有母亲}的，正如他规定祂是\Quote{没有父亲}而生的一样：好叫我们知道祂是\Quote{没有母亲}的，正如我们理解祂是\Quote{没有父亲}的一样。而且既然不可能相信祂完全\Quote{没有父亲}，我们也就不能理解祂完全\Quote{没有母亲}。那么，异端者，你为何没有如此完整无缺地引用你在宗徒书中所读到的呢？但是你插入一部分，省略一部分；并篡改真理之言，以便你能通过你的邪恶行为构建你的错误观念。我知道谁是你的师傅。我们必须相信你曾领受了\Emph{他的}教训，因为你正效法他的榜样。因为福音中魔鬼诱惑主时如此说：\Quote{祢若是天主子，就跳下去。因为经上记载：祂必为你吩咐祂的天使，在你的一切道路上保护你。} \BibleRef{路 4:9} 他说完这话，就删去了上下文和与之相关的部分，即：\Quote{你要践踏毒蛇和蜥蜴：你要踏碎狮子和巨龙。} 诚然，他狡猾地引用了前一句而省略了后一句：因为他引用前一句是为了欺骗祂；他闭口不谈后一句是为了避免谴责自己。因为他知道，在先知的话中，他自己就是由毒蛇和蜥蜴、狮子和巨龙所指代的。所以，你也提出一部分，省略一部分；引用一部分来欺骗；省略另一部分，害怕如果你引用全部，你也许会谴责你自己的欺骗。但现在该是转向进一步问题的时候了，因为我们在个别点上停留太久——如我们因渴望给予完整答复而被引导所做的那样——甚至超出了篇幅较长的书的限度。\par

% structure: p0050 source=h2 target=subsection reason=relative-heading-rank; base=h2

\subsection{第十七章}

\Emph{基督的光荣和尊荣不应如此归于圣神，以致否认它由基督自己而来，仿佛在祂内的一切崇高都属他人，并来自另一源头。}\par

你于是在另一篇讨论中说道——不，是在你的另一篇亵渎中——你将圣神与创造了祂人性的天主性分离。因为圣经说，那由玛利亚所生的，是出于圣神。\BibleRef{玛 1:20} 圣神又以\OriginalTerm{义德}{justitia}充满了那受造者；因为经上说：\BibleQuote{祂在肉身内出现，在圣神内成义}。\BibleRef{弟前 3:16} 同样，圣神也使祂为魔鬼所畏惧：祂说：\BibleQuote{因为我是靠着天主的神驱魔}。\BibleRef{路 11:20} 圣神也使祂的肉身成为圣殿：\BibleQuote{我看见圣神仿佛鸽子从天降下，停在祂身上}。\BibleRef{若 1:32} 同样，圣神赐予祂升入天堂。因为经上记载：\BibleQuote{祂藉圣神嘱咐了祂所拣选的宗徒以后，就被接升天}。\BibleRef{宗 1:2} 最后，是圣神赐予基督如此的荣耀。你的全部亵渎因而在于：基督自身毫无所有；祂，正如你所说，只是一个凡人，并没有从圣言——即天主子——那里领受任何东西；相反，祂内的一切都是圣神的恩赐。如果我们能证明你归之于圣神的一切，都是属于祂（基督）的，那么剩下的岂不是证明，你因此想要当作一个凡人看待的祂——因为你说祂的一切都是别人的——因此是天主，因为祂的一切都是祂自己的？我们不但要藉讨论和论证来证明这一点，更要藉天主性本身的声音：因为没有什么比神圣之事更能为天主作证。而且，既然没有什么比天主的荣耀更能认识自身，我们相信，关于天主，没有比那些天主亲自为自己作证的著作更权威的了。首先，关于你所说的圣神创造了祂的人性这一点；如果我们能够承认你并非出于不信而提出这一点，我们本可简单接受。因为我们并不否认主的肉身是因圣神而受孕的；但我们主张，身体是因圣神的合作而受孕，其方式使我们能说，祂的人性是由天主子为祂自己创造的，正如圣神自己在圣经中所说，作证说：\BibleQuote{智慧为自己建造了房屋}。\BibleRef{箴 9:1} 你于是看出，那由圣神受孕的，是由天主子建造并完成的：并非天主子的工作是一回事，圣神的工作是另一回事；而是因着天主性的合一与荣耀，圣神的作为就是天主子的建造；而天主子的建造就是圣神的合作。因此我们读到，不仅圣神降临于童贞女，而且至高者的能力也荫庇了童贞女；既然智慧本身是天主性的圆满，无人能怀疑当智慧为自己建造房屋时，整个天主性的圆满都在场。但你那亵渎的可悲顽固，试图将基督与天主子分离，却没有看出这完全是在将天主性的本性与其自身分离。除非你或许相信，房屋因此是由圣神为祂建造的，因为祂本身不足以为自己建造房屋。但相信祂——我们相信祂凭自己的意志创造了整个天地万物的宇宙——无法为自己建造一个身体，这既荒谬又狂妄；尤其因为圣神的能力就是祂的能力，并且天主圣三的天主性与荣耀如此合一不可分割，以致我们绝对不能认为在圣三的一位中有任何能与天主性的圆满分离的东西。因此，一旦确立了并掌握了这一点：即，按照圣经的信德，当圣神降临于（童贞女）且至高者的能力荫庇她时，智慧为自己建造了房屋；那么你亵渎的其他诽谤便化为乌有。因为毫无疑问，祂凭自己并在自己内创造了万有，因祂的名和信德，甚至信徒的信德也能成就一切。因为祂并不需要他人的帮助，正如那信靠祂能力的人也不需要。因此，至于你的断言：祂被圣神成义，圣神使祂为魔鬼所畏惧，祂的肉身成为圣神的圣殿，以及祂被圣神接升天堂，这一切都是亵渎而狂妄的：并非因为我们要相信，在祂自己所行的这一切中，圣神的合一与合作有所欠缺——因为天主性从不欠缺自身，圣三的能力常存在于救主的工程中——而是因为你坚持认为圣神扶助了主耶稣基督，仿佛祂是软弱无力的；并且圣神将那些祂无法为自己获得的东西赐给了祂。那么，要从神圣的见证学习信靠天主，不要将虚假与真理混杂：因为主题不容许如此，常识也厌恶将魔鬼的思想与神圣的见证混杂的做法。\par

% structure: p0053 source=h2 target=subsection reason=relative-heading-rank; base=h2

\subsection{第十八章}

\Emph{我们当如何理解宗徒的话：\Quote{祂在肉身内出现，在圣神内成义}等。}\par

因为首先，关于你的这一主张，即圣神以\OriginalTerm{义德}{justitia}充满了受造之物，以及你试图用宗徒的证据来证明这一点，宗徒说：\Quote{祂在肉身内显现，在圣神内被证明为义}，你每句陈述都出于不当的理解和狂妄的精神。因为你作此断言，即你坚持说祂是由圣神充满了义德，以表明祂是怎样缺乏义德，正如你声称充满义德是赐给祂的。至于你在这事上使用宗徒的证据，你歪曲了这段圣经经文的顺序和意义。因为宗徒的话并非如你所引用的那样残缺不全。宗徒说的是什么？\Quote{并且，毫无疑问，虔敬的奥迹是伟大的，祂在肉身内显现，在圣神内被证明为义。}\BibleRef{弟前 3:16}你看，宗徒宣告说虔敬的奥迹或圣事被证明为义。因为他并没有如此忘记自己的话和教导，以至于说祂缺乏义德，而他曾一直宣认祂为义德，说：\Quote{祂成了我们的义德、圣化和救赎。}\BibleRef{格前 1:30}在其他地方他也说：\Quote{但你们已经被洗净了，已经被圣化了，已经因我们的主耶稣基督的名成义了。}\BibleRef{格前 6:11}那么，祂需要被义德充满，这是何其远离祂呢？因为祂自己以义德充满万有，而祂的荣耀又怎能没有义德呢？祂的名字本身就使万有成义。你看，你的亵渎是多么愚蠢和狂妄，因为你试图从我们的主那里夺去祂不断倾注于所有信者身上的事物，而祂的源源供给却从未减少。\par

% structure: p0056 source=h2 target=subsection reason=relative-heading-rank; base=h2

\subsection{第十九章}

\Emph{那使祂被畏惧的，不仅是圣神，更是基督自己。}\par

你还说，圣神使祂被魔鬼所畏惧。为拒绝和驳斥这一点，尽管这话语的可怕性质已足够，我们仍要添加一些实例。我请问你，你说魔鬼畏惧祂不是祂自己的作为，而是别人的作为，并坚持这不是祂自己的能力而是一种恩赐，那么，请问，为什么连祂的名字也有那种能力——而按照你的说法，祂自己却没有那种能力呢？为什么因祂的名字魔鬼被驱逐，病人得痊愈，死人复活？因为宗徒伯多禄对坐在圣殿丽门前的那瘸子说：\Quote{因耶稣基督之名，起来行走吧！}\BibleRef{宗 3:6}又在\Place{约培}城，对那瘫痪在床上八年的病人说：\Quote{\Person{艾尼阿}，愿主耶稣基督治好你！起来，整理你的床褥吧！}\BibleRef{宗 9:34}保禄也对那占卜的神说：\Quote{我因耶稣基督之名命令你，从她身上出来！}\BibleRef{宗 16:18}魔鬼就出来了。但由此你该明白，这软弱与我们的主是何等不相称：因为那些祂藉自己的名字使之变强的人，我甚至不称他们为软弱，因为我们从未听说自主复活以来，有任何魔鬼或疾病能抵抗任何一位宗徒。那么，圣神如何使祂被畏惧呢？使别人被畏惧的祂，自己却需要被畏惧吗？或者说，祂在自己内是软弱的，而祂的信德即使通过他人之手也统御万有吗？最后，那些从天主那里获得能力的人，从不用这能力好像是自己的：而是将这能力归于那赐予者；因为能力本身除非通过赐予者的名字，否则没有任何效力。因此，宗徒们和所有天主的仆人们，从不在自己的名字下做任何事，而是在基督的名字和呼求下；因为能力本身从来源得到力量，除非来自本源，不能通过仆人的手赐予。那么，你——你说主与祂的仆人一样（因为宗徒们除了从主领受的以外一无所有，所以你得出主自己除了从圣神领受的以外也一无所有；因此你认为祂所拥有的一切，都不是作为主拥有，而是作为仆人领受的）——那么你告诉我，祂如何将这能力当作自己的来使用，而不是当作领受的东西呢？我们读到关于祂的什么？祂对瘫子说：\Quote{起来，拿起你的床，回家去吧！}\BibleRef{玛 9:6}又对为儿子求情的父亲说：\Quote{你回去吧！你的儿子活了。}\BibleRef{若 4:50}当一位母亲的独生子被抬出埋葬时，祂说：\Quote{青年人，我对你说，起来吧！}\BibleRef{路 7:14}那么，祂像那些从天主领受能力的人一样，祈求能力赐给祂，以便藉呼求神圣的名字行这些事吗？为什么祂自己不藉圣神的名字行事，如同宗徒们藉祂的名字行事呢？最后，福音本身关于祂说了什么？它说：\Quote{祂教导他们，像有权威的人，不像经师和法利塞人。}\BibleRef{玛 7:29}或者你硬要说祂是如此骄傲自负，以致将（按你的看法）从天主领受的能力归功于自己的力量？但事实上，能力除非通过其根源的名字，从不屈服于祂的仆人；如果行者将其任何部分声称是自己的，能力就毫无效力，这我们该如何解释呢？\par

% structure: p0059 source=h2 target=subsection reason=relative-heading-rank; base=h2

\subsection{第二十章}

\Emph{他以更强有力的论据来摧毁那个观念。}\par

但为何我们要如此长久地应对你的狂妄亵渎，用这些明显却仍有不足的论证？让我们听天主亲自向祂的门徒说话：\Quote{治好病人，复活死人，洁净癞病人，驱逐魔鬼。}\BibleRef{玛 10:8}又说：\Quote{因我的名，}祂说，\Quote{你们要驱逐魔鬼。}\BibleRef{谷 16:17}那位使自己的名成为权能者，岂有需要另一个名来行使权能？但还有更甚的：\Quote{看，}祂说，\Quote{我已赐给你们权柄，践踏蛇蝎和一切仇敌的权势。}\BibleRef{路 10:19}祂自称良善（祂实在如此），且心地谦卑。既然如此，祂命令别人因祂自己的名行事，而按照你们的说法，祂自己却要靠另一个名行事，这怎么可能？还是说，祂把似乎是自己的东西给了他人，而祂自己（按你们的说法）除非从另一个那里领受，就不拥有它？但请告诉我，哪位从天主领受权能的圣人曾这样行事？难道\Person{伯多禄}不会被当作疯子，\Person{若望}不会被视为狂人，\Person{保禄}不会被认为是丧心病狂，假如他们对病人说：\Quote{因我们的名，起来}；对瘸子说：\Quote{因我们的名，行走}；对死者说：\Quote{因我们的名，活着}；或对某人说：\Quote{我们赐你们权柄践踏蛇蝎和一切仇敌的权势}？那么你从这些看出你的疯狂：因为如果这些话出自人的自信，它们是疯狂的；同样，如果你看不出它们来自天主的权能，你也是完全疯狂的。因为你必须承认两个可能性之一：要么人能拥有并赐予天主权能，要么，既然没有人能做到，那能做到的那一位就是天主。因为没有人能凭自身的慷慨赐予天主权能，除非那一位本性就拥有它。\par

% structure: p0062 source=h2 target=subsection reason=relative-heading-rank; base=h2

\subsection{第二十一章}

\Emph{必须将祂的肉身和人性成为天主的圣殿，同等归于基督和圣神。}\par

但是，在你的亵渎中接着又说，祂的肉身成了圣神的圣殿，原因是若望曾说过：\Quote{我看见圣神从天降下，停在祂身上。}\BibleRef{若 1:32}因为你试图用圣经的权威来支持你这狂妄的言论：因此让我们看看，这神圣权威是否说了你所说的。\Quote{我看见，}它说，\Quote{圣神像鸽子降下，停在祂身上。}在这里，如果你能，分辨谁更有权能，谁更伟大，谁更应受尊荣？是降下的那一位，还是祂所降下的那一位？是带来尊荣的那一位，还是尊荣被带来的那一位？你在这段经文中哪里找到圣神使祂的肉身成为圣殿？或者，如果天主亲自降下来向人类显示天主，这如何减损天主的尊荣？因为毫无疑问，我们不应认为那被指出高等地位者，比那指出祂高等地位者更低。但切勿相信或制造在天主性中的任何分离：因为同一的天主性和同等的权能完全排除不平等的错误观念。因此在这件事上，有圣父、圣子和圣神的位格，而降下所向的是天主子，降下的是圣神，作证的是圣父，没有人更有尊荣，也没有人受到任何轻视，而是这一切同样归于天主性的圆满，因为天主圣三的每一位都包含整个圣三的光荣。因此无需再多说，只需展示你亵渎的起源和来源。因为荆棘和蒺藜从根部发芽，产生本性的枝条，并通过它们的特征显示其根源。同样，你们这些白拉奇异端的荆棘枝，在萌芽中显示出你们的父亲据说在根部所有的那同样的东西。因为他（正如他的追随者\Person{莱波里乌斯}所说）宣称我们的主通过祂的洗礼成了基督：而你们说，在祂的洗礼时，祂由圣神成了天主的圣殿。这些话不完全相同，但顽固思想完全相同。\par

% structure: p0065 source=h2 target=subsection reason=relative-heading-rank; base=h2

\subsection{第二十二章}

\Emph{基督升天不应只归于圣神。}\par

但是，你还在上述那些不虔之论上加上这条：即圣神将升天赐予了主；你以此亵渎的见解表明，你相信主耶稣基督是如此软弱无力，以致于若不是圣神将祂提升到天上，你想祂今日还会在地上。但为证明此断言，你举出圣经的一段话：因为你说\BibleQuote{祂藉圣神向祂所拣选的宗徒颁命之后，就被接升天。}\BibleRef{宗 1:2} 我该怎样称呼你？我该怎样看待你？你竟以篡改圣经的方式使它们的见证失去见证的效力？一种新的胆大妄为，它以其不虔的论证竭力操纵真理，以致使真理似乎证实虚假。因为宗徒大事录并非如你所臆断。圣经说的是什么？\BibleQuote{耶稣从开始行事与教导，直到祂藉圣神向祂所拣选的宗徒颁命的那一天，被接升天。}这是一个\OriginalTerm{倒置法}{Hyperbaton}的实例，应当这样理解：耶稣从开始行事与教导，直到祂被接升天的那一天，向祂藉圣神所拣选的宗徒颁命。因此，我们也许不必再在这件事上给你多于经文本身的答复，因为整段经文应该足以证明全部真理，既然它的割裂曾有利于你的虚假。但是你仍然认为我们的主耶稣基督除非被圣神提升就不能升天，请告诉我，祂自己是如何说的？\BibleQuote{除了那从天降下的、仍在天上的人子以外，没有人升过天。}\BibleRef{若 3:13} 那么承认你的见解是何等愚昧荒谬吧，你认为祂不能升天，而祂虽降下到地上，却从未离开过天堂。请说：对于祂而言，离开下界升上天堂，岂非可能？因为祂在地上时，仍能常在天上。但祂自己说的是什么？\BibleQuote{我升到我的父那里去。}\BibleRef{若 20:17} 祂难道暗示在这升天中会有别人的帮助介入吗？祂说祂将升天，这事实本身就显示了祂自己权能的效力。\Person{达味}论及主的升天也说：\BibleQuote{天主在欢呼声中上升，主在号角声中上升。}他清楚说明了，那上升者之光荣，在于上升的权能。\par

% structure: p0068 source=h2 target=subsection reason=relative-heading-rank; base=h2

\subsection{第二十三章}

\Emph{他继续同一论证，表明基督不需要别人的光荣，因祂有自己的光荣。}\par

但最后，让我们看看你总结你先前亵渎言论的附加部分。你的话是：\Quote{谁给了基督这样的光荣？}你提到光荣是为了贬低祂。因为当你断言主被赋予了光荣，说祂接受了光荣，你就亵渎地暗示祂需要光荣。因为你那乖谬的见解表明，给予者的慷慨显示了接受者的需要。哦，你可悲的不虔！那么，那神性本身曾预言主耶稣基督升天的话在哪里？说：\BibleQuote{抬起头来，光荣之王要进来。}当祂（依天启的方式）好像以询问者的角色自言自答时：\BibleQuote{谁是光荣之王？}祂立刻加上：\BibleQuote{强有力的主，大能的主，战争中的主。}以一场战斗的形像，显示主在祂胜利中的得胜。然后，为完成其阐释，祂重复了以上引用的谕言，以下结论显示了主进入天堂的威严，说：\BibleQuote{万军之主，祂就是光荣之王。}为使祂取了身体的事实不致干扰其大能神性的光荣，祂教导说，同一位既是万军之主，又是天上光荣之王，而祂先前曾宣布祂是下方战场的胜利者。如今你去说光荣是给予主的吧，尽管预言已说祂是光荣之王，且祂自己也如此为自己作证：\BibleQuote{当人子在自己的光荣中降来时。}\BibleRef{玛 25:31} 你若能，就反驳并否定这一点：即祂见证自己有光荣，而你说祂接受了别人的。尽管我们主张祂有自己的光荣，但我们并不否认祂的光荣属性与圣父和圣神所共有。因为天主所拥有的一切，都属于天主性；光荣的国度属于天主子，如此，它不会不被归于整个天主性。\par

% structure: p0071 source=h2 target=subsection reason=relative-heading-rank; base=h2

\subsection{第二十四章}

\Emph{他以\Person{圣希拉略}的权威支持这一教义。}\par

但现在是结束本书的时候了，是的，也是结束整部著作的时候了。然而，请允许我再引用几位圣德卓越的知名司铎的言论，借当今之信仰来支援我们已经由圣经权威所证明的事理。依拉略，一位以诸般美德和恩宠著称的人，其生活与辩才同样卓著；他身为教会的导师和司铎，不仅凭自身功德，也借他人的进步而前行，并在迫害风暴中坚守不移，以至因他不可征服的信仰的刚毅而获享宣信者的尊位。他在《论信仰》第一书中证明，主耶稣基督，出自真天主的真天主，在世界之前受生，随后又降生为人。又在第二书中说：\Quote{独一天主在圣童贞的胎中长成人的形体；那包含万有、万有都在其权能之下的那一位，按人类出生的律法而被生出。}又在同书中说：\Quote{天使作证说，那出生的那一位是天主与我们同在。}又在第十书中说：\Quote{我们已教导了由童贞所生、天主降生为人的奥迹。}又在同书中说：\Quote{因为当天主降生为人时，祂并非为了不再保持天主而受生。}又在同一位作者\WorkTitle{\BibleBook{玛窦}{福音}注释}的前言中说：\Quote{首先，我们需要知道，为了我们，独生天主应当被认识为降生为人。}又在随后说：\Quote{除是天主之外，祂还应降生为人，这是祂尚未是的。}又在同一处说：\Quote{于是这第三件事是合适的：既然天主在世上降生为人}等等。这里只是无数段落中的几段。但即便如此，从我们引用的这些段落中，你也看到他是多么清楚而明确地断言天主是由玛利亚所生。那么，你这句话又在何处：\Quote{受造物不能生出造物主；凡由血肉生的是血肉。}逐一引用每位作者关于此点的段落未免冗长。我必须尝试列举它们，而非解释它们：因为它们将充分说明自己。\par

% structure: p0074 source=h2 target=subsection reason=relative-heading-rank; base=h2

\subsection{第二十五章}

\Emph{他表明安博罗削与圣依拉略一致。}\par

安博罗削，那位卓著的天主司铎，他从未离开上主之手，始终像宝石一样在上主的手指上闪耀，他在给童贞女的书中这样说：\Quote{我的爱人白而且红。白，因为祂是父的荣耀；红，因为祂由童贞所生。但请记住，在祂内，天主性的标记比身体的奥迹更古老。因为祂并非从童贞开始存在，而是那位已经存在者进入了童贞。}又在圣诞节说：请看上主母亲的奇迹：童贞女怀孕，童贞女生子。她怀孕时是童贞女，有孕时是童贞女，生产后仍是童贞女。如\BibleBook{}{厄则克耳}所说：\Quote{门是关闭的，不打开，因为上主从中经过。}\BibleRef{则 44:2}多么辉煌的童贞，多么奇妙的多产！世界的主宰诞生了：而生祂的那一位竟没有哭喊。子宫空了，真正婴孩出生了，而童贞却未受损。当天主降生时，贞洁的德能理应变得更加伟大，祂的纯洁不应因那来医治败坏者的离去而受损。又在\WorkTitle{\BibleBook{路加}{福音}注释}中说，\Quote{有一位特别被拣选，来生出天主，她已许配给一个丈夫。}他确实宣称天主由童贞所生。他称玛利亚为天主之母。那么，你那可怕可咒的言论——问一个本性不同的人如何能成为母亲——又在哪里呢？但既然他们称她为母亲，那所生的是人性而非天主性。所以，那位卓越的信仰导师既说生祂的那一位是人，又说所生的那一位是天主：然而这并不足以成为不信的理由，而只是信仰的奇迹。\par

% structure: p0077 source=h2 target=subsection reason=relative-heading-rank; base=h2

\subsection{第二十六章}

\Emph{他在上述内容上又增加了圣热罗尼莫的见证。}\par

热罗尼莫，公教导师，其著作如同神圣明灯照耀全世界，他在致欧斯托钦的书中说：\Quote{天主子为我们的救恩成了人子。祂在子宫中等待十个月才出生；那手中托着世界的那一位，竟被容纳在狭小的马槽里。}又在\WorkTitle{\BibleBook{}{依撒意亚}注释}中说：\Quote{因为万军的上主，荣耀的君王，亲自降入童贞的子宫，进入并从永远关闭的东门出来。}\BibleRef{则 44:2}加俾额尔对童贞女说：\BibleQuote{圣神要临于你，至高者的能力要荫庇你，因此那要诞生的圣者，将称为天主子。}又在\BibleBook{}{箴言}中说：\BibleQuote{智慧建造了房屋。}请将此与你那学说——或者不如说你的亵渎——相比，你断言天主是月份的创造者，却不是月份所生的。因为看，热罗尼莫，一位学识极大、教义最纯正、最受认可的人，几乎用与你否认天主子是月份所生相同的言词证明，祂是月份所生的。因为他说，祂在子宫中等待十个月才出生。但或许这人的权威在你看来毫无价值。你可以认为，每个人都说同样的事、用同样的言词，因为凡不否认天主子是童贞所生的人，都承认祂是月份所生的。\par

% structure: p0080 source=h2 target=subsection reason=relative-heading-rank; base=h2

\subsection{第二十七章}

\Emph{他在上述内容上又添加了鲁菲努斯和真福奥斯定的见证。}\par

\Person{鲁菲诺}，一位基督徒哲学家，在教会圣师中不无地位，在其\WorkTitle{信经释义}中，就主的诞生作证如下。他说：\Quote{因为天主子生于童贞女，并非主要与血肉相连，而是生于灵魂，即血肉与天主之间的媒介。}他是否隐晦地见证天主生于人？希波勒吉斯的司铎\Person{奥古斯丁}说：\Quote{为使人能生于天主，天主首生于人：因为基督是天主。基督生于人时，只在地上需要一位母亲，因为祂在天上常有一位父亲，祂生于天主，借祂我们得以受造，也生于女人，借她我们得以再造。}又，在此处：\Quote{圣言成了血肉，居住在我们中间。那么，你何必惊讶人能够生于天主？请看天主自己是如何生于人的。}又在致\Person{沃卢西安}的书信中：\Quote{但梅瑟本人和其余先知最真实地预言了主基督，并赐给祂莫大荣耀：他们宣告祂要来，不是像他们那样，也不是仅仅在行奇迹的同样权能上更大，而明明白白是万有的主天主，并为了人而成为人。因此，祂自己也愿意行他们所做之事，以免祂自己没有行那些借他们而行之事，显得荒谬。但祂仍然有理由做一些特别的事：即生于童贞女，从死者中复活，升天。若有人以为这对天主来说太小，我不知道他还能期待什么。}\par

% structure: p0083 source=h2 target=subsection reason=relative-heading-rank; base=h2

\subsection{第二十八章}

\Emph{因为他将要提出希腊或东方主教的证词，他首先引用了纳齐盎的圣\Person{额我略}。}\par

但也许因为我们所列举的这些人都来自世界各地，他们的权威在你看来可能价值较低。这诚然是荒谬的，因为信仰并不因地方而改变，我们应考虑一个人\Emph{是什么}，而非\Emph{在哪里}；尤其因为宗教将众人联系在一起，在同一个信仰中的人也必在同一个身体内。但我们仍要为你举出一些你无法轻视的人，即来自东方的人。\Person{额我略}，这位知识与教义的最明亮的光，他虽已去世一段时间，但在权威与信仰中仍然活着；他虽在身体上离开教会已久，但在言语和权威上并未离弃它们。他说：\Quote{当那天主从童贞女出来时，在祂所取的人性中，祂存在于两个对立者之一中，即血肉和灵魂；一个被纳入天主，另一个被提升到神性的恩宠中。哦，新的、闻所未闻的混合！哦，奇妙而精妙的结合！那本来存在的，成为有；造物主受造；无限者被灵魂（即天主与血肉之间的媒介）所拥抱；使万有富足的，成为贫穷。}他又论主显说：\Quote{但发生了什么？关于我们、为我们做了什么？发生了一种新的、闻所未闻的本性转变：天主成了人。}又在此处说：\Quote{天主子也开始成为人子，并非从祂所是的改变，因为祂是不变的，而是取了祂所不是的：因为祂慈悲为怀，使那不能受拥抱的，如今得以被拥抱。}你看他多么庄严而高贵地肯定祂天主性的威严，为引入降生成人的屈尊：因为这位杰出的信仰导师深知，在天主入世赐予我们的一切恩惠中，这是首要的，即丝毫不减损祂的荣耀。因为凡天主赐予人的，应当在我们心中增加对祂的爱，而非减少我们对祂的尊敬。\par

% structure: p0086 source=h2 target=subsection reason=relative-heading-rank; base=h2

\subsection{第二十九章}

\Emph{其次，他引用了圣\Person{亚大纳削}的权威。}\par

\Person{亚大纳削}，亚历山大城的司铎，恒心与德行的卓越典范，异端迫害的风暴考验了他却未击垮他；他的生活始终如同明镜，在获得宣信者的尊位之前几乎已得到殉道的奖赏。让我们看看他对主耶稣基督和主的母亲的看法。他说：\Quote{因此，如我们常说的，这就是圣经的心意和印记：即在同一救主中，有两件事必须领会：（1）祂永远是天主，是圣父的圣子、圣言、光与智慧；（2）之后，为我们，祂从童贞女玛利亚——\OriginalTerm{天主之母}{Theotocos}——取了血肉，成了人。}又在其后说：\Quote{于是许多人是圣洁的、脱离罪恶的：\Person{耶肋米亚}也从母胎中被圣化，\Person{若望}还在母胎中时，听见\OriginalTerm{天主之母}{Theotocos}玛利亚的声音就欢喜踊跃。}他明确地说，天主、天主子——他借其言表明众人的信仰——即\Quote{圣言、与圣父的光与智慧}，为我们取了血肉；因此他称童贞玛利亚为\OriginalTerm{天主之母}{Theotocos}，因为她是天主之母。\par

% structure: p0089 source=h2 target=subsection reason=relative-heading-rank; base=h2

\subsection{第三十章}

\Emph{他又引用了圣\Person{金口若望}。}\par

至于君士坦丁堡主教职的荣耀——\Person{若望}，其圣洁生命未经历外邦人迫害即获得殉道者的赏报——请听他关于天主圣子降生成人的思想与教导：他说：\Quote{祂若以未经遮蔽的天主性降临，则天、地、海及任何受造物都不能容纳祂，但童贞玛利亚纯洁的子宫却承载了祂。}此人（若望）的信德与道理，即使你忽视其他人，也应当追随并持守，因为敬虔的子民出于对他的爱戴与深情选你作他们的主教。既然教会从安提约基雅教会（此教会从前亦选立了他）召你来作司铎，便相信将在你身上获得在他身上所失去的一切。我问你：难道这些人几乎是用先知般的灵讲说这一切，是为要挫败你的亵渎吗？因为你宣称我们的主、救主基督不是天主；他们却宣称主基督是真天主。你亵渎地妄称玛利亚是\OriginalTerm{基督之母}{Christotocos}而非\OriginalTerm{天主之母}{Theotocos}；他们并不否认她是基督之母，同时承认她是天主之母。不但实质，连词句都与你的亵渎相对立：好让我们清楚看见，天主早已为你的亵渎预备了一座不可攻克的堡垒，好使那必将到来的异端攻击，撞在早已预备的真理之墙上而粉碎。而你，最邪恶、无耻的著名城市的玷污者，你是公教圣善子民中灾难性的致命瘟疫，你竟敢在天主的教会中站立并施教，用你粗野亵渎的言语毁谤那从未动摇的信德与公教宣认的司铎们，并说君士坦丁堡城的子民是由于先前导师的过失而误入歧途？那么，你是先前主教的纠正者、古代司铎的控告者？你比\Person{额我略}更优秀，比\Person{聂克达里}更受认可，比\Person{若望}更伟大？以及东方各城的所有其他主教——虽然他们不如我所列举的那几位有名，却拥有同样的信德？就此而言，这已经足够：因为当涉及信德问题时，凡与最好者一致的人，都与最好者同样好。\par

% structure: p0092 source=h2 target=subsection reason=relative-heading-rank; base=h2

\subsection{第三十一章。}

\Emph{他哀叹君士坦丁堡因这位异端者所遭受的不幸厄运；同时敦促市民坚守古代公教及祖传的信德。}\par

\Emph{因此}，我——在名望与功德上都微不足道，虽不敢在那些杰出的君士坦丁堡主教中自居为师之位，却敢以弟子的热忱与热忱自许。因为我有幸蒙受真福若望主教的接纳而进入神圣职务，并奉献于天主；即使身体不在，我心仍在那里；虽然实际在场已不再与那可亲可敬的天主之民相混，但我仍藉精神与他们相连。因此，我与他们一同哀悼、同情，就在刚才，我藉着作品中悲痛的哀叹，如同为我自己的肢体一般，发出了我们共同忧伤的呼声（这是我仅能做到的）；因为诚如宗徒所说：\BibleQuote{若身体的一个部分受苦，所有的部分都一同受苦}\BibleRef{格前 12:26}，何况更大的部分受苦时，小的部分岂不更应同情？若在一个身体中，大的部分感受小的部分的痛苦，而小的部分却不感受大的部分的痛苦，那实在是不近人情。因此，我恳求和祈求你们——那些居住于\Place{君士坦丁堡}城内者，你们因着对祖国的爱而成为我的同胞，因着信德的合一而成为我的弟兄——要远离那吞噬天主之民如同吃食的豺狼（正如经上所载）。不要接触、不要品尝任何属于他的东西，因为这一切都通向死亡。要从他中间出来，保持分离，不要触摸不洁之物。要记念你们古代的导师和司铎：闻名于世的\Person{额我略}、以圣德著称的\Person{纳克塔留斯}、信德与纯洁的奇迹\Person{若望}。我说的是若望，那位像\BibleBook{若望}{福音}作者一样，的确是耶稣的门徒和宗徒，并可说是永远依偎在主胸怀与心里的那位。记念他，我说。跟随他。想想他的纯洁、他的信德、他的道理和圣德。永远记念他作为你们的导师和乳母，你们仿佛是在他的怀中成长。他是你们和我的共同导师，我们是他的门徒和学生。阅读他的著作。持守他的教导。拥抱他的信德和功德。因为虽然达到这个目标艰难而崇高，但即使是追随也是美好而崇高的。在最高尚的事情上，不仅达到目标，甚至尝试效仿也值得赞美。因为几乎没有人会完全错失他试图攀登和达到的目标中的所有部分。因此，他应常留在你们心中，几乎是眼前：他应活在你们心中和思想里。他本人会向你们推荐我所写的这些，因为正是他教导了我所写的内容；所以不要认为这是我自己的，而更多的是他的：因为溪流来自源头，凡你们认为属于门徒的一切，都应归功于老师的荣耀。但超出一切之上，我全心、全灵祈求你，啊，我们的主耶稣基督的天主圣父，求你以你浩大的恩宠，以你爱的恩赐充满我们所写的一切。并且，因为正如我们的主天主你的独生子亲自教导我们的，你如此爱了这个世界，以至于派遣你的独生子来拯救世界，求你将这份恩赐赐予你所救赎的子民，使他们在你的独生子的降生成人中，感知到你的恩赐和他的爱；并使众人明白真理：你的独生子，我们的主天主，为了我们而诞生、受苦并复活，愿他们如此爱这真理，以致他荣耀的屈尊增加我们的爱；不要让他的谦卑导致他在所有人心中的尊荣减少，而要永远使爱增加。愿我们都能正确而明智地领悟他神圣怜悯的恩惠，从而明白：我们欠天主的债越多，因为天主为了我们而将自己降得更低。\par

\SourceRecord{Translated by C.S. Gibson.；Nicene and Post-Nicene Fathers, Second Series；Vol. 11.；Edited by Philip Schaff and Henry Wace.；Buffalo, NY: Christian Literature Publishing Co.,；1894.；<http://www.newadvent.org/fathers/35097.htm>.}
